%% LyX 2.0.3 created this file.  For more info, see http://www.lyx.org/.
%% Do not edit unless you really know what you are doing.
\documentclass[oneside,english,a4paper]{amsart}
\usepackage[T1]{fontenc}
\usepackage[latin9]{inputenc}
\usepackage{verbatim}
\usepackage{amsthm}
\usepackage{amstext}
\usepackage{amssymb}
\usepackage{esint}

\makeatletter
%%%%%%%%%%%%%%%%%%%%%%%%%%%%%% Textclass specific LaTeX commands.
\numberwithin{equation}{section}
\numberwithin{figure}{section}
\numberwithin{table}{section}
  \theoremstyle{plain}
  \newtheorem*{thm*}{\protect\theoremname}
\theoremstyle{plain}
\newtheorem{thm}{\protect\theoremname}[section]
  \theoremstyle{definition}
  \newtheorem{defn}[thm]{\protect\definitionname}
  \theoremstyle{plain}
  \newtheorem{prop}[thm]{\protect\propositionname}
  \theoremstyle{remark}
  \newtheorem{rem}[thm]{\protect\remarkname}
  \theoremstyle{definition}
  \newtheorem{example}[thm]{\protect\examplename}
  \theoremstyle{plain}
  \newtheorem{lem}[thm]{\protect\lemmaname}
  \theoremstyle{plain}
  \newtheorem{cor}[thm]{\protect\corollaryname}

%%%%%%%%%%%%%%%%%%%%%%%%%%%%%% User specified LaTeX commands.
%\usepackage{MnSymbol}
%\usepackage[mathlf,textlf,minionint,frenchmath]{MinionPro}
\usepackage{extarrow}
\usepackage{enumerate}
\usepackage{fontenc}
\usepackage[T1]{fontenc}
\usepackage{graphicx}
\usepackage[colorlinks=true, linkcolor=blue]{hyperref}
\usepackage{latexsym}
\usepackage{mathrsfs}
\usepackage{mathtext}
\usepackage{tikz}
\usepackage{textcomp}
\usepackage{upgreek}
\usepackage{xy}


\usetikzlibrary{arrows}
\usetikzlibrary{matrix}
\usetikzlibrary{shapes}
\usetikzlibrary{snakes}
\usetikzlibrary{matrix}














\DeclareMathOperator{\Add}{\textup{Add}}
\DeclareMathOperator{\Aff}{\textup{Aff}}
\DeclareMathOperator{\Alg}{\textup{Alg}}
\DeclareMathOperator{\Ann}{\textup{Ann}}
\DeclareMathOperator{\Arr}{\textup{Arr}}
\DeclareMathOperator{\Art}{\textup{Art}}
\DeclareMathOperator{\Ass}{\textup{Ass}}
\DeclareMathOperator{\Aut}{\textup{Aut}}
\DeclareMathOperator{\Autsh}{\underline{\textup{Aut}}}
\DeclareMathOperator{\Bi}{\textup{B}}
\DeclareMathOperator{\CAdd}{\textup{CAdd}}
\DeclareMathOperator{\CAlg}{\textup{CAlg}}
\DeclareMathOperator{\CMon}{\textup{CMon}}
\DeclareMathOperator{\CPMon}{\textup{CPMon}}
\DeclareMathOperator{\CRings}{\textup{CRings}}
\DeclareMathOperator{\CSMon}{\textup{CSMon}}
\DeclareMathOperator{\CaCl}{\textup{CaCl}}
\DeclareMathOperator{\Cart}{\textup{Cart}}
\DeclareMathOperator{\Cl}{\textup{Cl}}
\DeclareMathOperator{\Coh}{\textup{Coh}}
\DeclareMathOperator{\Coker}{\textup{Coker}}
\DeclareMathOperator{\Cov}{\textup{Cov}}
\DeclareMathOperator{\Der}{\textup{Der}}
\DeclareMathOperator{\Div}{\textup{Div}}
\DeclareMathOperator{\End}{\textup{End}}
\DeclareMathOperator{\Endsh}{\underline{\textup{End}}}
\DeclareMathOperator{\Ext}{\textup{Ext}}
\DeclareMathOperator{\Extsh}{\underline{\textup{Ext}}}
\DeclareMathOperator{\FAdd}{\textup{FAdd}}
\DeclareMathOperator{\FCoh}{\textup{FCoh}}
\DeclareMathOperator{\FGrad}{\textup{FGrad}}
\DeclareMathOperator{\FLoc}{\textup{FLoc}}
\DeclareMathOperator{\FMod}{\textup{FMod}}
\DeclareMathOperator{\FPMon}{\textup{FPMon}}
\DeclareMathOperator{\FRep}{\textup{FRep}}
\DeclareMathOperator{\FSMon}{\textup{FSMon}}
\DeclareMathOperator{\FVect}{\textup{FVect}}
\DeclareMathOperator{\Fibr}{\textup{Fibr}}
\DeclareMathOperator{\Fix}{\textup{Fix}}
\DeclareMathOperator{\Fl}{\textup{Fl}}
\DeclareMathOperator{\Fr}{\textup{Fr}}
\DeclareMathOperator{\Funct}{\textup{Funct}}
\DeclareMathOperator{\GAlg}{\textup{GAlg}}
\DeclareMathOperator{\GExt}{\textup{GExt}}
\DeclareMathOperator{\GHom}{\textup{GHom}}
\DeclareMathOperator{\GL}{\textup{GL}}
\DeclareMathOperator{\GMod}{\textup{GMod}}
\DeclareMathOperator{\GRis}{\textup{GRis}}
\DeclareMathOperator{\GRiv}{\textup{GRiv}}
\DeclareMathOperator{\Gal}{\textup{Gal}}
\DeclareMathOperator{\Gl}{\textup{Gl}}
\DeclareMathOperator{\Grad}{\textup{Grad}}
\DeclareMathOperator{\Hilb}{\textup{Hilb}}
\DeclareMathOperator{\Hl}{\textup{H}}
\DeclareMathOperator{\Hom}{\textup{Hom}}
\DeclareMathOperator{\Homsh}{\underline{\textup{Hom}}}
\DeclareMathOperator{\ISym}{\textup{Sym}^*}
\DeclareMathOperator{\Imm}{\textup{Im}}
\DeclareMathOperator{\Irr}{\textup{Irr}}
\DeclareMathOperator{\Iso}{\textup{Iso}}
\DeclareMathOperator{\Isosh}{\underline{\textup{Iso}}}
\DeclareMathOperator{\Ker}{\textup{Ker}}
\DeclareMathOperator{\LAdd}{\textup{LAdd}}
\DeclareMathOperator{\LAlg}{\textup{LAlg}}
\DeclareMathOperator{\LMon}{\textup{LMon}}
\DeclareMathOperator{\LPMon}{\textup{LPMon}}
\DeclareMathOperator{\LRings}{\textup{LRings}}
\DeclareMathOperator{\LSMon}{\textup{LSMon}}
\DeclareMathOperator{\Left}{\textup{L}}
\DeclareMathOperator{\Lex}{\textup{Lex}}
\DeclareMathOperator{\Loc}{\textup{Loc}}
\DeclareMathOperator{\M}{\textup{M}}
\DeclareMathOperator{\ML}{\textup{ML}}
\DeclareMathOperator{\MLex}{\textup{MLex}}
\DeclareMathOperator{\Map}{\textup{Map}}
\DeclareMathOperator{\Mod}{\textup{Mod}}
\DeclareMathOperator{\Mon}{\textup{Mon}}
\DeclareMathOperator{\Ob}{\textup{Ob}}
\DeclareMathOperator{\Obj}{\textup{Obj}}
\DeclareMathOperator{\PDiv}{\textup{PDiv}}
\DeclareMathOperator{\PGL}{\textup{PGL}}
\DeclareMathOperator{\PML}{\textup{PML}}
\DeclareMathOperator{\PMLex}{\textup{PMLex}}
\DeclareMathOperator{\PMon}{\textup{PMon}}
\DeclareMathOperator{\Pic}{\textup{Pic}}
\DeclareMathOperator{\Picsh}{\underline{\textup{Pic}}}
\DeclareMathOperator{\Pro}{\textup{Pro}}
\DeclareMathOperator{\Proj}{\textup{Proj}}
\DeclareMathOperator{\QAdd}{\textup{QAdd}}
\DeclareMathOperator{\QAlg}{\textup{QAlg}}
\DeclareMathOperator{\QCoh}{\textup{QCoh}}
\DeclareMathOperator{\QMon}{\textup{QMon}}
\DeclareMathOperator{\QPMon}{\textup{QPMon}}
\DeclareMathOperator{\QRings}{\textup{QRings}}
\DeclareMathOperator{\QSMon}{\textup{QSMon}}
\DeclareMathOperator{\R}{\textup{R}}
\DeclareMathOperator{\Rep}{\textup{Rep}}
\DeclareMathOperator{\Rings}{\textup{Rings}}
\DeclareMathOperator{\Riv}{\textup{Riv}}
\DeclareMathOperator{\SFibr}{\textup{SFibr}}
\DeclareMathOperator{\SMLex}{\textup{SMLex}}
\DeclareMathOperator{\SMex}{\textup{SMex}}
\DeclareMathOperator{\SMon}{\textup{SMon}}
\DeclareMathOperator{\SchI}{\textup{SchI}}
\DeclareMathOperator{\Sh}{\textup{Sh}}
\DeclareMathOperator{\Soc}{\textup{Soc}}
\DeclareMathOperator{\Spec}{\textup{Spec}}
\DeclareMathOperator{\Specsh}{\underline{\textup{Spec}}}
\DeclareMathOperator{\Stab}{\textup{Stab}}
\DeclareMathOperator{\Supp}{\textup{Supp}}
\DeclareMathOperator{\Sym}{\textup{Sym}}
\DeclareMathOperator{\TMod}{\textup{TMod}}
\DeclareMathOperator{\Top}{\textup{Top}}
\DeclareMathOperator{\Tor}{\textup{Tor}}
\DeclareMathOperator{\Vect}{\textup{Vect}}
\DeclareMathOperator{\alt}{\textup{ht}}
\DeclareMathOperator{\car}{\textup{char}}
\DeclareMathOperator{\codim}{\textup{codim}}
\DeclareMathOperator{\degtr}{\textup{degtr}}
\DeclareMathOperator{\depth}{\textup{depth}}
\DeclareMathOperator{\divis}{\textup{div}}
\DeclareMathOperator{\et}{\textup{et}}
\DeclareMathOperator{\ffpSch}{\textup{ffpSch}}
\DeclareMathOperator{\h}{\textup{h}}
\DeclareMathOperator{\ilim}{\displaystyle{\lim_{\longrightarrow}}}
\DeclareMathOperator{\ind}{\textup{ind}}
\DeclareMathOperator{\indim}{\textup{inj dim}}
\DeclareMathOperator{\lf}{\textup{LF}}
\DeclareMathOperator{\op}{\textup{op}}
\DeclareMathOperator{\ord}{\textup{ord}}
\DeclareMathOperator{\pd}{\textup{pd}}
\DeclareMathOperator{\plim}{\displaystyle{\lim_{\longleftarrow}}}
\DeclareMathOperator{\pr}{\textup{pr}}
\DeclareMathOperator{\pt}{\textup{pt}}
\DeclareMathOperator{\rk}{\textup{rk}}
\DeclareMathOperator{\tr}{\textup{tr}}
\DeclareMathOperator{\type}{\textup{r}}
\DeclareMathOperator*{\colim}{\textup{colim}}
\usepackage[colorinlistoftodos]{todonotes}
%\usepackage[inline]{showlabels}
\usepackage{enumitem}
\setenumerate[1]{label={\arabic*)}} % Global setting
\usepackage[a4paper]{geometry}

\author[Tonini]{Fabio Tonini}

\address[Tonini]{Humboldt University of Berlin\\Unter den Linden 6\\
10099 Berlin\\Germany}
\email{tonini@mathematik.hu-berlin.de}

\makeatother

\usepackage{babel}
  \providecommand{\corollaryname}{Corollary}
  \providecommand{\definitionname}{Definition}
  \providecommand{\examplename}{Example}
  \providecommand{\lemmaname}{Lemma}
  \providecommand{\propositionname}{Proposition}
  \providecommand{\remarkname}{Remark}
  \providecommand{\theoremname}{Theorem}
\providecommand{\theoremname}{Theorem}

\begin{document}

\title{Sheafification functors and Tannaka's reconstruction}

\maketitle
\global\long\def\A{\mathbb{A}}


\global\long\def\Ab{(\textup{Ab})}


\global\long\def\C{\mathbb{C}}


\global\long\def\Cat{(\textup{cat})}


\global\long\def\Di#1{\textup{D}(#1)}


\global\long\def\E{\mathcal{E}}


\global\long\def\F{\mathbb{F}}


\global\long\def\GCov{G\textup{-Cov}}


\global\long\def\Gcat{(\textup{Galois cat})}


\global\long\def\Gfsets#1{#1\textup{-fsets}}


\global\long\def\Gm{\mathbb{G}_{m}}


\global\long\def\GrCov#1{\textup{D}(#1)\textup{-Cov}}


\global\long\def\Grp{(\textup{Grps})}


\global\long\def\Gsets#1{(#1\textup{-sets})}


\global\long\def\HCov{H\textup{-Cov}}


\global\long\def\MCov{\textup{D}(M)\textup{-Cov}}


\global\long\def\MHilb{M\textup{-Hilb}}


\global\long\def\N{\mathbb{N}}


\global\long\def\PGor{\textup{PGor}}


\global\long\def\PGrp{(\textup{Profinite Grp})}


\global\long\def\PP{\mathbb{P}}


\global\long\def\Pj{\mathbb{P}}


\global\long\def\Q{\mathbb{Q}}


\global\long\def\RCov#1{#1\textup{-Cov}}


\global\long\def\RR{\mathbb{R}}


\global\long\def\Sch{\textup{Sch}}


\global\long\def\WW{\textup{W}}


\global\long\def\Z{\mathbb{Z}}


\global\long\def\acts{\curvearrowright}


\global\long\def\alA{\mathscr{A}}


\global\long\def\alB{\mathscr{B}}


\global\long\def\arr{\longrightarrow}


\global\long\def\arrdi#1{\xlongrightarrow{#1}}


\global\long\def\catC{\mathscr{C}}


\global\long\def\catD{\mathscr{D}}


\global\long\def\catF{\mathscr{F}}


\global\long\def\catG{\mathscr{G}}


\global\long\def\comma{,\ }


\global\long\def\covU{\mathcal{U}}


\global\long\def\covV{\mathcal{V}}


\global\long\def\covW{\mathcal{W}}


\global\long\def\duale#1{{#1}^{\vee}}


\global\long\def\fasc#1{\widetilde{#1}}


\global\long\def\fsets{(\textup{f-sets})}


\global\long\def\iL{r\mathscr{L}}


\global\long\def\id{\textup{id}}


\global\long\def\la{\langle}


\global\long\def\odi#1{\mathcal{O}_{#1}}


\global\long\def\ra{\rangle}


\global\long\def\set{(\textup{Sets})}


\global\long\def\sets{(\textup{Sets})}


\global\long\def\shA{\mathcal{A}}


\global\long\def\shB{\mathcal{B}}


\global\long\def\shC{\mathcal{C}}


\global\long\def\shD{\mathcal{D}}


\global\long\def\shE{\mathcal{E}}


\global\long\def\shF{\mathcal{F}}


\global\long\def\shG{\mathcal{G}}


\global\long\def\shH{\mathcal{H}}


\global\long\def\shI{\mathcal{I}}


\global\long\def\shJ{\mathcal{J}}


\global\long\def\shK{\mathcal{K}}


\global\long\def\shL{\mathcal{L}}


\global\long\def\shM{\mathcal{M}}


\global\long\def\shN{\mathcal{N}}


\global\long\def\shO{\mathcal{O}}


\global\long\def\shP{\mathcal{P}}


\global\long\def\shQ{\mathcal{Q}}


\global\long\def\shR{\mathcal{R}}


\global\long\def\shS{\mathcal{S}}


\global\long\def\shT{\mathcal{T}}


\global\long\def\shU{\mathcal{U}}


\global\long\def\shV{\mathcal{V}}


\global\long\def\shW{\mathcal{W}}


\global\long\def\shX{\mathcal{X}}


\global\long\def\shY{\mathcal{Y}}


\global\long\def\shZ{\mathcal{Z}}


\global\long\def\st{\ | \ }


\global\long\def\stA{\mathcal{A}}


\global\long\def\stB{\mathcal{B}}


\global\long\def\stC{\mathcal{C}}


\global\long\def\stD{\mathcal{D}}


\global\long\def\stE{\mathcal{E}}


\global\long\def\stF{\mathcal{F}}


\global\long\def\stG{\mathcal{G}}


\global\long\def\stH{\mathcal{H}}


\global\long\def\stI{\mathcal{I}}


\global\long\def\stJ{\mathcal{J}}


\global\long\def\stK{\mathcal{K}}


\global\long\def\stL{\mathcal{L}}


\global\long\def\stM{\mathcal{M}}


\global\long\def\stN{\mathcal{N}}


\global\long\def\stO{\mathcal{O}}


\global\long\def\stP{\mathcal{P}}


\global\long\def\stQ{\mathcal{Q}}


\global\long\def\stR{\mathcal{R}}


\global\long\def\stS{\mathcal{S}}


\global\long\def\stT{\mathcal{T}}


\global\long\def\stU{\mathcal{U}}


\global\long\def\stV{\mathcal{V}}


\global\long\def\stW{\mathcal{W}}


\global\long\def\stX{\mathcal{X}}


\global\long\def\stY{\mathcal{Y}}


\global\long\def\stZ{\mathcal{Z}}


\global\long\def\then{\ \Longrightarrow\ }


\global\long\def\L{\textup{L}}


\global\long\def\l{\textup{l}}


\begin{comment}
\global\long\def\Add{\textup{Add}}


\global\long\def\Aff{\textup{Aff}}


\global\long\def\Alg{\textup{Alg}}


\global\long\def\Ann{\textup{Ann}}


\global\long\def\Arr{\textup{Arr}}


\global\long\def\Art{\textup{Art}}


\global\long\def\Ass{\textup{Ass}}


\global\long\def\Aut{\textup{Aut}}


\global\long\def\Autsh{\underline{\textup{Aut}}}


\global\long\def\Bi{\textup{B}}


\global\long\def\CAdd{\textup{CAdd}}


\global\long\def\CAlg{\textup{CAlg}}


\global\long\def\CMon{\textup{CMon}}


\global\long\def\CPMon{\textup{CPMon}}


\global\long\def\CRings{\textup{CRings}}


\global\long\def\CSMon{\textup{CSMon}}


\global\long\def\CaCl{\textup{CaCl}}


\global\long\def\Cart{\textup{Cart}}


\global\long\def\Cl{\textup{Cl}}


\global\long\def\Coh{\textup{Coh}}


\global\long\def\Coker{\textup{Coker}}


\global\long\def\Cov{\textup{Cov}}


\global\long\def\Der{\textup{Der}}


\global\long\def\Div{\textup{Div}}


\global\long\def\End{\textup{End}}


\global\long\def\Endsh{\underline{\textup{End}}}


\global\long\def\Ext{\textup{Ext}}


\global\long\def\Extsh{\underline{\textup{Ext}}}


\global\long\def\FAdd{\textup{FAdd}}


\global\long\def\FCoh{\textup{FCoh}}


\global\long\def\FGrad{\textup{FGrad}}


\global\long\def\FLoc{\textup{FLoc}}


\global\long\def\FMod{\textup{FMod}}


\global\long\def\FPMon{\textup{FPMon}}


\global\long\def\FRep{\textup{FRep}}


\global\long\def\FSMon{\textup{FSMon}}


\global\long\def\FVect{\textup{FVect}}


\global\long\def\Fibr{\textup{Fibr}}


\global\long\def\Fix{\textup{Fix}}


\global\long\def\Fl{\textup{Fl}}


\global\long\def\Fr{\textup{Fr}}


\global\long\def\Funct{\textup{Funct}}


\global\long\def\GAlg{\textup{GAlg}}


\global\long\def\GExt{\textup{GExt}}


\global\long\def\GHom{\textup{GHom}}


\global\long\def\GL{\textup{GL}}


\global\long\def\GMod{\textup{GMod}}


\global\long\def\GRis{\textup{GRis}}


\global\long\def\GRiv{\textup{GRiv}}


\global\long\def\Gal{\textup{Gal}}


\global\long\def\Gl{\textup{Gl}}


\global\long\def\Grad{\textup{Grad}}


\global\long\def\Hilb{\textup{Hilb}}


\global\long\def\Hl{\textup{H}}


\global\long\def\Hom{\textup{Hom}}


\global\long\def\Homsh{\underline{\textup{Hom}}}


\global\long\def\ISym{\textup{Sym}^{*}}


\global\long\def\Imm{\textup{Im}}


\global\long\def\Irr{\textup{Irr}}


\global\long\def\Iso{\textup{Iso}}


\global\long\def\Isosh{\underline{\textup{Iso}}}


\global\long\def\Ker{\textup{Ker}}


\global\long\def\LAdd{\textup{LAdd}}


\global\long\def\LAlg{\textup{LAlg}}


\global\long\def\LMon{\textup{LMon}}


\global\long\def\LPMon{\textup{LPMon}}


\global\long\def\LRings{\textup{LRings}}


\global\long\def\LSMon{\textup{LSMon}}


\global\long\def\Left{\textup{L}}


\global\long\def\Lex{\textup{Lex}}


\global\long\def\Loc{\textup{Loc}}


\global\long\def\M{\textup{M}}


\global\long\def\ML{\textup{ML}}


\global\long\def\MLex{\textup{MLex}}


\global\long\def\Map{\textup{Map}}


\global\long\def\Mod{\textup{Mod}}


\global\long\def\Mon{\textup{Mon}}


\global\long\def\Ob{\textup{Ob}}


\global\long\def\Obj{\textup{Obj}}


\global\long\def\PDiv{\textup{PDiv}}


\global\long\def\PGL{\textup{PGL}}


\global\long\def\PML{\textup{PML}}


\global\long\def\PMLex{\textup{PMLex}}


\global\long\def\PMon{\textup{PMon}}


\global\long\def\Pic{\textup{Pic}}


\global\long\def\Picsh{\underline{\textup{Pic}}}


\global\long\def\Pro{\textup{Pro}}


\global\long\def\Proj{\textup{Proj}}


\global\long\def\QAdd{\textup{QAdd}}


\global\long\def\QAlg{\textup{QAlg}}


\global\long\def\QCoh{\textup{QCoh}}


\global\long\def\QMon{\textup{QMon}}


\global\long\def\QPMon{\textup{QPMon}}


\global\long\def\QRings{\textup{QRings}}


\global\long\def\QSMon{\textup{QSMon}}


\global\long\def\R{\textup{R}}


\global\long\def\Rep{\textup{Rep}}


\global\long\def\Rings{\textup{Rings}}


\global\long\def\Riv{\textup{Riv}}


\global\long\def\SFibr{\textup{SFibr}}


\global\long\def\SMLex{\textup{SMLex}}


\global\long\def\SMex{\textup{SMex}}


\global\long\def\SMon{\textup{SMon}}


\global\long\def\SchI{\textup{SchI}}


\global\long\def\Sh{\textup{Sh}}


\global\long\def\Soc{\textup{Soc}}


\global\long\def\Spec{\textup{Spec}}


\global\long\def\Specsh{\underline{\textup{Spec}}}


\global\long\def\Stab{\textup{Stab}}


\global\long\def\Supp{\textup{Supp}}


\global\long\def\Sym{\textup{Sym}}


\global\long\def\TMod{\textup{TMod}}


\global\long\def\Top{\textup{Top}}


\global\long\def\Tor{\textup{Tor}}


\global\long\def\Vect{\textup{Vect}}


\global\long\def\alt{\textup{ht}}


\global\long\def\car{\textup{char}}


\global\long\def\codim{\textup{codim}}


\global\long\def\degtr{\textup{degtr}}


\global\long\def\depth{\textup{depth}}


\global\long\def\divis{\textup{div}}


\global\long\def\et{\textup{et}}


\global\long\def\ffpSch{\textup{ffpSch}}


\global\long\def\h{\textup{h}}


\global\long\def\ilim{{\displaystyle {\lim_{\longrightarrow}}}}


\global\long\def\ind{\textup{ind}}


\global\long\def\indim{\textup{inj dim}}


\global\long\def\lf{\textup{LF}}


\global\long\def\op{\textup{op}}


\global\long\def\ord{\textup{ord}}


\global\long\def\pd{\textup{pd}}


\global\long\def\plim{{\displaystyle {\lim_{\longleftarrow}}}}


\global\long\def\pr{\textup{pr}}


\global\long\def\pt{\textup{pt}}


\global\long\def\rk{\textup{rk}}


\global\long\def\tr{\textup{tr}}


\global\long\def\type{\textup{r}}


\global\long\def\colim{\textup{colim}}
\end{comment}



%\makeatletter 
%\providecommand\@dotsep{5} 
%\makeatother 
%\listoftodos\relax
\begin{abstract}
We introduce ``sheafification'' functors from categories of (lax
monoidal) linear functors to categories of quasi-coherent sheaves
(of algebras) of stacks. They generalize the homogeneous sheafification
of graded modules for projective schemes and have applications in
the theory of non-abelian Galois covers and of Cox rings and homogeneous
sheafification functors. Moreover, using this theory, we prove a non-neutral
form of Tannaka's reconstruction, extending the classical correspondence
between torsors and strong monoidal functors.
\end{abstract}

\section*{Introduction}

If $G$ is an affine group scheme over a field $k$, classical Tannaka's
reconstruction problem consists in reconstructing the group $G$ from
$\Rep^{G}k$, its category of finite representations: if $F\colon\Rep^{G}k\arr\Vect k$
is the forgetful functor then $G$ is canonically isomorphic to the
sheaf of automorphisms of $F$ (opportunely defined, see \cite[Proposition 2.8]{Deligne1982}).
More generally one can recover the stack $\Bi G$ of $G$-torsors
by looking at fiber functors. Given a ring $A$ denote by $\Loc A$
the category of locally free sheaves of finite rank over $A$, that
is finitely generated projective $A$-modules. If $\SMon_{k}^{G}$
is the stack over $k$ whose fiber over a $k$-algebra $A$ is the
category of $k$-linear, strong monoidal and exact functors $\Gamma\colon\Rep^{G}k\arr\Loc A$
then the functor \begin{align}\label{tannaka functor}
\Bi G\arr\SMon_{k}^{G}\comma(\Spec A\arrdi s\Bi G)\longmapsto(s_{|\Rep^{G}k}^{*}\colon\Rep^{G}k\arr\Loc A) \tag{$\star$}
\end{align}is an equivalence of categories (see \cite[Theorem 3.2]{Deligne1982}).
Here $\Rep^{G}k$ is thought of as the category of locally free sheaves
of finite rank on $\Bi G$. 

The above functor can be defined in a way more general context. Let
us introduce some notations and definitions. We fix a base commutative
ring $R$ and a category fibered in groupoids $\stX$ over $R$. We
say that $\stX$ is \emph{pseudo-algebraic} (resp. \emph{quasi-compact})
if there exists a scheme (resp. affine scheme) $X$ and a map $X\arr\stX$
representable by fpqc covering of algebraic spaces. We denote by $\QCoh\stX$
and $\Loc\stX$ the categories of quasi-coherent sheaves and locally
free sheaves of finite rank respectively (see Section \ref{sec: Preliminaries on sheaves}).
Given a full subcategory $\shC$ of $\QCoh\stX$ we say that $\shC$
generates $\QCoh\stX$ if all quasi-coherent sheaves on $\stX$ are
quotient of a (possibly infinite) direct sum of sheaves in $\shC$.
We say that $\stX$ satisfies (or has) the \emph{resolution property}
if $\Loc\stX$ generates $\QCoh\stX$.

Given a monoidal and additive full subcategory $\shC$ of $\Loc\stX$
and a category fibered in groupoids $\stY$ over $R$ define $ $$\text{Fib}_{\stX,\shC}(\stY)$
as the category of\emph{ }$R$-linear and strong monoidal functors
$\Gamma\colon\shC\arr\Loc\stY$ which are exact on right exact sequences
in $\shC$ (in the ambient abelian category $\QCoh\stX$). Denote
also by $\text{Fib}_{\stX,\shC}$ the stack over $R$ whose fiber
over an $R$-algebra $A$ is $\text{Fib}_{\stX,\shC}(\Spec A)$ and
by $\shP_{\shC}$ the functor 
\[
\shP_{\shC}\colon\stX\arr\text{Fib}_{\stX,\shC}\comma(\Spec A\arrdi s\stX)\longmapsto(s_{|\shC}^{*}\colon\shC\arr\Loc A)
\]
The functor (\ref{tannaka functor}) is obtained by taking $R=k$,
$\stX=\Bi G$ and $\shC=\Loc\stX$. We prove the following non-neutral
form of Tannaka's reconstruction.
\begin{thm*}
[\ref{thm: Tannaka reconstrion for Loc X}, \ref{cor:Lurie s point of view}]
Let $\stX$ be a quasi-compact stack over $R$ for the fpqc topology
with quasi-affine diagonal and $\shC\subseteq\Loc\stX$ be a full,
additive and monoidal subcategory with duals generating $\QCoh\stX$. 

If $\Gamma\colon\shC\arr\Mod A,$ where $A$ is an $R$-algebra, is
an $R$-linear, contravariant and strong monoidal functor such that
$\Gamma$, as well as $\Gamma\otimes_{A}k$ for all geometric points
$\Spec k\arr\Spec A$, is left exact on right exact sequences in $\shC$
then there exists $\Spec A\arrdi s\stX$ such that $\Gamma\simeq(s^{*})^{\vee}$. 

The functor $\shP_{\shC}\colon\stX\arr\text{Fib}_{\stX,\shC}$ is
an equivalence of stacks and, if $\stY$ is a category fibered in
groupoids, the functor 
\[
\Hom(\stY,\stX)\arr\text{Fib}_{\stX,\shC}(\stY)\comma(\stY\arrdi f\stX)\longmapsto f_{|\shC}^{*}\colon\shC\arr\Loc(\stY)
\]
is an equivalence of categories.
\end{thm*}
Notice that the two conclusions in the last statement are equivalent.
In the case $\shC=\Loc\stX$, the functor $\shP_{\Loc(\stX)}$ has
already been proved to be an equivalence in the neutral case, that
is $\stX=\Bi_{R}G$, where $G$ is a flat and affine group scheme
over $R$ (see \cite[Theorem 1.2]{Broshi2013}, where $R$ is a Dedekind
domain, and \cite[Theorem 1.3.2]{Schappi2011} for general rings $R$),
for particular quotient stacks over a field (see \cite{Savin2006}
and \ref{rem: cases for resolution property}) and for quasi-compact
and quasi-separated schemes (see \cite[Proposition 1.8]{Brandenburg2012}).
We also show an almost converse of Theorem above:
\begin{thm*}
{[}\ref{thm:when FibX,C has the resolution property}{]} Let $\stX$
be a quasi-compact category fibered in groupoids over $R$ admitting
a surjective (on equivalence classes of geometric points) map $X\arr\stX$
from a scheme whose connected components are open (e.g. a connected
or Noetherian algebraic stack) and let $\shC\subseteq\Loc\stX$ be
a full monoidal subcategory with duals such that $\Sym^{n}\E\in\shC$
for $n\in\N$ and $\E\in\shC$ if $\E$ has local rank not invertible
in $R$ (e.g. $\shC=\Loc(\stX)$ or $\shC$ consists of invertible
sheaves). Then $\text{Fib}_{\stX,\shC}$ is a quasi-compact stack
in groupoids for the fpqc topology over $R$ with affine diagonal
and with a collection of tautological locally free sheaves $\{\shG_{\E}\}_{\E\in\shC}$
generating $\QCoh(\text{Fib}_{\stX,\shC})$ and such that $\shP_{\shC}^{*}\shG_{\E}\simeq\E$.
\end{thm*}
For instance it follows that $\stX\arr\text{Fib}_{\stX,\Loc(\stX)}$
is universal among the maps $w\colon\stX\arr\stY$ where $\stY$ is
a quasi-compact stack over $R$ with quasi-affine diagonal and the
resolution property.

There are also variants of theory above where $\Loc(-)$ is replaced
by $\Coh(-)$ or $\QCoh(-)$ or the derived category $D(-)$ (see
\cite{Lurie2004,Schappi2012,Brandenburg2012,Brandenburg2014,Bhatt2014}).
Although not explicitly stated elsewhere, for stacks with the resolution
property and with affine diagonal (which is automatic in the algebraic
case, see \cite{Totaro2002}) those results and the fact that $\shP_{\Loc(\stX)}$
is an equivalence can be proved to be equivalent: one can pass from
quasi-coherent sheaves to locally free sheaves via dualizable objects
and, for the converse, extend functors from $\Loc(-)$ to $\QCoh(-)$
following the proof of \cite[Corollary 3.2]{Bhatt2014}. We complete
this picture by showing that in general the resolution property implies
the affineness of the diagonal (see \ref{cor:quasi-affine and resolution implies affine}).
One of the ingredients in the proof is the classification of quasi-compact
stacks whose quasi-coherent sheaves are generated by global sections,
called pseudo-affine. In the algebraic case those coincide with quasi-affine
schemes (see \cite[Proposition 3.1]{Gross2013}), while, in general,
we prove they are (arbitrary) intersection of quasi-compact open subschemes
(thought of as sheaves) of affine schemes (see \ref{prop:arbitrary intersection of quasi-compact open is pseudo-affine}).

The proof of Tannaka's reconstruction we present does not reduce to
the case of quasi-coherent sheaves as explained above but it follows
a different path. It is obtained by developing a theory of sheafification
functors which I think is interesting on its own and it is the heart
of the paper. In what follows $A$ will denote an $R$-algebra, $\stX$
a category fibered in groupoids over $R$ and $\shC$ a full subcategory
of $\QCoh\stX$. The idea is simple: if $s\colon\Spec A\arr\stX$
is a map (say quasi-affine so that $s_{*}$ preserves quasi-coherency)
then we have natural isomorphisms of $A$-modules
\[
s^{*}\E\simeq\Hom_{A}(s^{*}\duale{\E},A)\simeq\Hom_{\stX}(\duale{\E},s_{*}\odi{\Spec A})
\]
By passing from covariant functors to \emph{contravariant} ones we
can always define 
\[
\Omega^{\shF}\colon\shC\arr\Mod A\comma\Omega_{\E}^{\shF}=\Hom_{\stX}(\E,\shF)\text{ for }\shF\in\QCoh(\stX\times A)
\]
which are $R$-linear contravariant functors. If $\L_{R}(\shC,A)$
is the category of contravariant $R$-linear functors $\shC\arr\Mod A$
we obtain a functor
\[
\Omega^{*}\colon\QCoh(\stX\times A)\arr\L_{R}(\shC,A)
\]
Quite surprisingly (since we started from group schemes as motivation)
we recover also this well known situation (see \ref{example: quasi-projective case}
for details). If $X$ is a quasi-projective and quasi-compact scheme
over $R$ with very ample invertible sheaf $\odi X(1)$ consider $\shC_{X}=\{\odi X(n)\}_{n\in\Z}$,
set $S_{X}$ for the homogeneous coordinate ring of $(X,\odi X(1))$
and use $\GMod(-)$ to denote the category of graded modules. We have
that $\L_{R}(\shC_{X},A)$ is equivalent to $\GMod(S_{X}\otimes_{R}A)$
and $\Omega^{*}$ corresponds to
\[
\Gamma_{*}\colon\QCoh(X\times A)\arr\GMod(S_{X}\otimes_{R}A)\comma\Gamma_{*}(\shF)=\bigoplus_{n\in\Z}\Hl^{0}(X,\shF(n))
\]
It is a classical result, at least over a field, that the above functor
is fully faithful. Thus in general we can ask under what conditions
the functor $\Omega^{*}\colon\QCoh(\stX\times A)\arr\L_{R}(\shC,A)$
is fully faithful as well. Another property that the functor $\Gamma_{*}$
has is the existence of a left adjoint, namely the homogeneous sheafification
functor $\fasc -\colon\GMod(S_{X}\otimes_{R}A)\arr\QCoh(X\times A)$.
It turns out that also $\Omega^{*}\colon\QCoh(\stX\times A)\arr\L_{R}(\shC,A)$
has a left adjoint $\shF_{*,\shC}\colon\L_{R}(\shC,A)\arr\QCoh(\stX\times A)$
when $\shC$ is essentially small: by analogy we call this functor
a \emph{sheafification }functor.

In the context of Tannaka's reconstruction we also have that sheaves
of algebras correspond to (lax) monoidal functors $\Loc^{G}R\arr\Mod A$.
This is true in general. If $\shC$ is a monoidal subcategory of $\QCoh\stX$,
$\ML_{R}(\shC,A)$ denotes the category of $R$-linear, contravariant
and monoidal functors $\shC\arr\Mod A$ and $\QAlg(\stX\times A)$
the category of quasi-coherent sheaves of algebras on $\stX$ then
$\Omega^{*}$ extends to a functor $\Omega^{*}\colon\QAlg(\stX\times A)\arr\ML_{R}(\shC,A)$
and, if $\shC$ is essentially small, $\shF_{*,\shC}$ extends to
$\alA_{*,\shC}\colon\ML_{R}(\shC,A)\arr\QAlg(\stX\times A)$, which
is still a left adjoint of $\Omega^{*}$. 

Coming back to the fully faithfulness of $\Omega^{*}$ we prove the
following.
\begin{thm*}
[\ref{prop:excatness properties of Omega* and shF*},\ref{thm:general theorem for equivalences}]
Let $\stX$ be a pseudo-algebraic category fibered in groupoids over
$R$ and $\shC\subseteq\QCoh\stX$ be a full subcategory generating
$\QCoh\stX$. Then the functor $\Omega^{*}\colon\QCoh(\stX\times A)\arr\L_{R}(\shC,A)$
is fully faithful and, if $\shC$ is essentially small, then $\shF_{*,\shC}\colon\L_{R}(\shC,A)\arr\QCoh(\stX\times A)$
is exact and the natural map $\shG\arr\shF_{\Omega^{\shG},\shC}$
is an isomorphism.
\end{thm*}
If $X$ is quasi-projective then it is a classical fact that $\shC_{X}=\{\odi X(n)\}_{n\in\Z}$
generates $\QCoh X$ and thus we recover the classical properties
of $\Gamma_{*}\colon\QCoh(X\times A)\arr\GMod(S_{X}\otimes_{R}A)$.
Theorem above when $\shC$ consists of a single object is a rephrasing
of classical Gabriel-Popescu's theorem for the category $\QCoh\stX$
(see \ref{thm: Gabriel Popescu's theorem}). When $\shC$ is monoidal
and generates $\QCoh\stX$ we also have that $\Omega^{*}\colon\QAlg(\stX\times A)\arr\ML_{R}(\shC,A)$
is fully faithful (see \ref{thm:fundamental for sheaf of algebras and monoidal functors}).

The second problem we address is to describe the essential image of
$\Omega^{*}$. The main idea behind this description is that $\Hom_{\stX}(-,\shF)$
for $\shF\in\QCoh\stX$ is a left exact functor. Since the domain
$\shC$ of the functors $\Omega^{\shF}$ is not abelian we need an
ad hoc definition of exactness. A \emph{test sequence }in $\shC$
is an exact sequence 
\[
\shT_{*}\colon\bigoplus_{k\in K}\E_{k}\arrdi{\alpha}\bigoplus_{i\in I}\E_{i}\arr\E\arr0\text{ where }\E,\E_{i},\E_{k}\in\shC\text{ and }I,J\text{ are sets}
\]
such that $\alpha(\E_{k})$ is contained in a finite sum for all $k\in K$.
A test sequence is called finite if $K$ and $I$ are finite sets.
Given $\Gamma\in\L_{R}(\shC,A)$ we will say that $\Gamma$ is exact
on a test sequence $\shT_{*}$ in $\shC$ if the complex of $A$-modules
(see \ref{def: left exact functors})
\[
0\arr\Gamma_{\E}\arr\prod_{i\in I}\Gamma_{\E_{i}}\arr\prod_{k\in K}\Gamma_{\E_{k}}
\]
is exact. We denote by $\Lex_{R}(\shC,A)$ the full subcategory of
$\L_{R}(\shC,A)$ of functors which are exact an all test sequences.
We have the following:
\begin{thm*}
[\ref{thm:general theorem for equivalences}, \ref{ prop: Lex and left exact functors}]
Let $\stX$ be a pseudo-algebraic category fibered in groupoids over
$R$ and $\shC\subseteq\QCoh\stX$ be a full subcategory generating
$\QCoh\stX$. Then $\Lex_{R}(\shC,A)$ is the essential image of the
(fully faithful) functor $\Omega^{*}\colon\QCoh(\stX\times A)\arr\L_{R}(\shC,A)$.
If $\stX$ is quasi-compact and all sheaves in $\shC$ are finitely
presented then $\Lex_{R}(\shC,A)$ is the subcategory of $\L_{R}(\shC,A)$
of functors which are exact on finite test sequences.
\end{thm*}
In particular, when $\shC$ is essentially small, $\Omega^{*}\colon\QCoh(\stX\times A)\arr\Lex_{R}(\shC,A)$
and $\shF_{*,\shC}\colon\Lex_{R}(\shC,A)\arr\QCoh(\stX\times A)$
are quasi-inverses of each other. The two theorems above apply in
the following situations (see \ref{thm:equivalence for QCoh X}, \ref{thm:equivalence for noetherian algebraic stacks}
and \ref{thm:equivalence for Loc X}): if $\shC=\QCoh\stX$ then $\Lex_{R}(\QCoh\stX,A)$
is the category of contraviant, $R$-linear and left exact functors
$\QCoh\stX\arr\Mod A$ which transform direct sums into products,
if $\shC=\Coh(\stX)$ (resp. $\shC=\Loc(\stX)$) and $\stX$ is a
notherian algebraic stack (resp. $\stX$ is quasi-compact and has
the resolution property) then $\Lex_{R}(\shC,A)$ is the category
of contraviant, $R$-linear and left exact functors $\shC\arr\Mod A$.

When $\shC$ is essentially small there is another cohomological characterization
of the functors in $\Lex_{R}(\shC,A)$. A collection of maps $\shU=\{\E_{i}\arr\E\}$
in $\shC$ is called jointly surjective if the map $\bigoplus_{i\in I}\E_{i}\arr\E$
is surjective. Given such a collection $\shU$ we set $\Delta_{\shU}=\Imm(\bigoplus_{i\in I}\Omega^{\E_{i}}\arr\Omega^{\E})\in\L_{R}(\shC,R)$.
Denote by $\shC^{\oplus}$ the subcategory of $\QCoh\stX$ consisting
of all possible finite direct sums of sheaves in $\shC$. We have:
\begin{thm*}
[\ref{prop:test sequences and Extone}, \ref{ prop: Lex and left exact functors}
and \ref{prop:Lex and Grothendieck topology}] Let $\stX$ be a pseudo-algebraic
category fibered in groupoids over $R$ and $\shC\subseteq\QCoh\stX$
be a full and essentially small subcategory generating $\QCoh\stX$.
Then $\Lex_{R}(\shC,A)$ is the full subcategory of $\L_{R}(\shC,A)$
of functors $\Gamma$ satisfying 
\[
\Hom_{\L_{R}(\shC,R)}(\Omega^{\E}/\Delta_{\shU},\Gamma)=\Ext_{\L_{R}(\shC,R)}^{1}(\Omega^{\E}/\Delta_{\shU},\Gamma)=0
\]
for all jointly surjective collections of maps $\shU=\{\E_{i}\arr\E\}_{i\in I}$
in $\shC$. If $\stX$ is quasi-compact and the sheaves in $\shC$
are finitely presented we can consider only finite collections $\shU$.

We have $\Lex_{R}(\shC^{\oplus},A)\simeq\Lex_{R}(\shC,A)$ via the
restriction $\shC\arr\shC^{\oplus}$ and, if $\shC$ is additive and
$\shJ$ is the smallest Grothendieck topology on $\shC$ containing
the sieves $\Delta_{\shU}$ for all jointly surjective collections
$\shU=\{\E_{i}\arr\E\}_{i\in I}$ in $\shC$, then $\Lex_{R}(\shC,A)$
coincides with the category of sheaves of $A$-modules $\shC^{\op}\arr\Mod A$
on the site $(\shC,\shJ)$ which are $R$-linear.
\end{thm*}
Besides Tannaka's reconstruction problem, theory above has two other
applications. The first, which is also the original motivation, is
the theory of Galois cover. More precisely in my Ph.D. thesis \cite{Tonini2013}
I have worked out theory above in the case $\stX=\Bi G$ and $\shC=\Loc\stX$,
where $G$ is a finite, flat and finitely presented group scheme over
$R$. Notice that $\Bi G$ satisfies the resolution property in this
case (see \ref{rem: resolution property for BG cases}). The proof
presented in \cite{Tonini2013} makes use of representation theory
and can not be generalized to arbitrary categories fibered in groupoids.
The goal was to look at Galois covers with group $G$ as particular
monoidal functors, as $G$-torsors can be thought of as particular
strong monoidal functors, and the motivation was the study of non-abelian
Galois covers, where a direct approach as in the abelian case (see
\cite{Tonini2011}) fails due to the complexity of the representation
theory.

A second application is to the theory of Cox rings and homogeneous
sheafifications. The idea is to consider $\shC_{H}=\{\shL\}_{\shL\in H}\subseteq\Loc\stX$
where $H$ is a subgroup of $\Pic\stX$. As in the projective case
we have a homogeneous coordinate ring
\[
S_{H}=\bigoplus_{\shL\in H}\Hl^{0}(\stX,\shL)
\]
(opportunely defined), $\L_{R}(\shC_{H},A)$ is equivalent to $\GMod(S_{H}\otimes_{R}A)$,
the category of $H$-graded $(S_{H}\otimes_{R}A)$-modules, $\Omega^{*}$
corresponds to
\[
\Gamma_{*}\colon\QCoh(\stX\times A)\arr\GMod(S_{H}\otimes_{R}A)\comma\Gamma_{*}(\shF)=\bigoplus_{\shL\in H}\Hl^{0}(\stX,\shF\otimes\shL)
\]
and its adjoint $\shF_{*,\shC_{H}}$ behaves like a homogeneous sheafification.
Moreover in more concrete geometric situations, e.g. when $\stX$
is a normal variety, there are analogous constructions for reflexive
sheaves of rank $1$. We expect that this theory covers all known
cases where $\Gamma_{*}$ is proved to be fully faithful (see for
instance \cite[Appendix of Chapter 6]{Cox2011} and \cite[Section 2]{Kajiwara1998}).

Applications above are not described in the present paper, but, hopefully,
they will be subjects of future ones.

The outline of the paper is the following. The first section introduces
the notion and the basic properties of quasi-coherent sheaves on fibered
categories, while the second one is a general study of sheafification
functors. In the third section we study the fully faithfulness and
the essential image of the functor $\Omega^{*}\colon\QCoh(\stX\times A)\arr\L_{R}(\shC,A)$.
In the fourth section we rewrite the results obtained in the case
of the stack of $G$-torsor in terms of the representation theory
of $G$ and finally, in the last section, we prove the non-neutral
Tannaka's reconstruction.


\section*{Notation}

In this paper we work over a base commutative, associative ring $R$
with unity. If not stated otherwise a fiber category will be a category
fibered in groupoids over $\Aff/R$, the category of affine schemes
over $\Spec R$, or, equivalently, the opposite of the category of
$R$-algebras. A scheme or an algebraic space $X$ over $\Spec R$
will be thought of as the fibered category of maps from an affine
scheme to $X$, denoted by $\Aff/X$. A map or morphism of fibered
categories is a functor over $\Aff/R$. Recall that by the $2$-Yoneda
lemma objects of a fibered category $\stX$ can be thought of as maps
$T\arr\stX$ from an affine scheme. An fpqc stack will be a stack
for the fpqc topology.

A map $f\colon\stX'\arr\stX$ of fibered categories is called representable
if for all maps $T\arr\stX$ from an affine scheme (or an algebraic
space) the fiber product $T\times_{\stX}\stX'$ is (equivalent to)
an algebraic space.

Given a flat and affine group scheme $G$ over $T$ we denote by $\Bi_{R}G$
the stack of $G$-torsors for the \emph{fpqc} topology, which is an
fpqc stack with affine diagonal. When $G\arr\Spec R$ is finitely
presented (resp. smooth) then $\Bi_{R}G$ coincides with the stack
of $G$-torsors for the fppf (resp. �tale) topology.

By a ``subcategory'' of a given category we mean a ``full subcategory''
if not stated otherwise.


\section*{Acknowledgements}

I would like to thank Jarod Alper, David Rydh, Daniel Sch�ppi, Mattia
Talpo and Angelo Vistoli for the useful conversations I had with them
and all the suggestions they gave me. 


\section{\label{sec: Preliminaries on sheaves} Preliminaries on sheaves and
fibered categories}

Let $\pi\colon\stX\arr\Aff/R$ be a fibered category. There is a functor
of rings $\odi{\stX}\colon\stX^{\op}\arr\sets$ defined by $\odi{\stX}(\xi)=\Hl^{0}(\odi{\pi(\xi)})$,
so that $(\stX,\odi{\stX})$ is a ringed category. A presheaf of $\odi{\stX}$-modules
on $\stX$ is a functor $\shF\colon\stX^{\text{op}}\arr\sets$ together
with an $\Hl^{0}(\odi{\pi(\xi)})$-module structure on $\shF(\xi)$
for all $\xi\in\stX$ such that, if $\xi\arr\eta$ is a map on $\stX$,
then the map $\shF(\eta)\arr\shF(\xi)$ is $\Hl^{0}(\odi{\pi(\eta)})$-linear.
Morphism of presheaves of $\odi{\stX}$-modules are natural transformations
respecting the module structures. We denote by $\Mod\odi{\stX}$ the
category of presheaves of $\odi{\stX}$-modules.
\begin{defn}
A quasi-coherent sheaf over $\stX$ is a presheaf of $\odi{\stX}$-modules
such that for all maps $\xi\arr\eta$ in $\stX$ the induced map
\[
\shF(\eta)\otimes_{\Hl^{0}(\odi{\pi(\eta)})}\Hl^{0}(\odi{\pi(\xi)})\arr\shF(\xi)
\]
is an isomorphism. We denote by $\QCoh\stX$ the full subcategory
of $\Mod\odi{\stX}$ of quasi-coherent sheaves.
\end{defn}
Notice that by fpqc descent of modules a quasi-coherent sheaf is a
sheaf for the fpqc topology of $\stX$.

If $f\colon\stY\arr\stX$ is a morphism of fibered categories and
$\shF\in\Mod\odi{\stX}$ we define $f^{*}\shF=\shF\circ f\colon\stY^{\op}\arrdi f\stX^{\op}\arrdi{\shF}\sets$.
This association defines a functor $f^{*}\colon\Mod\odi{\stX}\arr\Mod\odi{\stY}$,
called the pull-back functor, and restricts to a functor $f^{*}\colon\QCoh\stX\arr\QCoh\stY$.
Notice that $f^{*}\odi{\stX}=\odi{\stY}$ tautologically.

The category $\Mod\odi{\stX}$ is an abelian category and cokernels
of maps between quasi-coherent sheaves are again quasi-coherent and
thus cokernels in $\QCoh\stX$. Notice that, since we have restricted
our fiber categories to affine schemes, a map $f\colon\shF\arr\shG$
of quasi-coherent sheaves is an epimorphism if and only if it is pointwise
surjective, that is surjective in $\Mod\odi{\stX}$. The category
$\QCoh\stX$ is $R$-linear but it is unclear if it is abelian. Moreover
kernels in $\Mod\odi{\stX}$ of maps between quasi-coherent sheaves
are almost never quasi-coherent, essentially because pull-backs are
not left exact. There is a natural condition on $\stX$ which allows
us to prove that $\QCoh\stX$ is an $R$-linear abelian category.
\begin{defn}
An \emph{fpqc atlas }(or simply \emph{atlas}) of a fibered category
$\stX$ is a representable fpqc covering $X\arr\stX$ from a scheme.
A fiber category is called \emph{pseudo-algebraic }if it has an atlas,
it is called \emph{quasi-compact }if it has an atlas from an affine
scheme.

Let$f\colon\stY\arr\stX$ be a morphism of fibered categories. The
map $f$ is called \emph{pseudo-algebraic} (resp. \emph{quasi-compact})
if for all maps $T\arr\stX$ from a scheme (resp. quasi-compact scheme)
the fiber product $T\times_{\stX}\stY$ is pseudo-algebraic (resp.
quasi-compact). It is called \emph{quasi-separated} if the diagonal
$\stY\arr\stY\times_{\stX}\stY$ is quasi-compact.
\end{defn}
If $\stX$ is pseudo-algebraic then the diagonal $\Delta_{\stX}\colon\stX\arr\stX\times_{R}\stX$
is not necessarily representable. It is unclear whether this is true
or not if $\stX$ is a fpqc stack, because it is not known if algebraic
spaces satisfy effective descent along fpqc coverings.

Let $f\colon\stY\arr\stX$ be a map of fibered categories. If $\stX$
and $f$ are pseudo-algebraic then $\stY$ is pseudo-algebraic. If
$\stY$ is pseudo-algebraic and $\Delta_{\stX}$ is representable
then $f$ is pseudo-algebraic.

If $\catC$ is a (not full) subcategory of $\stX$ one can analogously
define presheaves of $(\odi{\stX})_{|\catC}$-modules and quasi-coherent
sheaves on $\catC$ just replacing all occurrences of $\stX$ with
$\catC$. We denote by $\Mod(\odi{\stX})_{|\catC}$ and $\QCoh\catC$
the resulting categories.
\begin{defn}
We define $\stX_{\text{fl}}$ (resp. $\stX_{\text{sm}}$, $\stX_{\text{et}}$)
as the (not full) subcategory of $\stX$ of objects $\xi\colon\Spec B\arr\stX$
which are representable and flat (resp. smooth, �tale) and the arrows
are morphisms in $\stX$ whose underlying map of affine schemes is
flat (resp. smooth, �tale).
\end{defn}
If $X\arr\stX$ is an fpqc atlas then by definition $R=X\times_{\stX}X$
is an algebraic space and the two projections $\shR\rightrightarrows X$
extends to a groupoid in algebraic spaces. We denote by $\QCoh(R\rightrightarrows X)$
the category of quasi-coherent sheaves on $R\rightrightarrows X$
(see \cite[Tag 0440]{SP014}). By standard arguments of fpqc descent
for modules we have: %\todo{Motivare??}
\begin{prop}
\label{prop:reducing to atlases for quasi-geometric categories} If
$\stX$ admits an fpqc (resp. smooth, �tale) atlas then the restriction
$\QCoh\stX\arr\QCoh\stX_{\text{fl}}$ (resp. $\QCoh\stX_{\text{sm}}$,
$\QCoh\stX_{\text{et}}$) is an equivalence of categories. If $f\colon X\arr\stX$
is an fpqc atlas then $f^{*}\colon\QCoh\stX\arr\QCoh X$ is faithful
and it induces an equivalence $\QCoh\stX\arr\QCoh(R\rightrightarrows X)$.
\end{prop}
We see that if $\stX$ is pseudo-algebraic then $\QCoh\stX$ is equivalent
to an $R$-linear abelian category, namely $\QCoh(R\rightrightarrows X)$.
Moreover if $\alpha\colon\shF\arr\shG$ is a map of quasi-coherent
sheaves then $\Ker(\alpha)$ is defined by taking $\Ker(\alpha_{|\stX_{\text{fl}}})\in\QCoh\stX_{\text{fl}}$,
which is just given by $\Ker(\alpha_{|\stX_{\text{fl}}})(\Spec B\arrdi{\xi}\stX)=\Ker(\alpha(\xi)\colon\shF(\xi)\arr\shG(\xi))$
for $\xi\in\stX_{\text{fl}}$, and then extending it to the whole
$\stX$. If $\stX$ is an algebraic stack or a scheme we see that
$\QCoh\stX$ is equivalent to the usual category of quasi-coherent
sheaves via an $R$-linear and exact functor. If it is given a subcategory
$\shD$ of $\QCoh\stX$, an exact sequence of sheaves in $\shD$ will
always be an exact sequence in $\QCoh\stX$ of sheaves belonging to
$\shD$.

We now deal with the problem of defining a right adjoint of a pull-back
functor. Given $\shF\in\Mod\odi{\stX}$ we define the global section
$\shF(\stX)=\Hom(\odi{\stX},\shF)$ of $\shF$, also denoted by $\Hl^{0}(\stX,\shF)$
or simply $\Hl^{0}(\shF)$, which is an $\odi{\stX}(\stX)$-module.
More generally given a map of fibered categories $g\colon\stZ\arr\stX$
we define $\shF(\stZ)=(g^{*}\shF)(\stZ)$. If $\stZ=\Spec B$ is affine
we will often write $\shF(B)$ instead of $\shF(\Spec B)$.

Let $f\colon\stY\arr\stX$ be a map of fibered categories. Given $\shG\in\Mod\odi{\stY}$
and an object $\xi\colon T\arr\stX$ of $\stX$ we define 
\[
(f_{p}\shG)(\xi)=\shG(T\times_{\stX}\stY)
\]
Given another object $\xi'\colon T'\arr\stX$ and a morphism $\beta\colon\xi'\arr\xi$
in $\stX$ there is an induced morphism $(f_{p}\shG)(\xi)\arr(f_{p}\shG)(\xi')$.
This data define a functor $f_{p}\colon\Mod\odi{\stY}\arr\Mod\odi{\stX}$
and we have:
\begin{prop}
Let $f\colon\stY\arr\stX$ be a map of fibered categories. Then $f_{p}$
is a right adjoint of $f^{*}$ and, if   \[   \begin{tikzpicture}[xscale=1.5,yscale=-1.2]     \node (A0_0) at (0, 0) {$\stY'$};     \node (A0_1) at (1, 0) {$\stY$};     \node (A1_0) at (0, 1) {$\stX'$};     \node (A1_1) at (1, 1) {$\stX$};     \path (A0_0) edge [->]node [auto] {$\scriptstyle{g'}$} (A0_1);     \path (A1_0) edge [->]node [auto] {$\scriptstyle{g}$} (A1_1);     \path (A0_1) edge [->]node [auto] {$\scriptstyle{f}$} (A1_1);     \path (A0_0) edge [->]node [auto] {$\scriptstyle{f'}$} (A1_0);   \end{tikzpicture}   \] 
is a $2$-cartesian diagram of fibered categories, there is an isomorphism
of functors
\[
g^{*}f_{p}\arr f'_{p}g'^{*}\colon\Mod\odi{\stY}\arr\Mod\odi{\stX'}
\]
If $f$ is affine then $f_{p}(\QCoh\stY)\subseteq\QCoh\stX$ and $(f_{p})_{|\QCoh\stY}\colon\QCoh\stY\arr\QCoh\stX$
is right adjoint to $f^{*}\colon\QCoh\stX\arr\QCoh\stY$.\end{prop}
\begin{proof}
The adjunction between $f^{*}$ and $f_{p}$ can be found in \cite[Tag 00XF]{SP014}.
With notation from this reference, we have that $_{\eta}^{f}\shI\simeq T\times_{\stX}\stY$
for an object $\eta\colon T\arr\stX$. Moreover if $\shF\in\Mod\odi{\stY}$
then $\shF(T\times_{\stX}\stY)$ is the limit of $\shF_{|T\times_{\stX}\stY}$
over the whole $T\times_{\stX}\stY$. Thus what is denoted by $_{p}f$
is easily seen to be equivalent to our $f_{p}$ if we take limits
of $R$-modules and not of sets. The isomorphism for the base change
is tautological. For the last claim we can assume that $\stX$ is
an affine scheme in which case the result follows because (usual)
push-forwards commutes with arbitrary base changes.
\end{proof}
In general $f_{p}$ does not preserve quasi-coherent sheaves, even
if $f$ is a proper map of schemes. To get a right adjoint of pullback
we have to require more. 
\begin{defn}
A pseudo-algebraic map $f\colon\stY\arr\stX$ of fibered categories
is called \emph{flat} if given an object $\xi\colon T\arr\stX$ of
$\stX$ and an atlas $V\arr T\times_{\stX}\stY$ the resulting map
$V\arr T$ is flat.\end{defn}
\begin{prop}
\label{prop:properties of push-forward} Let $f\colon\stY\arr\stX$
be a quasi-compact and quasi-separated map of pseudo-algebraic fibered
categories. Then the composition $\QCoh\stY\arrdi{f_{p}}\Mod\odi{\stX}\arr\Mod(\odi{\stX})_{|\stX_{\text{fl}}}$
has values in $\QCoh\stX_{\text{fl}}$. The induced map $f_{*}\colon\QCoh\stY\arr\QCoh\stX$
is a right adjoint of $f^{*}\colon\QCoh\stX\arr\QCoh\stY$. If   \[   \begin{tikzpicture}[xscale=1.5,yscale=-1.2]     \node (A0_0) at (0, 0) {$\stY'$};     \node (A0_1) at (1, 0) {$\stY$};     \node (A1_0) at (0, 1) {$\stX'$};     \node (A1_1) at (1, 1) {$\stX$};     \path (A0_0) edge [->]node [auto] {$\scriptstyle{g'}$} (A0_1);     \path (A1_0) edge [->]node [auto] {$\scriptstyle{g}$} (A1_1);     \path (A0_1) edge [->]node [auto] {$\scriptstyle{f}$} (A1_1);     \path (A0_0) edge [->]node [auto] {$\scriptstyle{f'}$} (A1_0);   \end{tikzpicture}   \] 
is a $2$-cartesian diagram of fibered categories with $\stX'$ pseudo-algebraic
then $\stY'$ is pseudo-algebraic, $f'$ is quasi-compact and quasi-separated
and there is a natural transformation of functors 
\[
g^{*}f_{*}\arr f'_{*}g'^{*}\colon\QCoh\stX\arr\QCoh\stY'
\]
which is an isomorphism if $g$ is flat.\end{prop}
\begin{proof}
Consider the $2$-Cartesian diagram in the statement. The diagonal
of $f'$ is quasi-compact because it is base change of the diagonal
of $f$. To see that $f_{p}(\shF)_{|\stX_{\text{fl}}}$ is quasi-coherent
for $\shF\in\QCoh\stY$, we can assume $\stX=\Spec B$ affine and
that $\stY$ is quasi-compact with quasi-compact diagonal. If $U=\Spec A\arr\stY$
is a fpqc atlas, it follows that $R=U\times_{\stY}U$ is a quasi-compact
algebraic space. By covering $R$ by finitely many affine schemes
$\Spec A_{i}$ we can write $\shF(\stY)$ as kernel of a map $\shF(A)\arr\oplus_{i}\shF(A_{i})$.
If we base change along a flat map $B\arr B'$ it is now easy to see
that $\shF(\stY\times_{B}B')\simeq\shF(\stY)\otimes_{B}B'$, as required.

To define the natural transformation $\alpha\colon g^{*}f_{*}\arr f'_{*}g'^{*}$
notice that there is a natural map $f_{*}\shF\arr f_{p}\shF$ which
extends the identity on $\stX_{\text{fl}}$. Applying $g^{*}$ we
get $g^{*}f_{*}\shF\arr g^{*}f_{p}\shF\simeq f'_{p}g'^{*}\shF$ and
then, restricting to $\stX'_{\text{fl}}$, a map $(g^{*}f_{*}\shF)_{|\stX'_{\text{fl}}}\arr(f'_{*}g'^{*}\shF)_{|\stX'_{\text{fl}}}$.
Since both sides are in $\QCoh\stX'_{\text{fl}}$ this map uniquely
extends to a natural transformation $\alpha$ as required. Finally
assume that $g$ is flat and let $\xi\colon\Spec B\arr\stX'\in\stX'_{\text{fl}}$.
If the composition $\Spec B\arr\stX$ is in $\stX_{\text{fl}}$ then
one can easily check that $\alpha(\xi)$ is an isomorphism. Otherwise,
by definition of flatness, there exists an fpqc covering $\Spec B'\arr\Spec B$
whose composition $\xi'\colon\Spec B'\arr\stX'$ satisfies the previous
condition. Since $\alpha(\xi)\otimes_{B}B'\simeq\alpha(\xi')$ we
get the desired result.\end{proof}
\begin{rem}
There are set-theoretic problems in considering global sections of
presheaves and therefore push-forwards, because $\Mod\odi{\stX}$
is in general not locally small. The common way to solve this problem
is to use Grothendieck universes. Take a universe $U$ and define
rings inside $U$, so that $\Aff/R$ is small (with respect to a bigger
universe). Fibered categories should then be required to be small
too. In this situation it is easy to show that $\Mod\odi{\stX}$ is
locally small and therefore global sections and push-forwards are
well defined. With this approach we have to be careful in considering
(big) rings defined starting from some $\shF\in\Mod\odi{\stX}$: for
instance $\Spec\odi{\stX}(\stX)$ is in general not an object of $\Aff/R$.

Notice that global sections and pushforwards of quasi-coherent sheaves
are always well defined for a pseudo-algebraic fibered category and
a pseudo-algebraic map respectively. The reason is that if $\shF\in\QCoh\stX$
and $p\colon X\arr\stX$ is a fpqc atlas then $\shF(\stX)\arr(p^{*}\shF)(X)$
is injective and thus $\shF(\stX)$ is a set.

In the rest of the paper we will not be concerned about those set-theoretic
problems.\end{rem}
\begin{defn}
\label{rem: A module structures of odix modules} If $A$ is an $R$-algebra
and $\shF\in\Mod\odi{\stX}$ then a compatible $A$-module structure
on $\shF$ is the data of $A$-module structures on $\shF(\xi)$ commuting
with the $\Hl^{0}(\odi{\pi(\xi)})$-module structure on $\shF(\xi)$
for all $\xi\in\stX$ and such that, for all $\xi\arr\eta$ in $\stX$,
the map $\shF(\eta)\arr\shF(\xi)$ is $A$-linear. We define $\QCoh_{A}\stX$
as the category of quasi-coherent sheaves over $\stX$ with an $A$-module
structure. We also define $\stX_{A}$ as the fiber product $\Spec A\times_{R}\stX$.
\end{defn}
Notice that if $\stY\arrdi g\stX$ is a map of fibered categories
and $\shF$ is a presheaf of $\odi{\stX}$-modules with an $A$-module
structure then $g^{*}\shF$ inherits an $A$-module structure. In
particular $g^{*}\colon\QCoh\stX\arr\QCoh\stY$ extends to a functor
$\QCoh_{A}\stX\arr\QCoh_{A}\stY$.
\begin{prop}
Let $A$ be an $R$-algebra. Then the push-forward map $\QCoh\stX_{A}\arr\QCoh\stX$
extends naturally to an equivalence $\QCoh\stX_{A}\arr\QCoh_{A}\stX$.\end{prop}
\begin{proof}
The result is very simple if $\stX$ is an affine scheme. In general,
if we set $g\colon\stX_{A}\arr\stX$ for the projection and we consider
$\shG\in\QCoh\stX_{A}$, then $g_{*}\shG\in\QCoh\stX$ and it inherits
an $A$-module structure from the action of $A$ on $\shG$. Therefore
$g_{*}\shG\in\QCoh_{A}\stX$. If $h\colon\Spec B\arr\stX$ is a map
consider the diagrams   \[   \begin{tikzpicture}[xscale=3.4,yscale=-1.0]     
\node (A0_0) at (0.3, 0) {$\Spec (B\otimes_R A)$};     
\node (A0_1) at (1, 0) {$\stX_A$};     \node (A0_2) at (2, 0) {$\QCoh \stX_A$};     \node (A0_3) at (3, 0) {$\QCoh \Spec(B\otimes_R A)$};     \node (A1_0) at (0.3, 1) {$\Spec B$};    
\node (A1_1) at (1, 1) {$\stX$};     \node (A1_2) at (2, 1) {$\QCoh_A \stX$};     \node (A1_3) at (3, 1) {$\QCoh_A \Spec B$};     \path (A0_0) edge [->]node [auto] {$\scriptstyle{h'}$} (A0_1);     \path (A0_1) edge [->]node [auto] {$\scriptstyle{g}$} (A1_1);     \path (A1_0) edge [->]node [auto] {$\scriptstyle{h}$} (A1_1);     \path (A0_3) edge [->]node [auto] {$\scriptstyle{g'_*}$} (A1_3);     \path (A0_2) edge [->]node [auto] {$\scriptstyle{g_*}$} (A1_2);     \path (A0_0) edge [->]node [auto] {$\scriptstyle{g'}$} (A1_0);     \path (A0_2) edge [->]node [auto] {$\scriptstyle{h'^*}$} (A0_3);     \path (A1_2) edge [->]node [auto] {$\scriptstyle{h^*}$} (A1_3);   \end{tikzpicture}   \]  The second diagram is $2$-commutative and the last vertical map
is an equivalence. Using those diagrams it is easy to define a quasi-inverse
$\QCoh_{A}\stX\arr\QCoh\stX_{A}$ of $g_{*}$. %
\begin{comment}
\begin{proof}
We want to construct a quasi-inverse. Let $B$ an $A$-algebra, $M$
a $B$-module with an additional structure of $A$ module compatible
with the $B$-module structure, that is a ring homomorphism $A\arr\End_{B}(M)$.
We will denote this last structure by $a\star m$ for $a\in A$ and
$m\in M$. Then the abelian subgroup of $M$ generated by $a\star m-am$
is automatically a $B$-submodule of $M$ and we define 
\[
\shR(M)=M/\langle a\star m-am\rangle
\]
The $B$-module $\shR(M)$ has a unique $A$-module structure and
the projection $M\arr\shR(M)$ is $B$-linear and $A$-linear with
respect to the action $\star$. If $N$ is a $B$-module then $\Hom_{B}(\shR(M),N)\subseteq\Hom_{B}(M,N)$
is the set of maps $\phi\colon M\arr N$ such that $\phi(a\star m)=a\phi(m)$
for all $a\in A$, $m\in M$.

If $\shF\in\QCoh_{A}\stX$ and $U\arr\stX_{A}$ is a map from a scheme,
then $\Hl^{0}(\odi U)$ is an $A$-algebra and $\shF(U\arr\stX_{A}\arr\stX)$
is an $\Hl^{0}(\odi U)$-module with a compatible additional structure
of $A$-module. We define $\widetilde{\shR}(\shF)(U\arr\stX_{A})=\shR(\shF(U\arr\stX_{A}\arr\stX))$.
If $B,M$ are as above and $B\arr B'$ is a morphism then there is
a $B'$-linear isomorphism $\shR(M)\otimes_{B}B'\arr\shR(M\otimes_{B}B')$.
This tells us that $\widetilde{\shR}(\shF)$ is a functor of $\odi{\stX_{A}}$-module
$\widetilde{\shR}(\shF)\colon\stX_{A}^{\text{op}}\arr\Ab$ and that
its Zariski sheafification $\shR(\shF)\colon\stX_{A}^{\text{op}}\arr\Ab$
coincides with $\widetilde{\shR}(\shF)$ on the objects of $\stX_{A}$
over affine rings and that $\shR(\shF)\in\QCoh\stX_{A}$. Moreover
it is easy to check that $\shR\colon\QCoh_{A}\stX\arr\QCoh\stX_{A}$
is a functor. An isomorphism $g_{*}\circ\shR\simeq\id$ can be obtained
from the following fact: if $C$ is an $R$-algebra, $M'$ a $C$-module
with an $A$-module structure $A\arr\End_{C}(M')$, then $C\otimes_{R}A$
is an $A$-algebra, $M'\otimes_{C}(C\otimes_{R}A)=M'\otimes_{R}A$
is a $C\otimes_{R}A$-module with a compatible additional structure
and the map
\[
M'\arr M'\otimes_{R}A\arr\shR(M'\otimes_{R}A)
\]
is an isomorphism which is $C$-linear and $A$-linear. An isomorphism
$\shR\circ g_{*}\simeq\id$ can be obtained from the following fact:
if $D$ is an $A$-algebra and $M'$ is a $D$-module, then $D\otimes_{R}A$
is an $A$-algebra with respect to the second factor, $ $$M'\otimes_{D}(D\otimes_{R}A)=M'\otimes_{R}A$
is a $D\otimes_{R}A$-module with a compatible additional $A$-module
structure and the multiplication $M'\otimes_{R}A\arr M'$ factors
through a $D$-linear isomorphism $\shR(M'\otimes_{R}A)\arr M'$.
All the above claims became really easy using the description of $\Hom_{B}(\shR(M),N)$
given above.\end{proof}
\end{comment}

\end{proof}
We will almost always regard quasi-coherent sheaves over $\stX_{A}$
as objects of $\QCoh_{A}\stX$.
\begin{rem}
\label{rem: remark on the tensor product over A} If $\shF,\shG\in\QCoh_{A}\stX$
then $\shF\otimes_{\odi{\stX}}\shG$ does not correspond to the tensor
product in $\QCoh\stX_{A}$: $\shF\otimes_{\odi{\stX}}\shG$ has two
distinct $A$-module structures. Under the equivalence $\QCoh\stX_{A}\arr\QCoh_{A}\stX$
the tensor product of $\shF$ and $\shG$, that we will denote by
$\shF\otimes_{\odi{\stX_{A}}}\shG$, is given by 
\[
U\longmapsto\shF(U)\otimes_{\Hl^{0}(\odi U)}\shG(U)/\langle ax\otimes y-x\otimes ay\st x\in\shF(U),y\in\shG(U)\rangle
\]
\end{rem}
\begin{defn}
A locally free sheaf $\E$ (of rank $n$) over $\stX$ is a quasi-coherent
sheaf such that $\E(\Spec B\arr\stX)$ is a finitely generated projective
$B$-module (of rank $n$) for all maps $\Spec B\arr\stX$. We denote
by $\Loc\stX$ the subcategory of $\QCoh\stX$ of locally free sheaves.

We will say that a fiber category $\stX$ has the \emph{resolution
property }if $\Loc\stX$ generates $\QCoh\stX$.
\end{defn}

\section{Sheafification functors.}

In this section we define and describe particular functors that generalize
sheafification functors for affine schemes or projective schemes.
The idea is to interpret the category of modules or graded modules
respectively as a category of $R$-linear functors. More precisely:
\begin{defn}
Given a fibered category $\stX$ over a ring $R$, an $R$-algebra
$A$ and a subcategory $\shD$ of $\QCoh\stX$ we define $\L_{R}(\shD,A)$
as the category of contravariant $R$-linear functors $\Gamma\colon\shD\arr\Mod A$
and natural transformations as arrows. We define a functor $\Omega^{*}\colon\QCoh_{A}\stX\arr\L_{R}(\shD,A)$
by 
\[
\Omega_{-}^{\shF}=\Hom_{\stX}(-,\shF)\colon\shD\arr\Mod A
\]
The functor $\Omega^{*}$ is called the \emph{Yoneda functor }associated
with $\shD$. A left adjoint of $\Omega^{*}$ is called a \emph{sheafification
}functor associated with $\shD$. If $\shF\in\QCoh_{A}\stX$ we will
call $\Omega^{\shF}$ the Yoneda functor associated with $\shF$.\end{defn}
\begin{example}
\label{example: quasi-projective case} The analogy with the sheafification
functor associated with $\stX=\PP_{R}^{n}$ or any quasi-projective
and quasi-compact scheme over $R$ is the following. If $\shC=\{\odi{\stX}(n)\}_{n\in\Z}$
then with $\Gamma\in\L_{R}(\shC,R)$ we can associate the $\Z$-graded
$R$-module $M=\bigoplus_{n\in\Z}\Gamma_{\odi{\stX}(-n)}$. The functorial
properties of $\Gamma$ allow us to define a structure of graded $S$-module
on $M$, where $S=\bigoplus_{n\in\Z}\Hl^{0}(\odi{\stX}(n))$ is the
homogeneous coordinate ring of $\stX$. This associations extends
to an equivalence of categories between $\L_{R}(\shC,R)$ and the
category of graded $S$-modules. The functor $\Omega^{*}$ corresponds
to the functor $\Gamma_{*}$ which carries a quasi-coherent sheaf
$\shF$ on $\stX$ to the graded $S$-module $\bigoplus_{n\in\Z}\Hl^{0}(\shF(n))$.
Finally the sheafification functor from the graded $S$-modules to
$\QCoh\stX$ is left adjoint to $\Gamma_{*}$.
\end{example}
Let us fix an $R$-algebra $A$ and a fibered category $\pi\colon\stX\arr\Aff/R$.


\subsection{Sheafifying $R$-linear functors.}

In this section we want to explicitly describe sheafification functors
for small subcategories of $\QCoh\stX$. In particular we fix a \emph{small
}(and non empty) subcategory $\shC$ of $\QCoh\stX$.

In the construction of the sheafification functors we will make use
of the coend construction in the settings of categories enriched over
categories of modules over a ring. The general theory simplifies considerably
in this context and we will also apply such construction only in particular
cases. In the following remark we collect all the properties we will
need.
\begin{rem}
\label{rem: coends} Let $\stY$ be a fibered category over $R$,
$F\colon\shC\arr\QCoh\stY$ be an $R$-linear functor and $\Gamma\in\L_{R}(\shC,A)$.
The coend of the $R$-linear functor $\Gamma_{-}\otimes_{R}F_{-}\colon\shC^{\op}\otimes\shC\arr\QCoh_{A}\stY$,
denoted by
\[
\int^{\E\in\shC}\Gamma_{\E}\otimes_{R}F_{\E}\in\QCoh_{A}\stY
\]
is the cokernel of the map
\[
\bigoplus_{\E\arrdi u\overline{\E}}(\Gamma_{u}\otimes\id_{F_{\E}}-\id_{\Gamma_{\overline{\E}}}\otimes F_{u})\colon\bigoplus_{\E\arrdi u\overline{\E}}\Gamma_{\overline{\E}}\otimes_{R}F_{\E}\arr\bigoplus_{\E\in\shC}\Gamma_{\E}\otimes_{R}F_{\E}
\]
Moreover it comes equipped with an $A$-linear natural isomorphism
\[
\Hom_{\QCoh_{A}\stY}(\int^{\E\in\shC}\Gamma_{\E}\otimes_{R}F_{\E},\shH)\arrdi{\alpha}\Hom_{\L_{R}(\shC,A)}(\Gamma,\Hom_{\stY}(F_{-},\shH))\text{ for }\shH\in\QCoh_{A}\stY
\]
given by
\[
\alpha(\int^{\E\in\shC}\Gamma_{\E}\otimes_{R}F_{\E}\arrdi u\shH)(x)=u\circ p_{\overline{\E}}(x\otimes-)\colon F_{\overline{\E}}\arr\shH\text{ for }\overline{\E}\in\shC,x\in\Gamma_{\overline{\E}}
\]
where $p_{\E}\colon\Gamma_{\E}\otimes_{R}F_{\E}\arr\int^{\E\in\shC}\Gamma_{\E}\otimes_{R}F_{\E}$
for $\E\in\shC$ are the structure morphisms. Its inverse is uniquely
determined by the expression 
\[
\alpha^{-1}(\Gamma\arrdi v\Hom_{\stY}(F_{-},\shH))\circ p_{\E}\colon\Gamma_{\E}\otimes_{R}\E\arr\shH\comma x\otimes y\longmapsto v_{\E}(x)(y)\text{ for }\E\in\shC
\]
Natural transformations $F\arr F'$ and $\Gamma\arr\Gamma'$ yields
morphisms $\int^{\E\in\shC}\Gamma_{\E}\otimes_{R}F_{\E}\arr\int^{\E\in\shC}\Gamma_{\E}\otimes_{R}F'_{\E}$
and $\int^{\E\in\shC}\Gamma_{\E}\otimes_{R}F_{\E}\arr\int^{\E\in\shC}\Gamma'_{\E}\otimes_{R}F_{\E}$
respectively. Those can be defined either using Yoneda's lemma and
the above characterization of $\Hom(\int^{\E\in\shC}\Gamma_{\E}\otimes_{R}F_{\E},-)$
or directly using the description of $\int^{\E\in\shC}\Gamma_{\E}\otimes_{R}F_{\E}$
as a cokernel.

All the above claims are standard in the theory of coend in the enriched
settings (in our case enriched by $\Mod R$), but, in this simplified
context, it is elementary to prove them directly.
\end{rem}
We start by showing that $\shC$ (and therefore any essentially small
subcategory of $\QCoh\stX$) admits a sheafification functor.
\begin{prop}
\label{prop:adjoint of Omegastar with the cokernel} The Yoneda functor
$\Omega^{*}\colon\QCoh_{A}\stX\arr\L_{R}(\shC,A)$ has a left adjoint
$\shF_{-,\shC}\colon\L_{R}(\shC,A)\arr\QCoh_{A}\stX$ given by
\[
\shF_{\Gamma,\shC}\colon\stX^{\text{op}}\arr\sets\comma\stX\ni\xi\longmapsto\int^{\E\in\shC}\Gamma_{\E}\otimes_{R}\E(\xi)\in\Mod(\Hl^{0}(\odi{\pi(\xi)})\otimes_{R}A)
\]
where $\E(\xi)$ denotes the evaluation $\shC\arr\Mod\Hl^{0}(\odi{\pi(\xi)})$
of sheaves in $\xi\in\stX$. Alternatively 
\[
\shF_{\Gamma,\shC}=\int^{\E\in\shC}\Gamma_{\E}\otimes_{R}\E\in\QCoh_{A}\stX
\]
where $\E$ denotes the inclusion $\shC\arr\QCoh\stX$.\end{prop}
\begin{proof}
It is enough to apply \ref{rem: coends} with $\stY=\stX$ and $F\colon\shC\arr\QCoh\stX$
the inclusion. Using the description of coend as cokernel one can
check that the two functors defined in the statement are canonically
isomorphic.\end{proof}
\begin{defn}
\label{def: unit and counit for adjunction} We denote by $\gamma_{\Gamma}\colon\Gamma_{-}\arr\Omega_{-}^{\shF_{\Gamma,\shC}}=\Hom_{\stX}(-,\shF_{\Gamma,\shC})$
and $\delta_{\shG}\colon\shF_{\Omega^{\shG},\shC}\arr\shG$ for $\Gamma\in\L_{R}(\shC,A)$
and $\shG\in\QCoh_{A}\stX$ the unit and the counit of the adjunction
between $\Omega^{*}\colon\QCoh_{A}\stX\arr\L_{R}(\shC,A)$ and $\shF_{*,\shC}\colon\L_{R}(\shC,A)\arr\QCoh_{A}\stX$
respectively.

Given $\xi\in\stX$, $\E\in\shC$, $\psi\in\E(\xi)$ and $x\in\Gamma_{\E}$
we denote by $x_{\E,\psi}\in\shF_{\Gamma,\shC}(\xi)$ the image of
$x\otimes\psi$ under the map $\Gamma_{\E}\otimes_{R}\E(\xi)\arr\shF_{\Gamma,\shC}(\xi)$\end{defn}
\begin{prop}
\label{prop:FGammaC on affine schemes} Let $\Gamma\in\L_{R}(\shC,A)$.
The unit $\gamma_{\Gamma}\colon\Gamma_{-}\arr\Omega_{-}^{\shF_{\Gamma,\shC}}=\Hom_{\stX}(-,\shF_{\Gamma,\shC})$
is given by   \[   \begin{tikzpicture}[xscale=3.8,yscale=-0.6]     
\node (A0_0) at (0.2, 0) {$\Gamma_\E$};   
\node (A1_0) at (0.2, 1) {$x$}; 
\node (A0_1) at (1, 0) {$\Hom_{\stX}(\E,\shF_{\Gamma,\shC})$};     
\node (A1_1) at (1, 1) {$(\phi \longmapsto x_{\E,\phi})$};     
\path (A0_0) edge [->]node [auto] {$\scriptstyle{}$} (A0_1);     
\path (A1_0) edge [|->,gray]node [auto] {$\scriptstyle{}$} (A1_1);        
\end{tikzpicture}   \]  If $\shG\in\QCoh_{A}\stX$ the counit $\delta_{\shG}\colon\shF_{\Omega^{\shG},\shC}\arr\shG$
is given by 
\[
\shF_{\Omega^{\shG},\shC}(\xi)\ni x_{\E,\psi}\longmapsto x(\psi)\in\shG(\xi)\text{ for }\E\in\shC,x\in\Omega_{\E}^{\shG}=\Hom_{\stX}(\E,\shG),\xi\in\stX,\psi\in\E(\xi)
\]
\end{prop}
\begin{proof}
If $p_{\E}\colon\Gamma_{\E}\otimes_{R}\E\arr\shF_{\Gamma,\shC}$ for
$\E\in\shC$ are the maps associated with the coend defining $\shF_{\Gamma,\shC}$,
then $p_{\E}(x\otimes\psi)=x_{\E,\psi}$ for $x\in\Gamma_{\E},\psi\in\E(\xi)$
and $\xi\in\stX$. All the claims follows by a direct check using
the explicit description of the isomorphism
\[
\Hom_{\QCoh_{A}\stX}(\int^{\E\in\shC}\Gamma_{\E}\otimes_{R}\E,\shH)\arrdi{\alpha}\Hom_{\L_{R}(\shC,A)}(\Gamma,\Omega^{\shH})\text{ for }\shH\in\QCoh_{A}\stX
\]
and its inverse given in \ref{rem: coends}.
\end{proof}
Given a map $g\colon\stY\arr\stX$ of fibered categories we want to
express $g^{*}\shF_{\Gamma,\shC}\in\QCoh_{A}\stY$ for $\Gamma\in\L_{R}(\shC,A)$
as $\shF_{g^{*}\Gamma,g^{*}\shC}$ for a suitable choice of $g^{*}\shC\subseteq\QCoh\stY$
and $g^{*}\Gamma\in\L_{R}(g^{*}\shC,A)$.
\begin{defn}
Let $\stY$ be a fibered category, $g\colon\stY\arr\stX$ be a morphism
and $\shD$ be a subcategory of $\QCoh\stX$. We set $g^{*}\shD$
for the subcategory of $\QCoh\stY$ of sheaves $g^{*}\E$ for $\E\in\shD$.
If $\shD'\subseteq\QCoh\stY$ is a subcategory containing $g^{*}\shD$
we can define a restriction functor   \[   \begin{tikzpicture}[xscale=2.7,yscale=-0.6]     \node (A0_0) at (0, 0) {$\L_R(\shD',A)$};     \node (A0_1) at (1, 0) {$\L_R(\shD,A)$};     \node (A1_0) at (0, 1) {$\Gamma$};     \node (A1_1) at (1, 1) {$\Gamma\circ g^*$};     \path (A0_0) edge [->]node [auto] {$\scriptstyle{g_*}$} (A0_1);     \path (A1_0) edge [|->,gray]node [auto] {$\scriptstyle{}$} (A1_1);   \end{tikzpicture}   \] \end{defn}
\begin{prop}
\label{thm:adjoint of push forward} Let $\stY$ be a fibered category,
$g\colon\stY\arr\stX$ be a morphism and $\shD$ be a subcategory
of $\QCoh\stY$ such that $g^{*}\shC\subseteq\shD$. Then $g_{*}\colon\L_{R}(\shD,A)\arr\L_{R}(\shC,A)$
has a left adjoint $g^{*}\colon\L_{R}(\shC,A)\arr\L_{R}(\shD,A)$
and it is given by
\[
(g^{*}\Gamma)_{\shG}=\int^{\E\in\shC}\Gamma_{\E}\otimes_{R}\Hom_{\stY}(\shG,g^{*}\E)\in\Mod A\text{ for }\Gamma\in\L_{R}(\shC,A)\comma\shG\in\shD
\]
where $\Hom_{\stY}(\shG,g^{*}-)$ is thought of as a functor $\shC\arr\Mod R$.
If $\stY=\stX$ and $g=\id_{\stX}$, so that $\shC\subseteq\shD$
and $(\id_{\stX})_{*}\colon\L_{R}(\shD,A)\arr\L_{R}(\shC,A)$ is the
restriction, then the unit $\Gamma\arr(\id_{\stX}^{*}\Gamma)_{|\shC}$
is an isomorphism for $\Gamma\in\L_{R}(\shC,A)$.\end{prop}
\begin{proof}
Let $\Omega\in\L_{R}(\shD,A)$. $ $Applying \ref{rem: coends} with
$F=\Hom_{\stY}(\shG,g^{*}-)$ we get a bijection between $\Hom_{A}((g^{*}\Gamma)_{\shG},\Omega_{\shG})$
and the set of $A$-linear natural transformations $\Gamma_{\E}\arr\Hom_{R}(\Hom(\shG,g^{*}\E),\Omega_{\shG})$
for $\E\in\shC$. A natural transformation $g^{*}\Gamma\arr\Omega$
corresponds to a collection $\underline{\gamma}$ of $A$-linear maps
$\gamma_{\shG,\E}\colon\Gamma_{\E}\arr\Hom_{R}(\Hom(\shG,g^{*}\E),\Omega_{\shG})$
natural in $\E\in\shC$ and such that
\[
\gamma_{\shG,\E}(x)(\phi\circ u)=\Omega_{u}(\gamma_{\overline{\shG},\E}(x)(\phi))\text{ for }x\in\Gamma_{\E},\overline{\shG}\arrdi{\phi}g^{*}\E,\shG\arrdi u\overline{\shG}\in\shD
\]
Set $\mu_{\underline{\gamma}}(x)=\gamma_{g^{*}\E,\E}(x)(\id_{g^{*}\E})\in\Omega_{g^{*}\E}$
for $\E\in\shC,x\in\Gamma_{\E}$. If $\underline{\gamma}$ is a collection
as above it follows that $\gamma_{\shG,\E}(x)(\psi)=\Omega_{\psi}(\mu_{\underline{\gamma}}(x))$
for $\E\in\shC$, $\shG\in\shD$, $x\in\Gamma_{\E}$ and $\shG\arrdi{\psi}g^{*}\E$
and that $\mu_{\underline{\gamma}}\colon\Gamma_{\E}\arr\Omega_{g^{*}\E}=(g_{*}\Omega)_{\E}$
is $A$-linear and natural in $\E$. Conversely, given a morphism
$\mu\colon\Gamma\arr(g_{*}\Omega)$ in $\L_{R}(\shC,A)$ we can always
define a collection $\underline{\gamma}$ as above by setting $\gamma_{\shG,\E}(x)(\psi)=\Omega_{\psi}(\mu(x))$.
It is easy to check that $\underline{\gamma}$ induces a morphism
$g^{*}\Gamma\arr\Omega$ and that the above constructions yields a
bijection $\Hom(g^{*}\Gamma,\Omega)\simeq\Hom(\Gamma,g_{*}\Omega)$.

Assume now $\stY=\stX$ and $g=\id_{\stX}$ and let $\Gamma\in\L_{R}(\shC,A)$
and $\overline{\E}\in\shC$. Denote by $\alpha\colon\Gamma\arr(\id_{\stX}^{*}\Gamma)_{|\shC}$
the unit morphism. If $p_{\widetilde{\E}}\colon\Gamma_{\widetilde{\E}}\otimes\Hom_{\stX}(\overline{\E},\widetilde{\E})\arr(\id_{\stX}^{*}\Gamma)_{\overline{\E}}$
are the structure morphisms as in \ref{rem: coends}, then $\alpha_{\overline{\E}}=p_{\overline{\E}}(-\otimes\id_{\overline{\E}})$.
In particular, given $\shH\in\Mod A$ and using \ref{rem: coends},
the map $\Hom_{A}(\alpha_{\overline{\E}},\shH)\colon\Hom_{A}((\id_{\stX}^{*}\Gamma)_{\overline{\E}},\shH)\arr\Hom_{A}(\Gamma_{\overline{\E}},\shH)$
sends an $A$-linear natural transformation $\delta\colon\Gamma_{-}\arr\Hom_{R}(\Hom_{\stX}(\overline{\E},-),\shH)$
to $\Gamma_{\overline{\E}}\ni x\longmapsto\delta(x)(\id_{\overline{\E}})\in\shH$.
Since $\delta$ corresponds to an $R$-linear natural transformation
$\Hom_{\stX}(\overline{\E},-)\arr\Hom_{A}(\Gamma_{-},\shH)$, by Yoneda
lemma we see that $\Hom_{A}(\alpha_{\overline{\E}},\shH)$ is an isomorphism.
\end{proof}
The above proposition yields a natural extension of any $\Gamma\in\L_{R}(\shC,A)$
to a functor $\Gamma^{ex}\in\L_{R}(\QCoh\stX,A)$. By abuse of notation
we will denote them by the same symbol $\Gamma$. This means that
if $\Gamma\in L_{R}(\shC,A)$ and $\shG\in\QCoh\stX$ then we can
evaluate $\Gamma$ on $\shG$, writing $\Gamma_{\shG}$.

Given a map $g\colon\stY\arr\stX$ we will denote by $g^{*}\colon\L_{R}(\shC,A)\arr\L_{R}(g^{*}\shC,A)$
the left adjoint of the restriction $\L_{R}(g^{*}\shC,A)\arr\L_{R}(\shC,A)$.
So, given $\Gamma\in\L_{R}(\shC,A)$, $g^{*}\Gamma$ is a functor
$g^{*}\shC\arr\Mod A$ but it also defines a functor $\QCoh\stY\arr\Mod A$
denoted, by our convention, by the same symbol. By \ref{thm:adjoint of push forward}
the functor $g^{*}\Gamma\colon\QCoh\stY\arr\Mod A$ coincides with
the value of the left adjoint of the restriction $\L_{R}(\QCoh\stY,A)\arr\L_{R}(\shC,A)$.
\begin{rem}
\label{rem: extends to h0OX}Given $\Gamma\in\L_{R}(\shC,A)$ and
$\E\in\shC$ we have $R$-linear morphisms of rings 
\[
\Hl^{0}(\odi{\stX})\simeq\End_{\stX}(\odi{\stX})\arr\End_{\stX}(\E)\arr\End_{A}(\Gamma_{\E})
\]
This defines a lifting of $\Gamma$ to an $R$-linear functor $\Gamma\colon\shC\arr\Mod(\Hl^{0}(\odi{\stX})\otimes_{R}A)$
and an equivalence
\[
\L_{R}(\shC,A)\arr\L_{R}(\shC,\Hl^{0}(\odi{\stX})\otimes_{R}A)
\]
In particular, if $g\colon\Spec B\arr\stX$ is a map and $\Gamma\in\L_{R}(\shC,A)$
then $(g^{*}\Gamma)_{B}$ has a $B\otimes_{R}A$-module structure.
By \ref{prop:adjoint of Omegastar with the cokernel} and \ref{thm:adjoint of push forward}
there is a canonical $A$-linear isomorphism 
\[
\shF_{\Gamma,\shC}(B)\simeq(g^{*}\Gamma)_{B}
\]
and it is easy to see that it is also $B$-linear.

\begin{comment}
\begin{defn}
\label{def: valuation in odi X} We define a functor $(-)_{\odi{\stX}}\colon\L_{R}(\shC,A)\arr\Mod(\Hl^{0}(\odi{\stX})\otimes_{R}A)$
as the composition
\[
\L_{R}(\shC,A)\arr\L_{R}(\QCoh\stX,A)\arrdi{(-)_{\odi{\stX}}}\Mod(\Hl^{0}(\odi{\stX})\otimes_{R}A)
\]
where the first functor is the left adjoint of the restriction functor
$\L_{R}(\QCoh\stX,A)\arr\L_{R}(\shC,A)$. We will call (improperly)
it the evaluation in $\odi{\stX}$ and, given $\Gamma\in\L_{R}(\shC,A)$,
we will simply write $\Gamma_{\odi{\stX}}$, although $\odi{\stX}$
is not necessary an object of $\shC$. \end{defn}
\begin{rem}
By \ref{prop:find a directed category} the module $\Gamma_{\odi{\stX}}$
does not change if we think of $\Gamma\in\L_{R}(\shC,A)$ or its extension
$\Gamma\in\L_{R}(\shC^{\oplus},A)$. In particular, by \ref{thm:adjoint of push forward}
we have
\[
\Gamma_{\odi{\stX}}=\varinjlim_{\odi{\stX}\arr\E\in\shC^{\oplus}}\Gamma_{\E}
\]
When $\odi{\stX}\in\shC$, the category $J_{\odi{\stX},\shC}$ over
which the limit is taken has an initial object, namely $\odi{\stX}\arrdi{\id}\odi{\stX}$,
and therefore the two $A$-modules represented by the same symbol
$\Gamma_{\odi{\stX}}$ are canonically isomorphic. \end{rem}
\end{comment}
\end{rem}
\begin{prop}
Let $g\colon\stY\arr\stX$ be a morphism of fibered categories. Then
there exists an isomorphism $g^{*}\shF_{\Gamma,\shC}\simeq\shF_{g^{*}\Gamma,g^{*}\shC}$
natural in $\Gamma\in\L_{R}(\shC,A)$, that is a $2$-commutative
diagram   \[   \begin{tikzpicture}[xscale=3.1,yscale=-1.2]     \node (A0_0) at (0, 0) {$\L_R(\shC,A)$};     \node (A0_1) at (1, 0) {$\QCoh_A \stX$};     \node (A1_0) at (0, 1) {$\L_R(g^*\shC,A)$};     \node (A1_1) at (1, 1) {$\QCoh_A \stY$};     \path (A0_0) edge [->]node [auto] {$\scriptstyle{\shF_{-,\shC}}$} (A0_1);     \path (A0_0) edge [->]node [auto] {$\scriptstyle{g^*}$} (A1_0);     \path (A0_1) edge [->]node [auto] {$\scriptstyle{g^*}$} (A1_1);     \path (A1_0) edge [->]node [auto] {$\scriptstyle{\shF_{-,g^*\shC}}$} (A1_1);   \end{tikzpicture}   \] \end{prop}
\begin{proof}
Let $\Gamma\in\L_{R}(\shC,A)$. Let also $\Spec B\arrdi{\xi}\stY$
be a map and $N\in\Mod B\otimes_{R}A$. Denote by $F\colon\shC\arr\Mod B$
and $G\colon g^{*}\shC\arr\Mod B$ the functors obtained by taking
sections over $B$. In particular $F=G\circ g^{*}$. We have isomorphisms
\begin{alignat*}{1}
\Hom_{B\otimes_{R}A}(\shF_{g^{*}\Gamma,g^{*}\shC}(B),N) & \simeq\Hom_{\L_{R}(g^{*}\shC,A)}(g^{*}\Gamma,\Hom_{B}(G,N))\simeq\Hom_{\L_{R}(\shC,A)}(\Gamma,g_{*}\Hom_{B}(G,N))\\
 & =\Hom_{\L_{R}(\shC,A)}(\Gamma,\Hom_{B}(F,N))\simeq\Hom_{B\otimes_{R}A}((g^{*}\shF_{\Gamma,\shC})(B),N)
\end{alignat*}
In particular we get an isomorphism $(g^{*}\shF_{\Gamma,\shC})(B)\simeq\shF_{g^{*}\Gamma,g^{*}\shC}(B)$.
By a direct check it follows that this isomorphism is natural in $\Gamma$
and it also extends as an isomorphism of functors $g^{*}\shF_{\Gamma,\shC}\simeq\shF_{g^{*}\Gamma,g^{*}\shC}$.
\end{proof}
Notice that, as soon as $g^{*}\colon\QCoh\stX\arr\QCoh\stY$ has a
right adjoint $g_{*}\colon\QCoh\stY\arr\QCoh\stX$ (see \ref{prop:properties of push-forward}),
then $g_{*}\Omega^{\shG}\simeq\Omega^{g_{*}\shG}$ for $\shG\in\QCoh_{A}\stY$
and the above proposition follows taking the left adjoints.
\begin{rem}
\label{rem: LR and QCoh as stacks} If $A\arr A'$ is a morphism of
$R$-algebras then we have pull-back functors $\L_{R}(\shC,A)\arr\L_{R}(\shC,A')$
and $\QCoh_{A}\stX\arr\QCoh_{A'}\stX$. The first one is obtained
considering the tensor product $-\otimes_{A}A'$, while the second
one corresponds to the pullback $\QCoh\stX_{A}\arr\QCoh\stX_{A'}$
along the projection $\stX_{A'}\arr\stX_{A}$. Alternatively, those
functors are left adjoints to the restriction of scalars $\L_{R}(\shC,A')\arr\L_{R}(\shC,A)$
and $\QCoh_{A'}\stX\arr\QCoh_{A}\stX$ respectively. It is easy to
see that in this way we obtain two fpqc stacks (not in groupoids)
$\L_{R}(\shC,-)$ and $\QCoh_{-}\stX$ over the category of affine
$R$-schemes. Notice that the functor $\Omega^{*}\colon\QCoh_{-}\stX\arr\L_{R}(\shC,-)$
is not a morphism of stacks because $\Hom_{\stX}(\E,\shG)\otimes_{A}A'\not\simeq\Hom_{\stX}(\E,\shG\otimes_{A}A')$
in general for $\E\in\shC$ and $\shG\in\QCoh_{A}\stX$.\end{rem}
\begin{prop}
\label{prop:base change for FstarC} The functor $\shF_{*,\shC}\colon\L_{R}(\shC,-)\arr\QCoh_{-}\stX$
is a morphism of stacks.\end{prop}
\begin{proof}
Given a morphism $A\arr A'$ of $R$-algebras we have a $2$-commutative
diagram   \[   \begin{tikzpicture}[xscale=2.9,yscale=-1.0]     \node (A0_0) at (0, 0) {$\QCoh_{A'}\stX$};     \node (A0_1) at (1, 0) {$\L_R(\shC,A')$};     \node (A1_0) at (0, 1) {$\QCoh_A \stX$};     \node (A1_1) at (1, 1) {$\L_R(\shC,A)$};     \path (A0_0) edge [->]node [auto] {$\scriptstyle{\Omega^*}$} (A0_1);     \path (A0_0) edge [->]node [auto] {$\scriptstyle{}$} (A1_0);     \path (A0_1) edge [->]node [auto] {$\scriptstyle{}$} (A1_1);     \path (A1_0) edge [->]node [auto] {$\scriptstyle{\Omega^*}$} (A1_1);   \end{tikzpicture}   \] 
where the vertical arrows are obtained by restricting the scalars
from $A'$ to $A$. Using \ref{rem: LR and QCoh as stacks} and taking
the left adjoint functors of the functors in the diagram we exactly
get the $2$-commutative diagram expressing the fact that $\shF_{*,\shC}$
preserves Cartesian arrows.
\end{proof}
We conclude this section by showing that, when considering sheafification
functors $\shF_{-,\shC}$, we can always reduce problems to the case
when $\shC$ is an additive category. Moreover in this case the sections
of $\shF_{-,\shC}$ have a nice expression in terms of a direct limit.
\begin{defn}
Given a subcategory $\shD$ of $\QCoh\stX$ we denote by $\shD^{\oplus}$
the subcategory of $\QCoh\stX$ whose objects are all finite direct
sums of sheaves in $\shD$.
\end{defn}
Notice that if $\shD$ is small then $\shD^{\oplus}$ is small.
\begin{prop}
\label{prop:find a directed category} The restriction $\L_{R}(\shC^{\oplus},A)\arr\L_{R}(\shC,A)$
and its left adjoint are inverses of each other. In particular if
$\Gamma\in\L_{R}(\shC^{\oplus},A)$ then we have a canonical isomorphism
$\shF_{\Gamma,\shC^{\oplus}}\simeq\shF_{\Gamma_{|\shC},\shC}$.\end{prop}
\begin{proof}
Denote the restriction by $\alpha$ and its left adjoint by $\beta$.
From \ref{thm:adjoint of push forward} we know that $\alpha\beta(\Gamma)\simeq\Gamma$
for $\Gamma\in\L_{R}(\shC,A)$. Conversely, if $\Gamma\in\L_{R}(\shC^{\oplus},A)$,
we have a canonical morphism $\gamma\colon\beta\alpha(\Gamma)\arr\Gamma$.
If $\E\in\shC$ the map $\gamma_{\E}$ is easily seen to be an isomorphism
and by additivity this also follows for $\gamma$. For the last claim
it is enough to note that the composition $\QCoh_{A}\stX\arrdi{\Omega^{*}}\L_{R}(\shC^{\oplus},A)\arr\L_{R}(\shC,A)$
is exactly $\Omega^{*}\colon\QCoh_{A}\stX\arr\L_{R}(\shC,A)$.\end{proof}
\begin{rem}
\label{rem: natural transf are linear for additive categories} If
$\catC$ is an $R$-linear and additive category and $F,G\colon\catC\arr\Mod A$
are $R$-linear (covariant or contravariant) functors then any natural
transformation $\lambda\colon F\arr G$ of functors of sets is $R$-linear.
Indeed by considering $\catC^{\op}$ we can consider only covariant
functors. In this case it is easy to show that the maps $\lambda_{X}\colon F(X)\arr G(X)$
for $X\in\catC$ are $R$-linear using functoriality on the map $r\id_{X}\colon X\arr X$
for $r\in R$ and $\pr_{1},\pr_{2},\pr_{1}+\pr_{2}\colon X\oplus X\arr X$,
where $\pr_{*}$ are the projections.\end{rem}
\begin{defn}
Let $\Spec B\arr\stX$ be a map. We denote by $J_{B,\shC}$ the category
of pairs $(\E,\psi)$ where $\E\in\shC^{\oplus}$ and $\psi\in\E(B)$.
Given $\Gamma\in\L_{R}(\shC,A)$ we have a functor $\Gamma\colon J_{B,\shC}\arr\Mod A$
given by $\Gamma_{\E,\psi}=\Gamma_{\E}$.\end{defn}
\begin{prop}
\label{prop:FGammaCB as a direct limit} Let $\Spec B\arr\stX$ be
a map and $\Gamma\in\L_{R}(\shC,A)$. The category $J_{B,\shC}$ is
non-empty and for all $\xi,\xi'\in J_{B,\shC}$ there exists $\xi''\in J_{B,\shC}$
and maps $\xi''\arr\xi$, $\xi''\arr\xi'$. The $A$-linear maps $\Gamma_{\E,\psi}=\Gamma_{\E}\arr\shF_{\Gamma,\shC^{\oplus}}(B)\simeq\shF_{\Gamma,\shC}(B)$,
$x\longmapsto x_{\E,\psi}$ for $(\E,\psi)\in J_{B,\shC}$ (see \ref{def: unit and counit for adjunction})
induce an $A$-linear isomorphism 
\[
\varinjlim_{(\E,\psi)\in J_{B,\shC}^{\op}}\Gamma_{\E,\psi}\arr\shF_{\Gamma,\shC}(B)
\]
In particular all elements of $\shF_{\Gamma,\shC}(B)$ are of the
form $x_{\E,\psi}$ for some $\E\in\shC^{\oplus}$, $x\in\Gamma_{\E}$
and $\psi\in\E(B)$. The multiplication by $b\in B$ on the first
limit is induced by mapping $\Gamma_{\E,\psi}$ to $\Gamma_{\E,b\psi}$
using $\id_{\Gamma_{\E}}$ for $(\E,\psi)\in J_{B,\shC}$.\end{prop}
\begin{proof}
By \ref{prop:find a directed category} we can assume $\shC=\shC^{\oplus}$.
Denote by $\shH$ and $\alpha\colon\shH\arr\shF_{\Gamma,\shC}(B)$
the limit and the map in the statement respectively. The category
$J_{B,\shC}$ is not empty because $(\E,0)\in J_{B,\shC}$ for all
$\E\in\shC$ and the map $\alpha$ is well defined because, for all
$x\in\Gamma_{\overline{\E}}$ and for all $(\E,\psi)\arrdi u(\overline{\E},u(\psi))$
we have $x_{\overline{\E},u(\psi)}=(\Gamma_{u}(x))_{\E,\psi}$, by
definition of $\shF_{\Gamma,\shC}(B)$ as coend. Moreover if $(\E_{1},\psi_{1}),(\E_{2},\psi_{2})\in J_{B,\shC}$
then we have maps $\pr_{i}\colon(\E_{1}\oplus\E_{2},\psi_{1}\oplus\psi_{2})\arr(\E_{i},\psi_{i})$
for $i=1,2$, where $\pr_{i}$ is the projection. In particular any
element of $\shH$ comes from a map $\Gamma_{\E,\psi}\arr\shH$, where
$(\E,\psi)\in J_{B,\shC}$. Thus all elements of $\shF_{\Gamma,\shC}(B)$
are of the form $x_{\E,\psi}$ for some $\E\in\shC^{\oplus}$, $x\in\Gamma_{\E}$
and $\psi\in\E(B)$ provided that we prove that $\alpha$ is an isomorphism.

Given an $A$-module $N$ then $\Hom_{A}(\shH,N)$ is $A$-linearly
isomorphic to the set of natural transformations of sets $\beta_{\E}\colon\E(B)\arr\Hom_{A}(\Gamma_{\E},N)$.
Since $\shC$ is additive, by \ref{rem: natural transf are linear for additive categories},
those transformations are automatically $R$-linear. Given $b\in B$
there is an $A$-linear map $\phi_{b}\in\End_{A}(\shH)$ as described
in the last sentence in the statement. A direct check shows that,
given an $A$-module $N$, the map $\Hom_{A}(\phi_{b},N)\colon\Hom_{A}(\shH,N)\arr\Hom_{A}(\shH,N)$
sends an $R$-linear natural transformation $\beta_{\E}\colon\E(B)\arr\Hom_{A}(\Gamma_{\E},N)$
to $\beta_{\E}\circ b\id_{\E}$. This easily implies that $\phi\colon B\arr\End_{A}(\shH)$
makes $\shH$ into a $B\otimes_{R}A$-module and that $\alpha\colon\shH\arr\shF_{\Gamma,\shC}(B)$
is also $B$-linear. Since $B\otimes_{R}A$-linear maps $\shH\arr N$,
for $N\in\Mod B\otimes_{R}A$, corresponds to $B$-linear natural
transformations $\E(B)\arr\Hom_{A}(\Gamma_{\E},N)$ and thus to $A$-linear
natural transformations $\Gamma_{\E}\arr\Hom_{B}(\E(B),N)$ as described
above, one can check directly (using \ref{rem: coends}) that $\Hom_{B\otimes_{R}A}(\alpha,N)$
induces the identity on $\Hom_{\L_{R}(\shC,A)}(\Gamma_{-},\Hom_{B}(-(B),N))$.
This implies that $\alpha\colon\shH\arr\shF_{\Gamma,\shC}(B)$ is
an isomorphism.%
\begin{comment}
\begin{notation}
Given $\Gamma\colon\shC\arr\Mod A\in\L_{R}(\shC,A)$ we will denote
with the same symbol $\Gamma$ its extension to $\shC^{\oplus}\arr\Mod A$.\end{notation}
\begin{rem}
\label{rem: replacing C by Coplus} If $g\colon U\arr\stX$ is a map
from a scheme, the compositions of the restrictions $\L_{R}(\QCoh U,A)\arr\L_{R}(g^{*}\shC,A)\arr\L_{R}(\shC,A)$
and $\L_{R}(\QCoh U,A)\arr\L_{R}(\shC^{\oplus},A)\arr\L_{R}(\shC,A)$
are both isomorphic to the restriction $\L_{R}(\QCoh U,A)\arr\L_{R}(\shC,A)$.
By \ref{thm:adjoint of push forward} we have natural isomorphisms
of $\Hl^{0}(\odi U)\otimes_{R}A$-modules 
\[
\widetilde{\shF}_{\Gamma,\shC^{\oplus}}(U\arrdi g\stX)\simeq\widetilde{\shF}_{\Gamma,\shC}(U\arrdi g\stX)\text{ for }\Gamma\in\L_{R}(\shC,A)
\]
We can conclude that the functors $\L_{R}(\shC^{\oplus},A)\arr\L_{R}(\shC,A)\arrdi{\shF_{-,\shC}}\QCoh_{A}\stX$
and $\shF_{-,\shC^{\oplus}}\colon\L_{R}(\shC^{\oplus},A)\arr\QCoh_{A}\stX$
are isomorphic. By \ref{thm:adjoint of push forward} we can also
conclude that
\[
\shF_{\Gamma,\shC}(B)=\varinjlim_{(\E,x)\in J_{B,\shC^{\oplus}}}\Gamma_{\E,x}\text{ for }\Gamma\in\L_{R}(\shC,A)\comma\Spec B\arr\stX
\]

\end{rem}

\begin{rem}
\label{rem: elements of the limit when C is additive} If $\shC$
is an additive category, $g\colon\stY\arr\stX$ is a map of stacks,
$\shF\in\QCoh\stY$ and $\Gamma\in\L_{R}(\shC,A)$ then any element
of 
\[
(g^{*}\Gamma)_{\shF}=\varinjlim_{(\E,\phi)\in J_{\shF,\shC}}\Gamma_{\E,\phi}
\]
is of the form $x_{\E,\phi}$ for some $(\E,\phi)\in J_{\shF,\shC}$.
Indeed given $(\E,\phi),(\E',\phi')\in J_{\shF,\shC}$ and denoted
by $p\colon\E\oplus\E'\arr\E$, $p'\colon\E\oplus\E'\arr\E'$ the
projections, we have $(\E\oplus\E',\shF\arrdi{\phi\oplus\phi'}g^{*}\E\oplus g^{*}\E')\in J_{\shF,\shC}$
and arrows
\[
(\E\oplus\E',\phi\oplus\phi')\arrdi p(\E,\phi)\comma(\E\oplus\E',\phi\oplus\phi')\arrdi{p'}(\E',\phi')
\]
in $J_{\shF,\shC}$.\end{rem}
\begin{defn}
Given $\shF\in\QCoh\stY$ we also define
\[
J_{\shF,\shC}=\{\text{category of pairs }(\E,\shF\arr g^{*}\E)\text{ for }\E\in\shC\}
\]
Given $\Gamma\in\L_{R}(\shC,A)$ we define the functor $\Gamma\colon J_{\shF,\shC}\arr\Mod A$
as the composition $J_{\shF,\shC}\arrdi{\pr_{1}}\shC\arrdi{\Gamma}\Mod A$.
If $(\E,\phi)\in J_{\shF,\shC}$ we set $\Gamma_{\E,\phi}=\Gamma_{(\E,\phi)}=\Gamma_{\E}$
and given $x\in\Gamma_{\E}$ we set $x_{\E,\phi}\in\Gamma_{\E,\phi}$
for the element $x$ in the $A$-module $\Gamma_{\E,\phi}=\Gamma_{\E}$.
\end{defn}
Notice that $J_{\shF,\shC}$ is small and 
\[
\varinjlim_{(\E,\phi)\in J_{\shF,\shC}}\Gamma_{\shE,\phi}=\frac{{\displaystyle \bigoplus_{(\E,\phi)\in J_{\shF,\shC}}\Gamma_{\E,\phi}}}{\langle x_{\E',\phi'}-\Gamma_{u}(x)_{\E,\phi}\st(\E,\phi)\arrdi u(\E',\phi'),x\in\Gamma_{\E'}\rangle_{A}}\in\Mod A
\]
Denote by $\Delta_{\Gamma,\shF}$ the $A$-module above. Notice that
$J_{\shF,\shC}$ is not a filtered category in general. We remark
that the limit is \emph{direct}, although the functor $\Gamma\colon J_{\shF,\shC}\arr\Mod A$
is contravariant. Formally we should have written $J_{\shF,\shC}^{\op}$
in the expression above, but for simplicity we have preferred this
notation. Set also 
\[
N_{\Gamma,\shF}=\langle x_{\E,\phi}+x_{\E,\psi}-x_{\E,\phi+\psi}\st(\E,\phi),(\E,\psi)\in J_{\shF,\shC},x\in\Gamma_{\E}\rangle_{A}\subseteq\Delta_{\Gamma,\shF}
\]
Given a morphism $\omega\colon\shF\arr\shF'$ in $\shD$ we have a
functor and an $A$-linear map  \[   \begin{tikzpicture}[xscale=2.1,yscale=-0.6]     
\node (A0_0) at (0, 0) {$J_{\shF',\shC}$};     
\node (A0_1) at (1, 0) {$J_{\shF,\shC}$};     
\node (A0_2) at (2, 0) {$\Delta_{\Gamma,\shF'}$};    
\node (A0_3) at (3, 0) {$\Delta_{\Gamma,\shF}$};     
\node (A1_0) at (0, 1) {$(\E,\phi)$};     
\node (A1_1) at (1, 1) {$(\E,\phi\omega)$};
    \node (A1_2) at (2, 1) {$x_{\E,\phi}$};     
\node (A1_3) at (3, 1) {$x_{\E,\phi\omega}$};
    \path (A0_0) edge [->]node [auto] {$\scriptstyle{}$} (A0_1);     \path (A1_0) edge [|->,gray]node [auto] {$\scriptstyle{}$} (A1_1);     \path (A0_2) edge [->]node [auto] {$\scriptstyle{}$} (A0_3);     \path (A1_2) edge [|->,gray]node [auto] {$\scriptstyle{}$} (A1_3);   \end{tikzpicture}   \] which sends $N_{\Gamma,\shF'}$ to $N_{\Gamma,\shF}$.
\begin{prop}
Let $\stY$ be a fibered category, $g\colon\stY\arr\stX$ be a morphism
and $\shD$ be a subcategory of $\QCoh\stY$ such that $g^{*}\shC\subseteq\shD$.
The structures defined above makes $\Delta_{\Gamma,-}\colon\shD\arr\Mod A$
into a contravariant functor and $N_{\Gamma,-}\colon\shD\arr\Mod A$
into a subfunctor of $\Delta_{\Gamma,-}$. Moreover the formation
of $\Delta_{\Gamma,-}$ and $N_{\Gamma,-}$ is natural in $\Gamma\in\L_{R}(\shC,A)$.
The quotient functor
\[
g^{*}\Gamma=\Delta_{\Gamma,-}/N_{\Gamma,-}\colon\shD\arr\Mod A
\]
is $R$-linear, i.e. $g^{*}\Gamma\in\L_{R}(\shD,A)$ and the functor
$g^{*}\colon\L_{R}(\shC,A)\arr\L_{R}(\shD,A)$ is left adjoint to
the restriction $g_{*}\colon\L_{R}(\shD,A)\arr\L_{R}(\shC,A)$, i.e.
we have a natural isomorphism 
\[
\Hom_{\L_{R}(\shC,A)}(\Omega,g_{*}\Gamma)\simeq\Hom_{\L_{R}(\shD,A)}(g^{*}\Omega,\Gamma)\text{ for }\Omega\in\L_{R}(\shC,A),\Gamma\in\L_{R}(\shD,A)
\]
The functor $\Hom_{A}((g^{*}\Gamma)_{\shF},-)\colon\Mod A\arr\Mod A$
is isomorphic to the functor that maps $N\in\Mod A$ to the set of
$R$-linear natural transformations
\[
\Hom_{\stY}(\shF,g^{*}\E)\arr\Hom_{A}(\Gamma_{\E},N)\text{ for }\E\in\shC
\]
Moreover:
\begin{enumerate}
\item If $\shC$ is additive then $g^{*}\Gamma=\Delta_{\Gamma,-}$, that
is $N_{\Gamma,\shF}=0$ for all $\Gamma\in\L_{R}(\shC,A)$ and $\shF\in\shD$;
\item If $\stY=\stX$ and $g=\id_{\stX}$, so that $\shC\subseteq\shD$
and $g_{*}\colon\L_{R}(\shD,A)\arr\L_{R}(\shC,A)$ is the restriction
of functors, then for all $\Gamma\in\L_{R}(\shC,A)$ the functor $g^{*}\Gamma$
extends $\Gamma$, that is the natural transformation $\Gamma\arr g_{*}g^{*}\Gamma=(g^{*}\Gamma)_{|\shC}$
is an isomorphism.
\end{enumerate}
\end{prop}
\begin{proof}
Let $\Gamma\in\L_{R}(\shC,A)$ and for simplicity set $\Delta_{\shF}=\Delta_{\Gamma,\shF}$
and $N_{\shF}=N_{\Gamma,\shF}$ for $\shF\in\shD$. It is easy to
check that $\Delta\colon\shD\arr\Mod A$ is a contravariant functor
and $N\colon\shD\arr\Mod A$ a subfunctor of $\Delta$. Consider $\lambda\in R$
and $\shF\in\shD$. Given $(\E,\phi)\in J_{\shF,\shC}$ we have an
arrow $(\E,\phi)\arrdi{\lambda}(\E,\phi\lambda)$ and therefore 
\[
\Delta_{\lambda}(x_{\E,\phi})=x_{\E,\phi\lambda}=\Gamma_{\lambda\id_{\E}}(x)_{\E,\phi}=\lambda x_{\E,\phi}\text{ in }\Delta_{\shF}\text{ for }x\in\Gamma_{\E}
\]
Therefore $\Delta,N$ and $g^{*}\Gamma$ are all homogeneous, that
is they preserve the multiplication for $\lambda\in R$. Now let $u,v\colon\shF\arr\shF'$
in $\shD$. Given $(\E,\phi)\in J_{\shF',\shC}$ we have
\[
\Delta_{u}(x_{\E,\phi})=x_{\E,\phi u}\comma\Delta_{v}(x_{\E,\phi})=x_{\E,\phi v}\comma\Delta_{u+v}(x_{\E,\phi})=x_{\E,\phi(u+v)}=x_{\E,\phi u+\phi v}\text{ for }x\in\Gamma_{\E}
\]
in $\Delta_{\shF}$. In particular $(\Delta_{u}+\Delta_{v}-\Delta_{u+v})(x_{\E,\phi})\in N_{\shF}$
for all $x\in\Gamma_{\E}$ and we can conclude that $g^{*}\Gamma\colon\shD\arr\Mod A$
is $R$-linear. It is easy to check that $g^{*}$ defines a functor
$\L_{R}(\shC,A)\arr\L_{R}(\shD,A)$.

If $N$ is an $A$-module, a map $\alpha\colon\Delta_{\shF}\arr N$
is given by a coherent collection of $A$-linear maps $\alpha_{\E,\phi}\colon\Gamma_{\E,\phi}\arr N$
for $(\E,\phi)\in J_{\shF,\shC}$. This data is easily seen to correspond
to a natural transformation $\alpha_{\E}\colon\Hom_{\stY}(\shF,g^{*}\E)\arr\Hom_{A}(\Gamma_{\E},N)$
of functors $\shC\arr\sets$. Since the multiplication by $\lambda\in R$
on $\Delta_{\shF}$ coincides with $\Delta_{\lambda}$ we know that
the natural transformation $\alpha_{\E}$ is automatically homogeneous.
Finally $\alpha\colon\Delta_{\shF}\arr N$ factors through $(g^{*}\Gamma)_{\shF}$
if and only if $\alpha_{\E}$ is $R$-linear for all $\E\in\shC$,
that is $\alpha_{\E}$ is a natural transformation of functors $\shC\arr\Mod R$.
Given $\Omega\in\L_{R}(\shD,A)$ we define   \[   \begin{tikzpicture}[xscale=6.0,yscale=-0.6]     \node (A0_0) at (0, 0) {$\Hom_{\L_R(\shC,A)}(\Gamma,g_*\Omega)$};     \node (A0_1) at (1, 0) {$\Hom_{\L_R(\shD,A)}(g^*\Gamma,\Omega)$};     \node (A1_0) at (0, 1) {$\omega$};     \node (A1_1) at (1, 1) {$(\Gamma_{\E,\phi}=\Gamma_\E \arrdi{\omega_\E}(g_*\Omega)_\E=\Omega_{g^*\E}\arrdi{\Omega_\phi}\Omega_\shF)$};     \path (A0_0) edge [->]node [auto] {$\scriptstyle{\xi}$} (A0_1);     \path (A1_0) edge [|->,gray]node [auto] {$\scriptstyle{}$} (A1_1);   \end{tikzpicture}   \] that
is $\xi(\omega)_{\E,\phi}(x_{\E,\phi})=\Omega_{\phi}(\omega_{\E}(x))$
for $x\in\Gamma_{\E}$. Moreover we have a natural transformation
$\Gamma\arr g_{*}g^{*}\Gamma$ given by
\[
\Gamma_{\E}=\Gamma_{\E,\id_{g^{*}\E}}\arr\Delta_{g^{*}\E}=(g_{*}\Delta)_{\E}\text{ for }\E\in\shC
\]
which yields a map   \[   \begin{tikzpicture}[xscale=5.0,yscale=-0.6]     \node (A0_0) at (0, 0) {$\Hom_{\L_R(\shD,A)}(g^*\Gamma,\Omega)$};     \node (A0_1) at (1, 0) {$\Hom_{\L_R(\shC,A)}(\Gamma,g_*\Omega)$};     \node (A1_0) at (0, 1) {$\alpha$};     \node (A1_1) at (1, 1) {$(\Gamma\arr g_*g^*\Gamma\arrdi{g_*\alpha}g_*\Omega)$};     \path (A0_0) edge [->]node [auto] {$\scriptstyle{\zeta}$} (A0_1);     \path (A1_0) edge [|->,gray]node [auto] {$\scriptstyle{}$} (A1_1);   \end{tikzpicture}   \] 
so that $\zeta(\alpha)_{\E}(x)=\alpha_{\E,\id_{g^{*}\E}}(x_{\E,\id_{g^{*}\E}})\in\Omega_{g^{*}\E}=(g_{*}\Omega)_{\E}$
for $x\in\Gamma_{\E}$. We claim that the above maps are well defined
and inverses of each other. We start showing that $\xi(\omega)$ is
well defined. Given $x\in\Gamma_{\E'}$, $\shF\in\shD$ and $(\E,\phi)\arrdi u(\E',\phi')$
in $J_{\shF,\shC}$, so that $\phi'=g^{*}u\circ\phi$, we have
\[
\xi(\omega)_{\E,\phi}(\Gamma_{u}(x)_{\E,\phi})=\Omega_{\phi}(\omega_{\E}(\Gamma_{u}(x)))=\Omega_{\phi}(\Omega_{g^{*}u}(\omega_{\E'}(x)))=\Omega_{\phi'}(\omega_{\E'}(x))=\xi(\omega)_{\E',\phi'}(x_{\E',\phi'})
\]
in $\Omega_{\shF}$ and this proves that $\xi(\omega)$ defines a
map $\Delta_{\shF}\arr\Omega_{\shF}$. Since 
\[
\xi(\omega)(x_{\E,\phi}+x_{\E,\psi}-x_{\E,\phi+\psi})=(\Omega_{\phi}+\Omega_{\psi}-\Omega_{\phi+\psi})(\omega_{\E}(x))=0
\]
for $(\E,\phi),(\E,\psi)\in J_{\shF,\shC}$ and $\shF\in\shD$, we
see that $\xi(\omega)$ defines also a map $(g^{*}\Gamma)_{\shF}\arr\Omega_{\shF}$.
It is now easy to prove that $\xi(\omega)$ is a natural transformation.
The naturality of $\Gamma\arr g_{*}g^{*}\Gamma$ follows from the
fact that given $\E\arrdi u\E'$ in $\shC$ we always have an arrow
$(\E,\id_{g^{*}\E})\arrdi u(\E',g^{*}u)$ in $J_{g^{*}\E,\shC}$.

The proof that $\zeta\circ\xi=\id$ is really elementary. Also the
equality $\xi\circ\zeta=\id$ is very simple, taking into account
that if $\alpha\colon g^{*}\Gamma\arr\Omega$ and $(\E,\phi)\in J_{\shF,\shC}$
the naturality of $\alpha$ on the map $\phi\colon\shF\arr g^{*}\E$
in $\shD$ yields a commutative diagram   \[   \begin{tikzpicture}[xscale=2.0,yscale=-1.1]     \node (A0_0) at (0, 0) {$\Gamma_{\E,\id_{g^*\E}}$};     \node (A0_1) at (1, 0) {$(g^*\Gamma)_{g^*\E}$};     \node (A0_2) at (2, 0) {$\Omega_{g^*\E}$};     \node (A1_0) at (0, 1) {$\Gamma_{\E,\phi}$};     \node (A1_1) at (1, 1) {$(g^*\Gamma)_\shF$};     \node (A1_2) at (2, 1) {$\Omega_\shF$};     \path (A0_0) edge [->]node [auto] {$\scriptstyle{}$} (A0_1);     \path (A0_1) edge [->]node [auto] {$\scriptstyle{(g^*\Gamma)_\phi}$} (A1_1);     \path (A1_0) edge [->]node [auto] {$\scriptstyle{}$} (A1_1);     \path (A1_1) edge [->]node [auto] {$\scriptstyle{\alpha_\shF}$} (A1_2);     \path (A0_2) edge [->]node [auto] {$\scriptstyle{\Omega_\phi}$} (A1_2);     \path (A0_0) edge [->,bend right=40]node [auto] {$\scriptstyle{\alpha_{\E,\id_{g^*\E}}}$} (A0_2);     \path (A0_0) edge [->]node [auto] {$\scriptstyle{\id}$} (A1_0);     \path (A0_1) edge [->]node [auto,swap] {$\scriptstyle{\alpha_{g^*\E}}$} (A0_2);     \path (A1_0) edge [->,bend left=40]node [auto,swap] {$\scriptstyle{\alpha_{\E,\phi}}$} (A1_2);   \end{tikzpicture}   \] It
remains to prove the last two points.

$1)$ Let $(\E,\phi),(\E,\psi)\in J_{\shF,C}$, $x\in\Gamma_{\E}$.
We denote by $p_{1},p_{2},s\colon\E\oplus\E\arr\E$ the first projection,
second projection and the sum map respectively. Note that $(\E\oplus\E,\shF\arrdi{\phi\oplus\psi}g^{*}\E\oplus g^{*}\E)\in J_{\shF,\shC}$
and that we have arrows in $J_{\shF,\shC}$
\[
(\E\oplus\E,\phi\oplus\psi)\arrdi{p_{1}}(\E,\phi)\comma(\E\oplus\E,\phi\oplus\psi)\arrdi{p_{2}}(\E,\psi)\comma(\E\oplus\E,\phi\oplus\psi)\arrdi s(\E,\phi+\psi)
\]
Therefore
\[
x_{\E,\phi}+x_{\E,\psi}-x_{\E,\phi+\psi}=(\Gamma_{p_{1}}(x)+\Gamma_{p_{2}}(x)-\Gamma_{s}(x))_{\E\oplus\E,\phi\oplus\psi}=\Gamma_{p_{1}+p_{2}-s}(x)_{\E\oplus\E,\phi\oplus\psi}=0
\]
in $\Delta_{\shF}$ for all $x\in\Gamma_{\E}$ because $p_{1+}p_{2}=s$.

$2)$ If $\E'\in\shC$ then $J_{\E',\shC}$ has an initial object,
namely $(\E',\id_{\E'})$. Therefore the map $\Gamma_{\E'}=\Gamma_{\E',\id_{\E'}}\arr\Delta_{\E'}$
is an isomorphism. It remains to shows that $N_{\E'}=0$. Given $u,v\colon\E'\arr\E$
we have arrows
\[
(\E',\id_{\E'})\arrdi u(\E,u)\comma(\E',\id_{\E'})\arrdi v(\E,v)\comma(\E',\id_{\E'})\arrdi{u+v}(\E,u+v)
\]
in $J_{\E',\shC}$ and therefore 
\[
x_{\E,u}+x_{\E,v}-x_{\E,u+v}=(\Gamma_{u}(x)+\Gamma_{v}(x)-\Gamma_{u+v}(x))_{\E',\id_{\E'}}=0
\]
in $\Delta_{\E'}$ for all $x\in\Gamma_{\E}$.\end{proof}
\begin{rem}
\label{rem: map from as natural transformations} If $U\arrdi g\stX$
is a map from a scheme, then maps $\odi U\arr g^{*}\E$ corresponds
to elements of $\E(U)=\E(U\arrdi g\stX)$. In particular $J_{\odi U,\shC}$
can be described as the category of pairs $(\E,\phi)$ where $\E\in\shC$
and $\phi\in\E(U)$, while an arrow $(\E,\phi)\arrdi u(\overline{\E},\overline{\phi})$
is a morphism $\E\arrdi u\overline{\E}$ such that $u(\phi)=\overline{\phi}$. 

Moreover the functor $\Hom_{\Hl^{0}(\odi U)\otimes_{R}A}(\widetilde{\shF}_{\Gamma,\shC}(U),-)\colon\Mod(\Hl^{0}(\odi U)\otimes_{R}A)\arr\Mod(\Hl^{0}(\odi U)\otimes_{R}A)$
is isomorphic to the functor that sends $N\in\Mod(\Hl^{0}(\odi U)\otimes_{R}A)$
to the set of $\Hl^{0}(\odi U)$-linear natural transformations
\[
\E(U)\arr\Hom_{A}(\Gamma_{\E},N)\text{ for }\E\in\shC
\]
Indeed by \ref{thm:adjoint of push forward} an $A$-linear map $\widetilde{\shF}_{\Gamma,\shC}(U)\arr N$
corresponds to an $R$-linear natural transformation as above. Since
the $\Hl^{0}(\odi U)$-module structure on $\widetilde{\shF}_{\Gamma,\shC}(U)$
is given by $bx_{\E,\phi}=x_{\E,b\phi}$ where $b\in\Hl^{0}(\odi U)$,
$(\E,\phi)\in J_{\odi U,\shC}$ and $x\in\Gamma_{\E}$ we get the
desired result.
\end{rem}
If $N$ is an $\Hl^{0}(\odi V)\otimes_{R}A$-module then a $\Hl^{0}(\odi V)$-linear
natural transformation $\E(V)\arr\Hom_{A}(\Gamma_{\E},N)$ for $\E\in\shC$
induces an $\Hl^{0}(\odi U)$-linear natural transformation 
\[
\E(U)\arrdi{\E(\zeta)}\E(V)\arr\Hom_{A}(\Gamma_{\E},N)\text{ for }\E\in\shC
\]
By \ref{rem: map from as natural transformations} this induces an
$\Hl^{0}(\odi U)\otimes_{R}A$-linear morphism $\widetilde{\shF}_{\Gamma,\shC}(\zeta)\colon\widetilde{\shF}_{\Gamma,\shC}(U)\arr\widetilde{\shF}_{\Gamma,\shC}(V)$.
On generators we have
\[
\widetilde{\shF}_{\Gamma,\shC}(\zeta)(x_{\E,\phi})=x_{\E,\phi_{|V}}\text{ where }\E\in\shC,\phi\in\E(U),x\in\Gamma_{\E}
\]
where $\phi_{|V}$ denotes $\E(\zeta)(\phi)$.
\begin{defn}
Let $\Gamma\in\L_{R}(\shC,A)$, $g\colon U\arr\stX$ be a map from
a scheme and consider the functor $g^{*}\colon\L_{R}(\shC,A)\arr\L_{R}(g^{*}\shC,A)$.
We set
\[
\widetilde{\shF}_{\Gamma,\shC}(U\arrdi g\stX)=(g^{*}\Gamma)_{\odi U}=\int^{\E\in\shC}\Gamma_{\E}\otimes_{R}\E(U\arrdi g\stX)\in\Mod(\Hl^{0}(\odi U)\otimes_{R}A)
\]
where we have used the $\Hl^{0}(\odi U)$-linear natural isomorphism
$\Hom_{U}(\odi U,g^{*}\E)\simeq\E(U\arrdi g\stX)$. When $g$ is given
we will simply write $\widetilde{\shF}_{\Gamma,\shC}(U)$.
\end{defn}
Consider now $V\arrdi{\overline{g}}\stX$, $U\arrdi g\stX$ and an
arrow $\zeta\colon\overline{g}\arr g$ in $\Sch/\stX$. We obtain
a $\Hl^{0}(\odi U)$-linear natural transformation $\E(U)\arrdi{\E(\zeta)}\E(V)$
for $\E\in\shC$ and therefore an $\Hl^{0}(\odi U)$-linear morphism
$\widetilde{\shF}_{\Gamma,\shC}(\zeta)\colon\widetilde{\shF}_{\Gamma,\shC}(U)\arr\widetilde{\shF}_{\Gamma,\shC}(V)$.
\begin{lem}
\label{lem:properties of the presheaf FGammaC} The above structure
makes $\widetilde{\shF}_{\Gamma,\shC}$ into a presheaf of $\odi{\stX}$-modules
over $\stX$ with an $A$-module structure (see \ref{rem: A module structures of odix modules}).
If $\stY\arrdi g\stX$ is a map of fibered categories, then we have
an isomorphism
\[
g^{*}\widetilde{\shF}_{\Gamma,\shC}\simeq\widetilde{\shF}_{g^{*}\Gamma,g^{*}\shC}
\]
of presheaves of $\odi{\stY}$-modules with $A$-module structures.\end{lem}
\begin{proof}
It is very easy to verify that $\widetilde{\shF}_{\Gamma,\shC}$ is
a presheaf of $\odi{\stX}$-modules with $A$-module structure. Moreover
if $U\arrdi h\stY$ is a map we have
\[
g^{*}\widetilde{\shF}_{\Gamma,\shC}(U\arrdi h\stY)=\widetilde{\shF}_{\Gamma,\shC}(U\arrdi h\stY\arrdi g\stX)=((gh)^{*}\Gamma)_{\odi U}\simeq(h^{*}(g^{*}\Gamma))_{\odi U}=\widetilde{\shF}_{g^{*}\Gamma,g^{*}\shC}(U\arrdi h\stY)
\]
which yields the isomorphism for the last part of the statement.\end{proof}
\begin{defn}
Let $\Gamma\in\L_{R}(\shC,A)$. We define 
\[
\shF_{\Gamma,\shC}\colon(\Sch/\stX)^{\text{op}}\arr\Mod A
\]
as the Zariski sheafification of the functor $\widetilde{\shF}_{\Gamma,\shC}$. 
\end{defn}
We are going to prove that $\shF_{\Gamma,\shC}$ is a quasi-coherent
sheaf and that the association $\Gamma\arr\shF_{\Gamma,\shC}$ is
the left adjoint of the Yoneda functor $\Omega^{*}\colon\QCoh_{A}\stX\arr\L_{R}(\shC,A)$
or, in other words, the sheafification functor associated with $\shC$.
\begin{lem}
\label{lem:key lemma for the sheafification of FGammaC} If $\Gamma\in\L_{R}(\shC,A)$
then the functor $\widetilde{\shF}_{\Gamma,\shC}\colon\stX^{\text{op}}\arr\Ab$
restricted to $\Aff/\stX$ is quasi-coherent.\end{lem}
\begin{proof}
Let $B,B'$ be $R$-algebras, $\Spec B\arrdi{\xi}\stX$, $\Spec B'\arrdi{\eta}\stX$
be maps and $\zeta\colon\eta\arr\xi$ be an arrow in $\stX$. If $N\in\Mod B'\otimes_{R}A$
and we apply $\Hom_{B'\otimes_{R}A}(-,N)$ to the map $\widetilde{\shF}_{\Gamma,\shC}(\zeta)\otimes1\colon\widetilde{\shF}_{\Gamma,\shC}(B)\otimes_{B}B'\arr\widetilde{\shF}_{\Gamma,\shC}(B')$
we obtain a map sending an $A$-linear natural transformation $\Gamma_{\E}\arr\Hom_{B'}(\E(B'),N)$
to the $A$-linear natural transformation $\Gamma_{\E}\arr\Hom_{B}(\E(B),N)$
obtained by composing with $\E(\zeta)\colon\E(B)\arr\E(B')$. Since
$\E(\zeta)\otimes1\colon\E(B)\otimes_{B}B'\arr\E(B')$ is an isomorphism
for all $\E\in\shC$, we can conclude that the map $\Hom_{B'\otimes_{R}A}(\widetilde{\shF}_{\Gamma,\shC}(\zeta)\otimes1,N)$
is an isomorphism. By Yoneda's lemma it follows that $\widetilde{\shF}_{\Gamma,\shC}(\zeta)\otimes1$
is an isomorphism, as required.\end{proof}
\begin{cor}
\label{prop:FGammaC is a quasi-coherent sheaf} If $\Gamma\in\L_{R}(\shC,A)$
then $\shF_{\Gamma,\shC}\in\QCoh_{A}\stX$ and $\shF_{\Gamma,\shC}$
and $\widetilde{\shF}_{\Gamma,\shC}$ coincides on the category of
affine schemes over $\stX$. The association
\[
\shF_{-,\shC}\colon\L_{R}(\stC,A)\arr\QCoh_{A}\stX
\]
defines a functor. If $g\colon\stY\arr\stX$ is a map of fibered categories
we have a natural isomorphism $g^{*}\shF_{-,\shC}\arr\shF_{g^{*}-,g^{*}\shC}$.\end{cor}
\begin{proof}
Taking into account \ref{lem:properties of the presheaf FGammaC},
since $(\widetilde{\shF}_{\Gamma,\shC})_{|\Aff/\stX}$ is a sheaf
if follows that $(\shF_{\Gamma,\shC})_{|\Aff/\stX}=(\widetilde{\shF}_{\Gamma,\shC})_{|\Aff/\stX}$
and that $\shF_{\Gamma,\shC}$ is quasi-coherent. It is elementary
to check that $\shF_{*,\shC}$ defines a functor as in the statement.
Finally the last part follows from \ref{lem:properties of the presheaf FGammaC}.
\end{proof}
Let $\Gamma\in\L_{R}(\shC,A)$. We define a natural transformation
$\gamma_{\Gamma,-}\colon\Gamma\arr\Omega^{\shF_{\Gamma,\shC}}$ of
functors in $\L_{R}(\shC,A)$ by   \[   \begin{tikzpicture}[xscale=3.8,yscale=-0.6]     
\node (A0_0) at (0.2, 0) {$\Gamma_\E$};   
\node (A1_0) at (0.2, 1) {$x$}; 
\node (A0_1) at (1, 0) {$\Hom_{\stX}(\E,\shF_{\Gamma,\shC})$};     
\node (A1_1) at (1, 1) {$(\phi \longmapsto x_{\E,\phi})$};     
\path (A0_0) edge [->]node [auto] {$\scriptstyle{}$} (A0_1);     
\path (A1_0) edge [|->,gray]node [auto] {$\scriptstyle{}$} (A1_1);        
\end{tikzpicture}   \]  It is not difficult to check that this map is well defined, $A$-linear
and natural in $\E\in\shC$. 

Let $\shG\in\QCoh_{A}\stX$. We define a morphism of quasi-coherent
sheaves $\delta_{\shG}\colon\shF_{\Omega^{\shG},\shC}\arr\shG$ by
the expression 
\[
\delta_{\shG}(x_{\E,\phi})=x(\phi)
\]
where $(\E,\phi)\in J_{\odi U,\shC}$, $U\arr\stX$ is a map and $x\in\Omega_{\E,\phi}^{\shG}=\Hom_{\stX}(\E,\shG)$.
The map $\delta_{\shG}$ is the map induced by the morphism $\widetilde{\shF}_{\Omega^{\shG},\shC}(U)\arr\shG(U)$
corresponding to the $\Hl^{0}(\odi U)$-linear natural transformation
\[
\E(U)\arr\Hom_{A}(\Omega_{\E}^{\shG},\shG(U))=\Hom_{A}(\Hom_{\stX}(\E,\shG),\shG(U))\comma\phi\longmapsto(x\longmapsto x(\phi))
\]
It is now easy to verify that $\delta_{\shG}$ is well defined.
\begin{prop}
\label{prop:Fstar and Omegastar are adjoint} The functor $\shF_{*,\shC}\colon\L_{R}(\shC,A)\arr\QCoh_{A}\stX$
is left adjoint to the functor $\Omega^{*}\colon\QCoh_{A}\stX\arr\L_{R}(\shC,A)$.\end{prop}
\begin{proof}
The natural transformations $\gamma_{\Gamma,-}\colon\Gamma\arr\Omega^{\shF_{\Gamma,\shC}}$
for $\Gamma\in\L_{R}(\shC,A)$ and the map $\delta_{\shG}\colon\shF_{\Omega^{\shG},\shC}\arr\shG$
for $\shG\in\QCoh_{A}\stX$ induce maps
\[
\Hom(\shF_{\Gamma,\shC},\shG)\arrdi{\xi}\Hom(\Gamma,\Omega^{\shG})\text{ and }\Hom(\Gamma,\Omega^{\shG})\arrdi{\eta}\Hom(\shF_{\Gamma,\shC},\shG)
\]
and we are going to prove that they are inverses of each other. In
what follows $\psi$ will be a generic section of $\E$ over some
map $\Spec B\arr\stX$. Given $\alpha\colon\Gamma\arr\Omega^{\shG}=\Hom(-,\shG)$
the map $\eta(\alpha)\colon\shF_{\Gamma,\shC}\arr\shG$ is given by
the collection of maps   \[   \begin{tikzpicture}[xscale=3.0,yscale=-0.6]     
\node (bo) at (-0.4, 0) {$\eta(\alpha)_{\E,\psi}\ \ \colon$};     
\node (A0_0) at (0, 0) {$\Gamma_{\E,\psi}$};     
\node (A0_1) at (1, 0) {$\Omega^\shG_{\E,\psi}=\Hom(\E,\shG)$};     \node (A0_2) at (2, 0) {$\shG$};     \node (A1_1) at (1, 1) {$\phi$};     \node (A1_2) at (2, 1) {$\phi(\psi)$};     \path (A0_0) edge [->]node [auto] {$\scriptstyle{\alpha_\E}$} (A0_1);     \path (A0_1) edge [->]node [auto] {$\scriptstyle{}$} (A0_2);     \path (A1_1) edge [|->,gray]node [auto] {$\scriptstyle{}$} (A1_2);   \end{tikzpicture}   \]  so that $\eta(\alpha)_{\E,\psi}(x_{\E,\psi})=[\alpha_{\E}(x)](\psi)$
for $x\in\Gamma_{\E}$. Instead given $\beta\colon\shF_{\Gamma,\shC}\arr\shG$
the map $\xi(\beta)\colon\Gamma\arr\Omega^{\shG}=\Hom(-,\shG)$ is
given by maps   \[   \begin{tikzpicture}[xscale=3.8,yscale=-0.6]     
\node (bo) at (-0.05, 0) {$\xi(\beta)_\E \ \ \colon $};   
\node (A0_0) at (0.2, 0) {$\Gamma_\E$};   
\node (A1_0) at (0.2, 1) {$x$}; 
\node (A0_1) at (1, 0) {$\Hom(\E,\shF_{\Gamma,\shC})$};     \node (A0_2) at (2, 0) {$\Hom(\E,\shG)$};     
     \node (A1_1) at (1, 1) {$(\psi \longmapsto x_{\E,\psi})$};     \node (A1_2) at (2, 1) {$(\psi \longmapsto \beta_{\E,\psi}(x_{\E,\psi}))$};     \path (A0_0) edge [->]node [auto] {$\scriptstyle{}$} (A0_1);     \path (A1_0) edge [|->,gray]node [auto] {$\scriptstyle{}$} (A1_1);     \path (A0_1) edge [->]node [auto] {$\scriptstyle{}$} (A0_2);     \path (A1_1) edge [|->,gray]node [auto] {$\scriptstyle{}$} (A1_2);   \end{tikzpicture}   \]  so that $[\xi(\beta)_{\E}(x)](\psi)=\beta_{\E,\psi}(x_{\E,\psi})$
for $x\in\Gamma_{\E}$. With those descriptions it is easy to check
that $\xi$ and $\eta$ are inverses of each other. \end{proof}
\begin{prop}
\label{prop:presentation of Fgammashc} If $\Gamma\in\L_{R}(\shC,A)$
the sequence   \[   \begin{tikzpicture}[xscale=4.3,yscale=-0.8]     
\node (A0_0) at (-0.15, 0) {$x\otimes\alpha\otimes y$};     
\node (A0_1) at (1, 0) {$x\otimes \alpha(y)-\Gamma_\alpha (x)\otimes y$};     
\node (A1_0) at (-0.15, 1) {$\displaystyle{\bigoplus_{\E,\overline \E\in \shC} \Gamma_{\overline \E}\otimes \Hom(\E,\overline\E)\otimes \E}$};     
\node (A1_1) at (1, 1) {$\displaystyle{\bigoplus_{\E \in \shC} \Gamma_\E\otimes  \E}$};     \node (A1_2) at (1.7, 1) {$\shF_{\Gamma,\shC}$};     
\node (A1_3) at (2.2, 1) {$0$};     \node (A2_1) at (1, 2) {$x\otimes \phi$};     
\node (A2_2) at (1.7, 2) {$x_{\E,\phi}$};     
\path (A0_0) edge [|->,gray]node [auto] {$\scriptstyle{}$} (A0_1);     
\path (A1_0) edge [->]node [auto] {$\scriptstyle{\omega_\Gamma}$} (A1_1);     \path (A2_1) edge [|->,gray]node [auto] {$\scriptstyle{}$} (A2_2);     \path (A1_2) edge [->]node [auto] {$\scriptstyle{}$} (A1_3);     \path (A1_1) edge [->]node [auto] {$\scriptstyle{}$} (A1_2);   \end{tikzpicture}   \]  is exact.\end{prop}
\begin{proof}
By \ref{prop:adjoint of Omegastar with the cokernel} we have an isomorphism
\[
\Hom(\Gamma,\Omega^{\shG})\arr\Hom(\Coker\omega_{\Gamma},\shG)
\]
The map $\Gamma\arr\Omega^{\shF_{\Gamma,\shC}}$ induces a map $\Coker\omega_{\Gamma}\arr\shF_{\Gamma,\shC}$
which is an isomorphism thanks to \ref{prop:Fstar and Omegastar are adjoint}.
A direct check shows that this map coincides with the second map in
the statement.\end{proof}
\begin{example}
\label{exam: semisimple decomposition}Assume that $\Hom(\E,\overline{\E})=0$
for all $\E\neq\overline{\E}\in\shC$ and $\End_{\stX}(\E)=\Hl^{0}(\odi{\stX})$
for all $\E\in\shC$. It is easy to check that 
\[
\L_{R}(\shC,A)\simeq\prod_{\E\in\shC}\Mod(\Hl^{0}(\odi{\stX})\otimes_{R}A)
\]
and, by \ref{prop:presentation of Fgammashc} we see that $\shF_{\Gamma,\shC}\simeq\bigoplus_{\E\in\shC}\Gamma_{\E}\otimes\E\text{ for }\Gamma\in\L_{R}(\shC,A)$.
Thus $\QCoh\stX_{A}\arrdi{\Omega^{*}}\L_{R}(\shC,A)\arrdi{\shF_{*,\shC}}\QCoh\stX_{A}$
correspond to   \[   \begin{tikzpicture}[xscale=4.1,yscale=-0.8]     \node (A0_1) at (1, 0) {$(M_\E)_{\E\in \shC}$};     \node (A0_2) at (2, 0) {$\displaystyle\bigoplus_{\E\in \shC} M_\E \otimes_R \E$};     \node (A1_0) at (0, 1) {$\QCoh \stX_A$};     \node (A1_1) at (1, 1) {$\displaystyle \prod_{\E\in \shC} \Mod(\Hl^0(\odi \stX)\otimes_R A)$};     \node (A1_2) at (2, 1) {$\QCoh \stX_A$};     \node (A2_0) at (0, 2) {$\shF$};     \node (A2_1) at (1, 2) {$(\Hom_\stX(\E,\shF))_{\E\in \shC}$};     \path (A0_1) edge [|->,gray]node [auto] {$\scriptstyle{}$} (A0_2);     \path (A1_0) edge [->]node [auto] {$\scriptstyle{}$} (A1_1);     \path (A2_0) edge [|->,gray]node [auto] {$\scriptstyle{}$} (A2_1);     \path (A1_1) edge [->]node [auto] {$\scriptstyle{}$} (A1_2);   \end{tikzpicture}   \] 
\end{example}

\begin{example}
\label{examp: the case C equal OX} Let $\shC_{\stX}=\{\odi{\stX}\}$.
By \ref{exam: semisimple decomposition} we see that $\L_{R}(\shC_{\stX},A)\simeq\Mod(\odi{\stX}(\stX)\otimes_{R}A)$
and $\Omega^{*}$ and $\shF_{*,\shC_{\stX}}$ correspond to $\Hl^{0}(-)$
and $\pi_{A}^{*}$ respectively, where $\pi_{A}\colon\stX_{A}\arr\Spec(\odi{\stX}(\stX)\otimes_{R}A)$
is the structure morphism.
\end{example}
Sometimes it will be useful to assume that the category $\shC$ is
additive. We conclude this section by showing that we can always reduce
the problems to this case.
\end{comment}

\end{proof}

\subsection{Sheafifying $R$-linear monoidal functors.}

In this section we show how ``ring structures'' on a quasi-coherent
sheaf over $\stX$ correspond to ``monoidal'' structures on the
corresponding Yoneda functor.

We start setting up some definitions:
\begin{defn}
\label{def: pseudo monoidal commutativi associative} Let $\catC$
and $\catD$ be $R$-linear symmetric monoidal categories. A (contravariant)
\emph{pseudo-monoidal }functor $\Omega\colon\catC\arr\catD$ is an
$R$-linear (and contravariant) functor together with a natural transformation
\[
\iota_{V,W}^{\Omega}\colon\Omega_{V}\otimes\Omega_{W}\arr\Omega_{V\otimes W}\text{ for }V,W\in\catC
\]
A (contravariant) pseudo-monoidal functor $\Omega\colon\catC\arr\catD$
is
\begin{enumerate}
\item \emph{\label{enu:commutative for functors} symmetric} or commutative
if for all $V,W\in\catC$ the following diagram is commutative$$
\begin{tikzpicture}[xscale=2.7,yscale=-1.2]     \node (A0_1) at (1, 0) {$\Omega_V\otimes \Omega_W$};     \node (A0_2) at (2, 0) {$\Omega_{V\otimes W}$};     \node (A1_1) at (1, 1) {$\Omega_W\otimes \Omega_V$};     \node (A1_2) at (2, 1) {$\Omega_{W\otimes V}$};     \path (A0_1) edge [->]node [auto] {$\scriptstyle{\iota_{V,W}^\Omega}$} (A0_2);     \path (A0_2) edge [->]node [auto] {$\scriptstyle{}$} (A1_2);     \path (A1_1) edge [->]node [auto] {$\scriptstyle{\iota_{W,V}^\Omega}$} (A1_2);     \path (A0_1) edge [->]node [auto] {$\scriptstyle{}$} (A1_1);   \end{tikzpicture}
$$where the vertical arrows are the obvious isomorphisms;
\item \emph{\label{enu:associative for functors} associative} if for all
$V,W,Z\in\catC$ the following diagram is commutative $$
 \begin{tikzpicture}[xscale=4,yscale=-1.2]     \node (A0_0) at (0, 0) {$\Omega_V\otimes \Omega_W\otimes \Omega_Z$};     \node (A0_1) at (1, 0) {$\Omega_{V\otimes W}\otimes \Omega_Z$};     \node (A1_0) at (0, 1) {$\Omega_V\otimes \Omega_{W\otimes Z}$};     \node (A1_1) at (1, 1) {$\Omega_{V\otimes W\otimes Z}$};     \path (A0_0) edge [->]node [auto] {$\scriptstyle{\iota_{V,W}^\Omega \otimes \id}$} (A0_1);     \path (A1_0) edge [->]node [auto] {$\scriptstyle{\iota_{V,W\otimes Z}^\Omega}$} (A1_1);     \path (A0_0) edge [->]node [auto] {$\scriptstyle{\id\otimes \iota_{W,Z}^\Omega}$} (A1_0);     \path (A0_1) edge [->]node [auto] {$\scriptstyle{\iota_{V\otimes W,Z}^\Omega}$} (A1_1);   \end{tikzpicture}
$$
\end{enumerate}
If $I$ and $J$ are the unit objects of $\catC$ and $\catD$ respectively,
a unity for $\Omega$ is a morphism $1\colon J\arr\Omega_{I}$ such
that, for all $V\in\catC$, the compositions
\[
\Omega_{V}\otimes J\arrdi{\id\otimes1}\Omega_{V}\otimes\Omega_{I}\arrdi{\iota_{V,I}^{\Omega}}\Omega_{V\otimes I}\arr\Omega_{V}\text{ and }J\otimes\Omega_{V}\arrdi{1\otimes\id}\Omega_{I}\otimes\Omega_{V}\arrdi{\iota_{I,V}^{\Omega}}\Omega_{I\otimes V}\arr\Omega_{V}
\]
coincide with the natural isomorphisms $\Omega_{V}\otimes J\arr\Omega_{V}$
and $J\otimes\Omega_{V}\arr\Omega_{V}$ respectively. A (contravariant)\emph{
monoidal} functor $\Omega\colon\catC\arr\catD$ is a symmetric and
associative pseudo-monoidal (contravariant) functor with a unity $1$.
\end{defn}
A morphism of pseudo-monoidal functors $(\Omega,\iota^{\Omega})\arr(\Gamma,\iota^{\Gamma})$,
called a monoidal morphism or tranformation, is a natural transformation
$\Omega\arr\Gamma$ which commutes with the monoidal structures $\iota^{*}$.
A morphism of monoidal functors is a monoidal transformation preserving
the unities.
\begin{defn}
\label{def: monoidal subcategories} We define the categories:
\begin{itemize}
\item $\Rings_{A}\stX$, whose objects are $\alB\in\QCoh_{A}\stX$ with
an $A$-linear map $m\colon\alB\otimes_{\odi{\stX_{A}}}\alB\arr\alB$,
called the multiplication;
\item $\QAlg_{A}\stX$, as the (not full) subcategory of $\Rings_{A}\stX$
whose objects are $\alB$ with a commutative, associative multiplication
with a unity and the arrows are morphisms preserving unities;
\end{itemize}
We also set $\Rings\stX=\Rings_{R}\stX$ and $\QAlg\stX=\QAlg_{R}\stX$. 

Let $\shD$ be a monoidal subcategory of $\QCoh\stX$, that is a subcategory
such that $\odi{\stX}\in\shD$ and for all $\E,\E'\in\shD$ we have
$\E\otimes\E'\in\shD$. We define the category $\PML_{R}(\shD,A)$
(resp. $\ML_{R}(\shD,A)$), whose objects are $\Gamma\in\L_{R}(\shD,A)$
with a pseudo-monoidal (resp. monoidal) structure.\end{defn}
\begin{rem}
If $q\colon\stX_{A}\arr\stX$ is the projection, the equivalence $\QCoh\stX_{A}\arr\QCoh_{A}\stX$
extends to an equivalence 
\[
q_{*}\colon\QRings\stX_{A}\arr\QRings_{A}\stX
\]
Indeed, if $\shG\in\QCoh\stX_{A}$ then $q_{*}(\shG\otimes_{\odi{\stX_{A}}}\shG)=q_{*}\shG\otimes_{\odi{\stX_{A}}}q_{*}\shG$
(see \ref{rem: remark on the tensor product over A}). We will use
the following notation, which is somehow implicit in the definition
of $\QAlg_{A}\stX$: a sheaf $\alB\in\QRings_{A}\stX$ with $\alB\simeq q_{*}\alB'$
is associative (resp. commutative, has a unity, ...) if $\alB'$ has
the same property.

If $\alB\in\QRings_{A}\stX$ with multiplication $m$, then the composition
$\alB\otimes_{\odi{\stX}}\alB\arr\alB\otimes_{\odi{\stX_{A}}}\alB\arrdi m\alB$
induces a ring structure on $\alB$ as an $\odi{\stX}$-module, i.e.
$\alB\in\QRings\stX$. Moreover $\alB\in\QRings_{A}\stX$ is associative
(resp. commutative, has a unity) if and only if $\alB$ has the same
property. If $\alB\in\QAlg_{A}\stX$ we can form the relative spectrum
$\Spec\alB$ over $\stX_{A}$ and also over $\stX$. The final result
is the same.
\end{rem}

\begin{rem}
For $\alB\in\Rings_{A}\stX$ or $\Gamma\in\PMon_{R}(\shD,A)$ having
a unity is a property, not an additional datum. Indeed in both cases
unities are unique.
\end{rem}
Let $\shD$ be a monoidal subcategory of $\QCoh\stX$. If $\alB\in\Rings_{A}\stX$
with multiplication $m$, we endow $\Omega^{\alB}\in\L_{R}(\shD,A)$
with the pseudo monoidal structure
\[
\iota^{\alB}\colon\Hom(\E,\alB)\otimes_{A}\Hom(\overline{\E},\alB)\arr\Hom(\E\otimes\overline{\E},\alB\otimes_{\odi{\stX_{A}}}\alB)\arrdi{\Hom(\E\otimes\overline{\E},m)}\Hom(\E\otimes\overline{\E},\alB)\comma\E,\overline{\E}\in\shD
\]
that is $\iota_{\E,\overline{\E}}^{\alB}(\phi\otimes\psi)=m\circ(\phi\otimes\psi)$.
If $1\in\alB$ is a unity then we set $1_{\alB}=1\in\Omega_{\odi{\stX}}^{\alB}=\Hl^{0}(\alB)$
\begin{prop}
\label{prop:Omegastar on algebras} The structures defined above yield
an extension of the functor $\Omega^{*}\colon\QCoh_{A}\stX\arr\L_{R}(\shD,A)$
to a functor $\Omega^{*}\colon\QRings_{A}\stX\arr\PML_{R}(\shD,A)$.
Moreover if $\alB\in\Rings_{A}\stX$ is associative (resp. commutative,
has a unity $1\in\alB$) then $\Omega^{\alB}$ is associative (resp.
symmetric, has unity $1_{\alB}\in\Omega_{\odi{\stX}}^{\alB}$). In
particular we also get a functor $\Omega^{*}\colon\QAlg_{A}\stX\arr\ML_{R}(\shD,A)$.\end{prop}
\begin{proof}
The proof is elementary and it consists of applying the definitions.
\end{proof}
Let $\shC$ be a small monoidal subcategory of $\QCoh\stX$. Given
$\Gamma\in\PMon_{R}(\shC,A)$ with monoidal structure $\iota$, we
define the multiplication $m_{\Gamma}\colon\alA_{\Gamma,\shC}\otimes_{\odi{\stX_{A}}}\alA_{\Gamma,\shC}\arr\alA_{\Gamma,\shC}$
on $\alA_{\Gamma,\shC}=\shF_{\Gamma,\shC}$ by
\[
\alA_{\Gamma,\shC}(B)\otimes_{B\otimes_{R}A}\alA_{\Gamma,\shC}(B)\ni x_{\E,\psi}\otimes\overline{x}_{\overline{\E},\overline{\psi}}\arr\iota_{\E,\overline{\E}}(x\otimes\overline{x})_{\E\otimes\overline{\E},\psi\otimes\overline{\psi}}\in\alA_{\Gamma,\shC}(B)
\]
where $\Spec B\arr\stX$ is a map, $\E,\overline{\E}\in\shC$, $\psi\in\E(B)$,
$\overline{\psi}\in\overline{\E}(B)$, $x\in\Gamma_{\E}$, $\overline{x}\in\Gamma_{\overline{\E}}$
(see \ref{def: unit and counit for adjunction}). We continue to denote
by $\alA_{\Gamma,\shC}$ the sheaf $\shF_{\Gamma,\shC}$ together
with the multiplication map just defined. If $1\in\Gamma_{\odi{\stX}}$
is a unity we set $1_{\Gamma}\in\alA_{\Gamma,\shC}$ the image of
$1$ under the morphism $\Gamma_{\odi{\stX}}\arr\Omega_{\odi{\stX}}^{\alA_{\Gamma,\shC}}=\Hl^{0}(\alA_{\Gamma,\shC})$.
\begin{prop}
\label{prop: from sheaves to algebras}The structures defined above
yield an extension of the functor $\shF_{*,\shC}\colon\L_{R}(\shC,A)\arr\QCoh_{A}\stX$
to a functor $\alA_{*,\shC}\colon\PML_{R}(\shC,A)\arr\Rings_{A}\stX$
which is left adjoint to $\Omega^{*}\colon\Rings_{A}\stX\arr\PML_{R}(\shC,A)$.
More precisely, if $\alB\in\Rings_{A}\stX$ then the morphism $\delta_{\alB}\colon\alA_{\Omega^{\alB},\shC}\arr\alB$
preserves multiplications and unities, while if $\Gamma\in\PMon_{R}(\shC,A)$
then the natural transformation $\gamma_{\Gamma}\colon\Gamma\arr\Omega^{\alA_{\Gamma,\shC}}$
is monoidal and preserves unities (see \ref{def: unit and counit for adjunction}). 

If $\Gamma\in\PMon_{R}(\shC,A)$ is associative (resp. symmetric,
has a unity $1\in\Gamma_{\odi{\stX}}$) then $\alA_{\Gamma,\shC}$
is associative (resp. commutative, has unity $1_{\Gamma}\in\alA_{\Gamma,\shC}$).
In particular we get a functor $\alA_{*,\shC}\colon\ML_{R}(\shC,A)\arr\QAlg_{A}\stX$
which is left adjoint to $\Omega^{*}\colon\QAlg_{A}\stX\arr\ML_{R}(\shC,A)$.\end{prop}
\begin{proof}
A lot of properties have to be checked. All of them are very easy,
because they consist only in the application of the definitions. We
therefore leave them to the reader.
\end{proof}

\section{Yoneda embeddings.}

In this section we address the problem of when the Yoneda functor
$\QCoh_{A}\stX\arr\L_{R}(\shC,A)$ is fully faithful and describe
its essential image. This will led us to the notion of generating
category and left exactness for functors in $\L_{R}(\shC,A)$.

We fix an $R$-algebra $A$ and a fibered category $\stX$. We assume
that our fibered category $\stX$ is pseudo-algebraic, that is we
assume there exists a representable map $X\arr\stX$ from a scheme
which is an fpqc covering. We also fix a small subcategory $\shC$
of $\QCoh\stX$.
\begin{defn}
Let $\shD\subseteq\QCoh\stX$ be a subcategory. A sheaf $\shG\in\QCoh\stX$
is generated by $\shD$ if there exists a surjective morphism 
\[
\bigoplus_{i\in I}\E_{i}\arr\shG
\]
where $I$ is a set and $\E_{i}\in\shD$ for all $i\in I$. A sheaf
$\shG\in\QCoh_{A}\stX$ is generated by $\shD$ if it is so as an
object of $\QCoh\stX$. Equivalently, a sheaf $\shG\in\QCoh\stX_{A}$
is generated by $\shD$ if $h_{*}\shG\in\QCoh\stX$ is so, where $h\colon\stX_{A}\arr\stX$
is the projection. We define $\QCoh_{A}^{\shD}\stX$ as the subcategory
of $\QCoh_{A}\stX$ of sheaves $\shG$ generated by $\shD$ and such
that, for all maps $\E\arrdi{\psi}\shG$ with $\E\in\shD^{\oplus}$,
also $\Ker\psi$ is generated by $\shD$.

If $\shD'$ is another subcategory of $\QCoh\stX$ we will say that
$\shD$ generates $\shD'$ if all quasi-coherent sheaves in $\shD'$
are generated by $\shD$.\end{defn}
\begin{rem}
\label{rem: local characterization of generation by D} Consider a
set of morphisms $\{U_{j}=\Spec B_{j}\arr\stX\}_{j\in J}$ such that
$\sqcup_{j}U_{j}\arr\stX$ is an atlas. By \ref{prop:reducing to atlases for quasi-geometric categories}
we have the following characterizations. If $\shG\in\QCoh\stX$ then
$\shG$ is generated by $\shD$ if and only if 
\[
\forall j\in J\comma x\in\shG(B_{j})\comma\exists\E\in\shD^{\oplus}\comma\phi\in\E(B_{j})\comma\E\arrdi u\shG\text{ such that }u(\phi)=x
\]
If $\E\in\shD^{\oplus}$ and $\psi\colon\E\arr\shG$ is a map then
$\Ker\psi$ is generated by $\shD$ if and only
\[
\forall j\in J\comma y\in\E(B_{j})\text{ with }\psi(y)=0\comma\exists\overline{\E}\in\shD^{\oplus}\comma\phi\in\overline{\E}(B_{j})\comma\overline{\E}\arrdi v\E\text{ such that }\psi v=0\comma v(\phi)=y
\]
In particular if $\shG\in\QCoh^{\shD}\stX$, $\shH\in\QCoh\stX$ is
generated by $\shD$ and $\shH\arrdi{\alpha}\shG$ is a map then $\Ker\alpha$
is generated by $\shD$.\end{rem}
\begin{prop}
If $\shD$ is a subcategory of $\QCoh\stX$ then the category $\QCoh_{A}^{\shD}\stX$
is stable by direct sums. In particular $\shD\subseteq\QCoh^{\shD}\stX\iff\shD^{\oplus}\subseteq\QCoh^{\shD}\stX$.\end{prop}
\begin{proof}
Let $\shF,\shG\in\QCoh_{A}^{\shD}\stX$. Clearly $\shF\oplus\shG$
is generated by $\shD$. Now consider a map $\alpha\colon\E\arr\shF\oplus\shG$
with $\E\in\shD^{\oplus}$ and write $\alpha=\phi\oplus\psi$. By
\ref{rem: local characterization of generation by D} it follows that
$\Ker\alpha=\Ker(\phi_{|\Ker(\psi)}\colon\Ker\psi\arr\shF)$ is generated
by $\shD$ because $\Ker(\psi)$ is generated by $\shD$ and $\shF\in\QCoh_{A}^{\shD}\stX$.
\end{proof}
If $g\colon U=\Spec B\arr\stX$ is a map from a scheme then $(J_{B,\shC})^{\op}$
(see \ref{prop:FGammaCB as a direct limit}) is not a filtered category
in general. Thus if $\Gamma\in\L_{R}(\shC,A)$ and $x_{\E,\phi}\in\shF_{\Gamma,\shC}(B)$
it is very difficult to understand when $x_{\E,\phi}$ is zero in
$\shF_{\Gamma,\shC}(B)$. Luckily, under some hypothesis this is possible.
\begin{lem}
\label{lem:when J is filtered} Assume $\shC\subseteq\QCoh^{\shC}\stX$.
Then for all flat and representable maps $g\colon\Spec B\arr\stX$
the category $(J_{B,\shC^{\oplus}})^{\text{op}}$ is filtered. In
this case, given $\Gamma\in\L_{R}(\shC,A)$, an element $x_{\E,\phi}\in\shF_{\Gamma,\shC}(B)$
for $(\E,\phi)\in J_{B,\shC}$ is zero if and only if there exists
$(\overline{\E},\overline{\phi})\arrdi u(\E,\phi)$ in $J_{B,\shC}$
such that $\Gamma_{u}(x)=0$.\end{lem}
\begin{proof}
By \ref{prop:find a directed category} we can assume $\shC=\shC^{\oplus}$.
By \ref{prop:FGammaCB as a direct limit} we have to prove that for
all maps $\alpha,\beta\colon(\E,\phi)\arr(\overline{\E},\overline{\phi})$
in $J_{B,\shC}$ there exists $u\colon(\E',\phi')\arr(\E,\phi)$ in
$J_{B,\shC}$ such that $\alpha u=\beta u$. Let $\shK=\Ker((\alpha-\beta)\colon\E\arr\overline{\E})$,
so that $\phi\in\shK(B)$ since $g$ is flat. By assumption $\shK$
is generated by $\shC$ and, since $\shC$ is additive, there exist
$\E'\in\shC$, a map $u\colon\E'\arr\shK$ and $\phi'\in\E'(B)$ such
that $u(\phi')=\phi$. So $(\E',\phi')\arrdi u(\E,\phi)$ is an arrow
in $J_{B,\shC}$ such that $(\alpha-\beta)u=0$ as required. The last
claim follows from \ref{prop:FGammaCB as a direct limit} and the
fact that $J_{B,\shC}^{\op}$ is filtered.
\end{proof}
In what follows we work out sufficient (and sometimes necessary) conditions
for the surjectivity or injectivity of $\delta_{\shG}\colon\shF_{\Omega^{\shG},\shC}\arr\shG$.
Recall that $\delta_{\shG}(u_{\E,\phi})=u(\phi)$ for $\E\in\shC^{\oplus}$,
$\Spec B\arr\stX$, $\phi\in\E(B)$, $u\in\Omega_{\E}^{\shG}=\Hom_{\stX}(\E,\shG)$
(see \ref{prop:FGammaC on affine schemes}).
\begin{lem}
\label{lem:surjectivity natural map F} If $\shG\in\QCoh_{A}\stX$
the map $\delta_{\shG}\colon\shF_{\Omega^{\shG},\shC}\arr\shG$ is
surjective if and only if $\shG$ is generated by $\shC$.\end{lem}
\begin{proof}
By \ref{prop:find a directed category} we can assume $\shC=\shC^{\oplus}$.
Let $\{g_{i}\colon U_{i}=\Spec B_{i}\arr\stX\}$ be a set of maps
such that $\sqcup_{i}U_{i}\arr\stX$ is an atlas. By \ref{prop:reducing to atlases for quasi-geometric categories}
$\delta_{\shG}$ is surjective if and only if $\delta_{\shG,U_{i}}\colon\shF_{\Omega^{\shG},\shC}(U_{i})\arr\shG(U_{i})$
is surjective for all $i\in I$. By \ref{prop:FGammaCB as a direct limit}
$\Imm\delta_{\shG,U_{i}}$ is the set of elements of $\shG(U_{i})$
of the form $\delta_{\shG}(u_{\E,\phi})=u(\phi)$ for $\E\in\shC^{\oplus}$,
$\phi\in\E(U_{i})$, $u\in\Omega_{\E}^{\shG}=\Hom_{\stX}(\E,\shG)$.
So the claim follows from \ref{rem: local characterization of generation by D}.\end{proof}
\begin{lem}
\label{lem:injectivity natural map F} Let $\shG\in\QCoh_{A}\stX$.
If for all maps $\E\arrdi{\phi}\shG$ with $\E\in\shC^{\oplus}$ the
kernel $\Ker\phi$ is generated by $\shC$ then the map $\delta_{\shG}\colon\shF_{\Omega^{\shG},\shC}\arr\shG$
is injective. The converse holds if $\shC\subseteq\QCoh^{\shC}\stX$.\end{lem}
\begin{proof}
By \ref{prop:find a directed category} we can assume $\shC=\shC^{\oplus}$.
Let $\{g_{i}\colon U_{i}=\Spec B_{i}\arr\stX\}$ be a set of maps
such that $\sqcup_{i}U_{i}\arr\stX$ is an atlas. By \ref{prop:reducing to atlases for quasi-geometric categories}
$\delta_{\shG}$ is injective if and only if $\delta_{\shG,U_{i}}\colon\shF_{\Omega^{\shG},\shC}(U_{i})\arr\shG(U_{i})$
is injective for all $i\in I$. We start proving that $\delta_{\shG}$
is injective if the hypothesis in the statement is fulfilled. So let
$z\in\Ker\delta_{\shG,U_{i}}$. By \ref{prop:FGammaCB as a direct limit}
there exists $\E\in\shC$, $\phi\in\E(U_{i})$ and $u\colon\E\arr\shG$
such that $z=u_{\E,\phi}$ and $\delta_{\shG}(u_{\E,\phi})=u(\phi)=0$.
Set $\shK=\Ker u$. Since $\phi\in\shK(U_{i})$ and, by hypothesis,
$\shK$ is generated by $\shC$ there exist $\overline{\E}\in\shC$,
$\overline{\phi}\in\overline{\E}(U_{i})$ and a map $v\colon\overline{\E}\arr\shK$
such that $v(\overline{\phi})=\phi$. If we denote by $v$ also the
composition $\overline{\E}\arr\shK\arr\E$ we have
\[
u_{\E,\phi}=(\Omega_{v}^{\shG}(u))_{\overline{\E},\overline{\phi}}=(uv)_{\overline{\E},\overline{\phi}}=0_{\overline{\E},\overline{\phi}}=0\text{ in \ensuremath{\shF_{\Gamma,\shC}(U_{i})}}
\]
Assume now that $\delta_{\shG}$ is injective and $\shC\subseteq\QCoh^{\shC}\stX$
and let $u\colon\E\arr\shG$ be a map with $\E\in\shC$. We have to
prove that $\shK=\Ker u$ is generated by $\shC$. If $\phi\in\shK(U_{i})\subseteq\E(U_{i})$,
then $u(\phi)=\delta_{\shG}(u_{\E,\phi})=0$. So $u_{\E,\phi}=0$
and the conclusion follows from \ref{rem: local characterization of generation by D}
and \ref{lem:when J is filtered}.
\end{proof}
In general we can still conclude that:
\begin{prop}
\label{prop: delta E is an isomorphism} If $\E\in\shC^{\oplus}$
then the map $\delta_{\E}\colon\shF_{\Omega^{\E},\shC}\arr\E$ is
an isomorphism.\end{prop}
\begin{proof}
By \ref{prop:find a directed category} we can assume $\shC=\shC^{\oplus}$.
Let $\shH\in\QCoh\stX$. The map 
\[
\Hom_{\stX}(\E,\shH)\arrdi{\Omega_{*}}\Hom_{\L_{R}(\shC,R)}(\Omega^{\E},\Omega^{\shH})\simeq\Hom_{\stX}(\shF_{\Omega^{\E},\shC},\shH)
\]
maps $\id_{\E}$ to $\delta_{\E}$ and thus is induced by $\delta_{\E}$.
By the enriched Yoneda's lemma or a direct check we see that the above
map and therefore $\delta_{\E}$ are isomorphisms. %
\begin{comment}
\label{rem: remark of counits of adjoints}If $F\colon\catC\arr\catD$
and $G\colon\catD\arr\catC$ are adjoint functors with counit $\delta\colon FG\arr\id_{\catD}$
and unit $\gamma\colon\id_{\catC}\arr GF$ and we consider the subcategories
\[
\catC'=\{X\in\catC\st\gamma_{X}\text{ isomorphism }\}\comma\catD'=\{Y\in\catD\st\delta_{Y}\text{ isomorphism }\}
\]
of $\catC$, $\catD$ respectively, then $F_{|\catC'}\colon\catC'\arr\catD'$
and $G_{|\catD'}\colon\catD'\arr\catC'$ are well defined and quasi-inverses
of each other. This follows using that the compositions 
\[
F(X)\arrdi{F(\gamma_{X})}FGF(X)\arrdi{\delta_{F(X)}}F(X)\text{ and }G(Y)\arrdi{\gamma_{G(Y)}}GFG(Y)\arrdi{G(\delta_{Y})}G(Y)
\]
for $X\in\catC$ and $Y\in\catD$ are the identities.
\end{comment}
\end{proof}
\begin{thm}
\label{thm:when Omegastar is fully faithful} Let $\shD_{A}$ be the
subcategory of $\QCoh_{A}\stX$ of sheaves $\shG$ such that $\delta_{\shG}\colon\shF_{\Omega^{\shG},\shC}\arr\shG$
is an isomorphism. Then $\shD_{A}$ is an additive category containing
$\QCoh_{A}^{\shC}\stX$, $\shC^{\oplus}\subseteq\shD_{R}$ and the
functor
\[
\Omega^{*}\colon\shD_{A}\arr\L_{R}(\shC,A)
\]
is fully faithful. Moreover if $\shC\subseteq\QCoh^{\shC}\stX$ then
$\shD_{A}=\QCoh_{A}^{\shC}\stX$.\end{thm}
\begin{proof}
The category $\shD_{A}$ is additive because $\Omega^{*}$ and $\shF_{*,\shC}$
are additive. All the other claims follows from \ref{lem:surjectivity natural map F},
\ref{lem:injectivity natural map F}, \ref{prop: delta E is an isomorphism}
and the fact that $\delta_{\shG}$ is the counit of an adjiunction. 
\end{proof}
Now we want to address the problem of what is the essential image
of the Yoneda functor $\Omega^{*}\colon\QCoh_{A}\stX\arr\L_{R}(\shC,A)$.
We will see that if $\shF\in\QCoh_{A}\stX$ the associated Yoneda
functor $\Omega^{\shF}$ is always ``left exact'' and we will give
sufficient conditions assuring that ``left exact'' functors in $\L_{R}(\shC,A)$
are Yoneda functors associated with some quasi-coherent sheaf on $\stX$.
Since $\shC$ is not abelian, we introduce an ad hoc notion of left
exactness.
\begin{defn}
\label{def: left exact functors} Let $\shD$ be a subcategory of
$\QCoh\stX$. A \emph{test sequence }for $\shD$ is an exact sequence
\begin{equation}
\bigoplus_{k\in K}\E_{k}\arr\bigoplus_{j\in J}\E_{j}\arr\E\arr0\text{ with }\E,\E_{j},\E_{k}\in\shD\text{ for all }j\in J,k\in K\label{eq:exact sequence for functor exactness}
\end{equation}
in $\QCoh\stX$ given by maps $u_{j}\colon\E_{j}\arr\E$, $u_{kj}\colon\E_{k}\arr\E_{j}$
such that for all $k\in K$ the set $\{j\in J\st u_{kj}\neq0\}$ is
finite. We will also say that it is a test sequence for $\E\in\shD$.
A finite test sequence for $\E\in\shD$ is an exact sequence
\[
\E''\arr\E'\arr\E\arr0\text{ with }\E',\E''\in\shD^{\oplus}
\]
Given $\Gamma\in\L_{R}(\shD,A)$ we say that $\Gamma$ is exact on
the test sequence (\ref{eq:exact sequence for functor exactness})
if the sequence \begin{equation}
\begin{tikzpicture}[xscale=2.9,yscale=-0.8]     
\node (A0_2) at (2, 0) {$(x_j )_j$};     
\node (A0_3) at (3, 0) {$(\sum_j \Gamma_{u_{kj}}(x_j))_k$};     
\node (A1_0) at (0.4, 1) {$0$};     
\node (A1_1) at (1.1, 1) {$\Gamma_\E$};     
\node (A1_2) at (2, 1) {$\displaystyle{\prod_{j\in J}\Gamma_{\E_j}}$};     
\node (A1_3) at (3, 1) {$\displaystyle{\prod_{k\in K}\Gamma_{\E_k}}$};    
\node (A2_1) at (1.1, 2) {$x$};     
\node (A2_2) at (2, 2) {$(\Gamma_{u_j}(x))_j$};     
\path (A0_2) edge [|->,gray]node [auto] {$\scriptstyle{}$} (A0_3);     
\path (A1_0) edge [->]node [auto] {$\scriptstyle{}$} (A1_1);     
\path (A2_1) edge [|->,gray]node [auto] {$\scriptstyle{}$} (A2_2);     
\path (A1_2) edge [->]node [auto] {$\scriptstyle{}$} (A1_3);     
\path (A1_1) edge [->]node [auto] {$\scriptstyle{}$} (A1_2);   
\end{tikzpicture}
\label{Gamma on a test sequence}
\end{equation} is exact. We say that $\Gamma$ is left exact if it is left exact
on all short exact sequences in $\shD$. We define $\Lex_{R}(\shD,A)$
as the subcategory of $\L_{R}(\shD,A)$ of functors exact on all test
sequences in $\shD$.\end{defn}
\begin{rem}
Notice that the sequence (\ref{Gamma on a test sequence}) is a complex
because $\Gamma$ is $R$-linear. Moreover we should warn the reader
that the sequence (\ref{Gamma on a test sequence}) is not obtained
applying $\Gamma$ on the test sequence, unless $J$ and $K$ are
finite, even if $\Gamma$ is defined (or extended) to the whole $\QCoh\stX$.
The problem is that $\Gamma$ does not necessarily transforms infinite
direct sums in products.\end{rem}
\begin{prop}
\label{prop:OmegaF is left exact} Let $\shD\subseteq\QCoh\stX$ be
a subcategory. If $\shF\in\QCoh_{A}\stX$ then $\Omega^{\shF}\in\Lex_{R}(\shD,A)$.\end{prop}
\begin{proof}
It is enough to apply $\Hom_{\stX}(-,\shF)$ to the test sequence
(\ref{eq:exact sequence for functor exactness}) and observe that
\[
\Hom_{\stX}(\bigoplus_{i}\E_{i},\shF)\simeq\prod_{i}\Hom(\E_{i},\shF)\simeq\prod_{i}\Omega_{\E_{i}}^{\shF}
\]
\end{proof}
\begin{prop}
\label{prop:excatness properties of Omega* and shF*} The functor
$\Omega^{*}\colon\QCoh_{A}\stX\arr\L_{R}(\shC,A)$ is left exact.
If $\shC\subseteq\QCoh^{\shC}\stX$ then $\shF_{*,\shC}\colon\L_{R}(\shC,A)\arr\QCoh_{A}\stX$
is exact.\end{prop}
\begin{proof}
For the first claim it is enough to use that $\Hom_{\stX}(\E,-)$
is left exact. For the last part of the statement consider a set of
maps $\{U_{i}=\Spec B_{i}\arrdi{g_{i}}\stX\}_{i\in I}$ such that
$\sqcup_{i}U_{i}\arr\stX$ is an atlas. Let also $\Gamma'\arr\Gamma\arr\Gamma''$
be an exact sequence in $\L_{R}(\shC,A)$. By \ref{prop:FGammaCB as a direct limit}
the sequence $\shF_{\Gamma',\shC}(B_{i})\arr\shF_{\Gamma,\shC}(B_{i})\arr\shF_{\Gamma'',\shC}(B_{i})$
are exact for all $i\in I$ because limit of exact sequences $\Gamma_{\E,\phi}'\arr\Gamma_{\E,\phi}\arr\Gamma_{\E,\phi}''$
over the category $(J_{B_{i},\shC})^{\op}$, which is filtered thanks
to \ref{lem:when J is filtered}. Applying \ref{prop:reducing to atlases for quasi-geometric categories}
we get the result.
\end{proof}
Recall that if $\Gamma\in\L_{R}(\shC,A)$ then $\gamma_{\Gamma,\E}\colon\Gamma_{\E}\arr\Omega_{\E}^{\shF_{\Gamma,\shC}}=\Hom_{\stX}(\E,\shF_{\Gamma,\shC})$
is given by $\gamma_{\Gamma,\E}(x)(\phi)=x_{\E,\phi}$ for $\E\in\shC$,
$x\in\Gamma_{\E}$, $\Spec B\arr\stX$ and $\phi\in\E(B)$ (see \ref{prop:FGammaC on affine schemes}).
\begin{lem}
\label{lem:injectivity natural transformation Omega} Assume $\shC\subseteq\QCoh^{\shC}\stX$.
If $\Gamma\in\Lex_{R}(\shC,A)$ and the map $\Omega^{*}\circ\shF_{*,\shC}(\gamma_{\Gamma})\colon\Omega^{\shF_{\Gamma,\shC}}\arr\Omega^{\shF_{\Omega^{\shF_{\Gamma,\shC}},\shC}}$
is an isomorphism then the natural transformation $\gamma_{\Gamma}\colon\Gamma\arr\Omega^{\shF_{\Gamma,\shC}}$
is an isomorphism.\end{lem}
\begin{proof}
Let $\{U_{i}=\Spec B_{i}\arrdi{g_{i}}\stX\}_{i\in I}$ be a set of
maps such that $\sqcup_{i}U_{i}\arr\stX$ is an atlas and let $\Psi\in\L_{R}(\shC,A)$
and $x\in\Ker\gamma_{\Psi,\E}$ for some $\E\in\shC$. We are going
to prove that there exists a surjective map $\mu=\oplus_{j}\mu_{j}\colon\oplus_{j\in J}\E_{j}\arr\E$
with $\E_{j},\E\in\shC$ such that $\Psi_{u_{j}}(x)=0$ for all $j$.

If $\phi\in\E(U_{i})$, by \ref{lem:when J is filtered} and the fact
that $\gamma_{\Psi,\E}(x)(\phi)=x_{\E,\phi}$ is zero in $\shF_{\Psi,\shC}(U_{i})$,
there exists $(\E_{\phi},y_{\phi})\in J_{B_{i},\shC}$ and a map $(\E_{\phi},y_{\phi})\arrdi{u_{\phi}}(\E,\phi)$
such that $\Psi_{u_{\phi}}(x)=0$. Consider the induced map
\[
\bigoplus_{i\in I}\bigoplus_{\phi\in\E(U_{i})}\E_{\phi}\arr\E
\]
which is surjective by \ref{prop:reducing to atlases for quasi-geometric categories}.
Writing all the $\E_{\phi}\in\shC^{\oplus}$ as sums of sheaves in
$\shC$ we get the desired surjective map.

We return now to the proof of the statement. Given $x\in\Ker\gamma_{\Gamma,\E}$
and considering a surjection $\mu$ as above for $\Psi=\Gamma$, we
can conclude that $x=0$ by \ref{prop:test sequences and Extone}.
This means that the natural transformation $\gamma_{\Gamma}\colon\Gamma\arr\Omega^{\shF_{\Gamma,\shC}}$
is injective. Set now $\Pi=\Coker\gamma_{\Gamma,\shC}$. By \ref{prop:excatness properties of Omega* and shF*}
we have an exact sequence
\[
0\arr\shF_{\Gamma,\shC}\arrdi{\shF_{*,\shC}(\gamma_{\Gamma})}\shF_{\Omega^{\shF_{\Gamma,\shC}}}\arr\shF_{\Pi,\shC}\arr0
\]
This is a split sequence because the composition of $\shF_{*,\shC}(\gamma_{\Gamma})\colon\shF_{\Gamma,\shC}\arr\shF_{\Omega^{\shF_{\Gamma,\shC}},\shC}$
and $\delta_{\shF_{\Gamma,\shC}}\colon\shF_{\Omega^{\shF_{\Gamma,\shC}},\shC}\arr\shF_{\Gamma,\shC}$
is the identity. So $\Omega^{*}$ maintains the exactness of the above
sequence and therefore 
\[
\Omega^{\shF_{\Pi,\shC}}=\Coker(\Omega^{*}\circ\shF_{*,\shC}(\gamma_{\Gamma}))=0
\]
We want to prove that $\Pi=0$. Let $x\in\Pi_{\E}$ for $\E\in\shC$.
Since $\Omega^{\shF_{\Pi,\shC}}=0$ we have $x\in\Ker\gamma_{\Pi,\E}$.
Consider a surjection $\mu=\oplus_{j}\mu_{j}$ constructed as above
starting from $x\in\Pi_{\E}$ and $\Psi=\Pi$. By \ref{rem: local characterization of generation by D}
$\mu$ can be extended to a test sequence $\oplus_{k}\E_{k}\arr\oplus_{j}\E_{i}\arrdi{\mu}\E$
because $\shC\subseteq\QCoh^{\shC}\stX$. Since $\Omega^{\shF_{\Gamma,\shC}}\in\Lex_{R}(\shC,A)$
by \ref{prop:OmegaF is left exact} we get a commutative diagram   \[   \begin{tikzpicture}[xscale=2.0,yscale=-1.4]     
\node (A0_1) at (1, 0) {$0$};     
\node (A0_2) at (2, 0) {$0$};     
\node (A0_3) at (3, 0) {$0$};     
\node (A1_0) at (0, 1) {$0$};     
\node (A1_1) at (1, 1) {$\Gamma_\E$};     
\node (A1_2) at (2, 1) {$\Omega^{\shF_{\Gamma,\shC}}_\E$};     
\node (A1_3) at (3, 1) {$\Pi_\E$};     
\node (A1_4) at (4, 1) {$0$};     
\node (A2_0) at (0, 2) {$0$};     
\node (A2_1) at (1, 2) {$\displaystyle\prod_{j\in J}\Gamma_{\E_j}$};     
\node (A2_2) at (2, 2) {$\displaystyle\prod_{j\in J}\Omega^{\shF_{\Gamma,\shC}}_{\E_j}$};     
\node (A2_3) at (3, 2) {$\displaystyle\prod_{j\in J}\Pi_{\E_j}$};     
\node (A2_4) at (4, 2) {$0$};     \node (A3_0) at (0, 3) {$0$};     
\node (A3_1) at (1, 3) {$\displaystyle\prod_{k\in K}\Gamma_{\E_k}$};     
\node (A3_2) at (2, 3) {$\displaystyle\prod_{k\in K}\Omega^{\shF_{\Gamma,\shC}}_{\E_k}$};    
%\node (A3_3) at (3, 3) {$\displaystyle\prod_{k\in K}\Delta_{\E_k}$};     
%\node (A3_4) at (4, 3) {$0$};     
\path (A2_1) edge [->]node [auto] {$\scriptstyle{}$} (A3_1);    
\path (A2_2) edge [->]node [auto] {$\scriptstyle{}$} (A2_3);     
\path (A1_2) edge [->]node [auto] {$\scriptstyle{}$} (A1_3);     
\path (A0_3) edge [->]node [auto] {$\scriptstyle{}$} (A1_3);     
\path (A2_1) edge [->]node [auto] {$\scriptstyle{}$} (A2_2);     
\path (A1_0) edge [->]node [auto] {$\scriptstyle{}$} (A1_1);     
\path (A3_0) edge [->]node [auto] {$\scriptstyle{}$} (A3_1);     
\path (A1_3) edge [->]node [auto] {$\scriptstyle{\beta}$} (A2_3);     
\path (A1_1) edge [->]node [auto] {$\scriptstyle{}$} (A1_2);     
\path (A2_3) edge [->]node [auto] {$\scriptstyle{}$} (A2_4);     
%\path (A2_3) edge [->]node [auto] {$\scriptstyle{}$} (A3_3);     
\path (A2_2) edge [->]node [auto] {$\scriptstyle{}$} (A3_2);     
%\path (A3_2) edge [->]node [auto] {$\scriptstyle{}$} (A3_3);     
\path (A1_3) edge [->]node [auto] {$\scriptstyle{}$} (A1_4);     
\path (A1_1) edge [->]node [auto] {$\scriptstyle{}$} (A2_1);     
\path (A0_2) edge [->]node [auto] {$\scriptstyle{}$} (A1_2);     
%\path (A3_3) edge [->]node [auto] {$\scriptstyle{}$} (A3_4);     
\path (A0_1) edge [->]node [auto] {$\scriptstyle{}$} (A1_1);     
\path (A2_0) edge [->]node [auto] {$\scriptstyle{}$} (A2_1);     
\path (A3_1) edge [->]node [auto] {$\scriptstyle{}$} (A3_2);     
\path (A1_2) edge [->]node [auto] {$\scriptstyle{}$} (A2_2);   
\end{tikzpicture}   \]  in which all the rows and the first two columns are exact. By diagram
chasing it is easy to conclude that $\beta$ is injective. Since by
construction $\beta(x)=0$ we can conclude that $x=0$.\end{proof}
\begin{defn}
We define $\Lex_{R}^{\shC}(A)$ as the subcategory of $\Lex_{R}(\shC,A)$
of functors $\Gamma$ such that $\shF_{\Gamma,\shC}\in\QCoh_{A}^{\shC}\stX$.\end{defn}
\begin{thm}
\label{thm:relative theorem for equivalence} Assume $\shC\subseteq\QCoh^{\shC}\stX$.
Then the functors 
\[
\Omega^{*}\colon\QCoh_{A}^{\shC}\stX\arr\Lex_{R}^{\shC}(A)\text{ and }\shF_{*,\shC}\colon\Lex_{R}^{\shC}(A)\arr\QCoh_{A}^{\shC}\stX
\]
are quasi-inverses of each other. \end{thm}
\begin{proof}
Let $\Gamma\in\L_{R}(\shC,A)$ be such that $\shF_{\Gamma,\shC}\in\QCoh_{A}^{\shC}\stX$.
The composition
\[
\shF_{\Gamma,\shC}\arrdi{\shF_{*,\shC}(\gamma_{\Gamma})}\shF_{\Omega^{\shF_{\Gamma,\shC}},\shC}\arrdi{\delta_{\shF_{\Gamma,\shC}}}\shF_{\Gamma,\shC}
\]
is the identity and $\delta_{\shF_{\Gamma,\shC}}$ is an isomorphism
since $\shF_{\Gamma,\shC}\in\QCoh_{A}^{\shC}\stX$ by \ref{thm:when Omegastar is fully faithful}.
Thus $\shF_{*,\shC}(\gamma_{\Gamma})$ and therefore $\Omega^{*}\circ\shF_{*,\shC}(\gamma_{\Gamma})$
are isomorphisms. By \ref{lem:injectivity natural transformation Omega}
we can conclude that if $\Gamma\in\Lex_{R}^{\shC}(A)$ then $\gamma_{\Gamma}\colon\Gamma\arr\Omega^{\shF_{\Gamma,\shC}}$
is an isomorphism. The result then follows from \ref{thm:when Omegastar is fully faithful}
and \ref{prop:OmegaF is left exact}.
\end{proof}
The following result allow us to extend results from small subcategories
of $\QCoh\stX$ to any subcategory.
\begin{prop}
\label{prop: QCoh X is Grothendieck} The category $\QCoh\stX$ is
generated by a small subcategory. Equivalently $\QCoh\stX$ has a
generator, that is there exists $\E\in\QCoh\stX$ such that $\{\E\}$
generates $\QCoh\stX$. %
\begin{comment}
\begin{proof}
We want to describe a quasi-coherent sheaf as descent datum. Denote
by $\pi\colon\Sch/\stX\arr\Sch/R$ the structure map and let $U_{i}=\Spec B_{i}\arr\stX$
be representable and flat morphisms such that $\sqcup_{i}U_{i}\arr\stX$
is an fpqc covering. Consider the set of maps $Z$ obtained iterating
the following procedure starting from $\{U_{i}\arr\stX\}$. Given
a set $H\subseteq\Sch/\stX$ of representable and flat maps let $H'\subseteq\Sch/\stX$
the union of $H$ and of the following sets for all $U\arr\stX,V\arr\stX\in H$:
pick an etal� covering $\{\Spec C_{i}\arr U\times_{\stX}V\}_{i\in I}$
and add the maps $\{\Spec C_{i}\arr\stX\}$, which are still representable
and flat. Now let $\stZ\subseteq\Sch/\stX$ be the subcategory with
$Z$ as set of objects and with arrows the morphisms $z\arr z'$ such
that $\pi(z)\arr\pi(z')$ is flat. Thus $\stZ$ is a small category
of flat and representable maps from an affine scheme containing the
$U_{i}\arr\stX$ and such that, for all $U\arr\stX$ and $V\arr\stX$
there exists an etal� covering of $\{W_{j}\arr U\times_{\stX}V\}$
such that the projections $W_{j}\arr U$ and $W_{j}\arr V$ are in
$\stZ$. Given $z\in\stZ$ set $B_{z}=\odi{\stX}(z)$. Let $D(\stX)$
be the category whose objects are tuples $(\underline{M},\underline{\xi})$
where $\underline{M}=(M_{z})_{z}\in\prod_{z\in Z}\Mod B_{z}$ and
$\underline{\xi}=(\xi_{\phi})_{\phi\in\Arr(\stZ)}$ is a collection
of $B_{\overline{z}}$-linear maps $\xi_{\phi}\colon M_{\overline{z}}\arr M_{z}$
such that $M_{\overline{z}}\otimes_{B_{\overline{z}}}B_{z}\arr M_{z}$
is an isomorphism for all $\phi\in\Hom_{\stZ}(z,\overline{z})$, $z,\overline{z}\in\stZ$,
$\underline{\xi}$ is compatible with compositions and $\xi_{\id_{z}}=\id_{M_{z}}$
for all $z\in\stZ$. The arrows are the obvious ones. Then the association
$\QCoh\stX\arr D(\stX)$ which sends a sheaf $\shF\colon(\Sch/\stX)^{\text{op}}\simeq\stX^{\text{op}}\arr\Mod R$
to the object $((F(z))_{z\in\stZ},(F(\phi))_{\phi\in\Arr(\stZ)})$
is an equivalence of categories. With abuse of notation we will call
the objects of $D(\stX)$ quasi-coherent sheaves. 

Given $\underline{N}\in\prod_{z\in\stZ}\Mod B_{z}$ we will use the
notation $\underline{N}=(N_{z})_{z\in\stZ}$. Moreover if $\underline{N},\underline{N}'\in\prod_{z\in\stZ}\Mod B_{z}$
we will write $\underline{N}<\underline{N}'$ if $N_{z}$ is a sub
$B_{z}$-module of $N_{z}'$ for all $z\in\stZ$. If $N$ is a $B$-module
for some ring $B$ and $\aleph$ is a cardinal we write $g(N)\leq\aleph$
meaning that $N$ has a system of generators of cardinality less or
equal than $\aleph$. Similarly we write $g(\underline{N})\leq\aleph$
if $g(N_{z})\leq\aleph$ for all $z\in\stZ$.

Let $\aleph$ be an infinite cardinal such that $ $$|\Arr(\stZ)|\leq\aleph$
and set 
\[
\shC=\{(\underline{H},\underline{\eta})\in D(\stX)\st g(\underline{H})\leq\aleph\}
\]
Our aim is to prove that $\shC$ generates $\QCoh\stX$ and that it
is small. We start with the last part. Given $\underline{H}\in\prod_{z\in\stZ}\Mod B_{z}$
the collection of $\underline{\eta}$ such that $(\underline{H},\underline{\eta})\in\shC$
forms a set. Thus it is enough to prove that $\shC'=\{\underline{H}\in\prod_{z\in\stZ}\Mod B_{z}\st g(\underline{H})\leq\aleph\}$
is a set. Writing $\shC'$ as a product, we reduce the problem to
show that $\shC_{B}=\{H\in\Mod B\st g(\underline{H})\leq\aleph\}$
is a set for all ring $B$. This follows because the set of submodules
of $B^{\aleph}$ has a map onto $\shC_{B}$.

It remains to prove that $\shC$ generates $\QCoh\stX$. Let $(\underline{M},\underline{\xi})\in D(\stX)$
and $\underline{N}<\underline{M}$. We want to define a new $J(\underline{N})<\underline{M}$.
We set $J(\underline{N})=\underline{N}'$ where
\[
N'_{z}=\langle\xi_{\phi}(N_{\overline{z}})\text{ for all z\ensuremath{\arrdi{\phi}}}\overline{z}\in\stZ\rangle_{B_{z}}
\]
Note that $\underline{N}<J(\underline{N})$. Set $J^{n}(\underline{N})=J(J^{n-1}(\underline{N}))$
and $J^{\infty}(\underline{N})=\cup_{n\in\N}J^{n}(\underline{N})$,
where the union is taken over each component. If $\overline{\aleph}$
is an infinite cardinal such that $g(\underline{N})\leq\overline{\aleph}$
and $|\Arr(\stZ)|\leq\overline{\aleph}$ it follows that $g(J(\underline{N}))\leq\overline{\aleph}$
and $g(J^{\infty}(\underline{N}))\leq\overline{\aleph}$. Note that
quasi-coherent ``subsheaves'' correspond to $\underline{H}<\underline{M}$
such that $\langle\xi_{\phi}(H_{\overline{z}})\rangle_{B_{z}}=H_{z}$
for all $\phi\in\Hom_{\stZ}(z,\overline{z})$, $z,\overline{z}\in\stZ$.
It is easy to check that $J^{\infty}(\underline{N})$ corresponds
to a sub quasi-coherent sheaf of $(\underline{M},\underline{\xi})\in D(\stX)$.
Let now $\shF\in\QCoh\stX$ and $(\underline{M},\underline{\xi})$
be the corresponding object in $D(\stX)$. If $z_{0}\in\stZ$, $m\in\shF(z_{0})=M_{z_{0}}$
and we consider $\underline{N}<\underline{M}$ given by $N_{z_{0}}=\langle m\rangle_{B_{z_{0}}}$,
$N_{z}=0$ if $z\neq z_{0}$ the object $(J^{\infty}(\underline{N}),\underline{\xi})\in D(\stX)$
belongs to $\shC$ and it corresponds to a subsheaf $\shG\subseteq\shF$
such that $m\in N_{z_{0}}<J^{\infty}(\underline{N})_{z_{0}}=\shG(z_{0})$.
Thus all the sections of $\shF$ on some $z\in\stZ$ are contained
in a subsheaf $\shG\in\shF$ such that $\shG\in\shC$. This implies
that $\shC$ generates $\shF$.\end{proof}
\end{comment}
\end{prop}
\begin{proof}
Follows from \ref{prop:reducing to atlases for quasi-geometric categories}
and \cite[Tag 0780]{SP014}.\end{proof}
\begin{rem}
\label{rem: finding small subcategory generating QCoh X} If $\shD\subseteq\QCoh\stX$
generates $\QCoh\stX$ there always exists a small subcategory $\overline{\shD}\subseteq\shD$
that generates $\QCoh\stX$. Indeed if $\E$ is a generator of $\QCoh\stX$
it is enough to take a subset of sheaves in $\shD$ that generates
$\E$.\end{rem}
\begin{thm}
\label{thm:general theorem for equivalences} Let $\shD\subseteq\QCoh\stX$
be a subcategory that generates $\QCoh\stX$. Then the functor
\[
\Omega^{*}\colon\QCoh_{A}\stX\arr\Lex_{R}(\shD,A)
\]
is an equivalence of categories and, if $\shD$ is small, $\shF_{*,\shD}\colon\Lex_{R}(\shD,A)\arr\QCoh_{A}\stX$
is a quasi-inverse. In particular if $\overline{\shD}\subseteq\shD$
is a subcategory that generates $\QCoh\stX$ the restriction functor
$\Lex_{R}(\shD,A)\arr\Lex_{R}(\overline{\shD},A)$ is an equivalence.\end{thm}
\begin{proof}
If $\shD$ is small we have $\QCoh_{A}\stX=\QCoh_{A}^{\shD}\stX$,
$\Lex_{R}(\shD,A)=\Lex_{R}^{\shD}(A)$ and everything follows from
\ref{thm:relative theorem for equivalence}. In particular the restriction
$\Lex_{R}(\shD,A)\arr\Lex_{R}(\overline{\shD},A)$ is an equivalence
if $\overline{\shD}\subseteq\shD$, they are small and they generate
$\QCoh\stX$. Assume now that $\shD$ is not necessarily small and
consider a small subcategory $\shC\subseteq\shD$ that generates $\QCoh\stX$,
which exists thanks to \ref{rem: finding small subcategory generating QCoh X}.
The proof of the statement is complete if we prove that the restriction
functor $\Lex_{R}(\shD,A)\arr\Lex_{R}(\shC,A)$ is an equivalence.
Given a set $I\subseteq\shD$ we set $\shC_{I}=\shC\cup I\subseteq\shD$.
For all sets $I$, quasi-coherent sheaves are generated by $\shC_{I}$.
Note that we have the restriction functor $-_{|C_{I}}\colon\Lex_{R}(\shD,A)\arr\Lex_{R}(\shC_{I},A)$
and the composition $\QCoh_{A}\stX\arrdi{\Omega^{*}}\Lex_{R}(\shD,A)\arrdi{-_{|\shC_{I}}}\Lex_{R}(C_{I},A)$
is an equivalence for all $I$. In particular $-_{|C}\colon\Lex_{R}(\shD,A)\arr\Lex_{R}(\shC,A)$
is essentially surjective. We will conclude by proving that it is
fully faithful. Let $\Gamma,\Gamma'\in\Lex_{R}(\shD,A)$. If $\overline{I}\subseteq I$
the restriction functor $\Lex_{R}(\shC_{I},A)\arr\Lex_{R}(\shC_{\overline{I}},A)$
is an equivalence of categories. In particular the map $\Hom(\Gamma_{|\shC_{I}},\Gamma'_{|\shC_{I}})\arr\Hom(\Gamma_{|\shC_{\overline{I}}},\Gamma'_{|\shC_{\overline{I}}})$
is bijective. Using this, it is elementary to prove that also the
map $\Hom(\Gamma,\Gamma')\arr\Hom(\Gamma_{|\shC},\Gamma'_{|\shC})$
is bijective.
\end{proof}
As a corollary we recover Gabriel-Popescu's theorem for the category
$\QCoh\stX$.
\begin{thm}
\label{thm: Gabriel Popescu's theorem}{[}Gabriel-Popescu's theorem{]}
If $\E$ is a generator of $\QCoh\stX$ then the functor
\[
\Hom_{\stX}(\E,-)\colon\QCoh\stX\arr\Mod_{\text{right}}(\End_{\stX}(\E))
\]
is fully faithful and has an exact left adjoint.\end{thm}
\begin{proof}
It follows from \ref{thm:general theorem for equivalences} and \ref{prop:excatness properties of Omega* and shF*}
with $\shD=\shC=\{\E\}$.
\end{proof}
Before showing other applications of \ref{thm:general theorem for equivalences}
we want to present a cohomological interpretation of the functors
in $\Lex_{R}(\shC,A)$, which will allow us to show that it is often
enough to consider just finite test sequences instead of arbitrary
test sequences.
\begin{rem}
In an abelian category $\alA$, given $X,Y\in\alA$ we can always
define the abelian group $\Ext^{1}(X,Y)$ as the group of extensions
(regardless if $\alA$ has enough injectives) and it has the usual
nice properties on short exact sequences. In order to avoid set-theoretic
problems one should require that $\alA$ is locally small and that,
given $X,Y\in\alA$, $\Ext^{1}(X,Y)$ is a set. This is the case for
$\alA=\L_{R}(\shC,R)$, for instance by looking at the cardinalities
of the $\Gamma_{\E}$ for $\Gamma\in\L_{R}(\shC,R)$ and $\E\in\shC$.\end{rem}
\begin{defn}
\label{def: cohomologically left exact} Given a surjective map $\mu=\oplus_{j}\mu_{j}\colon\oplus_{j}\E_{j}\arr\E$
with $\E,\E_{j}\in\shC$ we set $\Omega^{\mu}=\oplus_{j}\Omega^{\mu_{j}}\colon\oplus_{j}\Omega^{\E_{j}}\arr\Omega^{\E}$.
A functor $\Gamma\in\L_{R}(\shC,A)$ is\emph{ cohomologically left
exact }on $\mu$ if
\begin{equation}
\Hom_{\L_{R}(\shC,R)}(\Omega^{\E}/\Imm(\Omega^{\mu}),\Gamma)=\Ext_{\L_{R}(\shC,R)}^{1}(\Omega^{\E}/\Imm(\Omega^{\mu}),\Gamma)=0\label{eq:cohomological vanishing left exact functors}
\end{equation}
It is cohomologically left exact if it is so on all surjections $\mu$
as above.
\end{defn}
We setup some notation. We denote by $\Phi_{\shC}(\E)$ for $\E\in\shC$
the set of subfunctor of $\Omega^{\E}$of the form $\Imm(\Omega^{\mu})$
for some surjective map $\mu\colon\oplus_{j}\E_{i}\arr\E$ with $\E_{j}\in\shC$
and by $\Phi_{\shC}$ the disjoint union of all the $\Phi_{\shC}(\E)$.
Given $\Delta=\Imm(\Omega^{\mu})\in\Phi_{\shC}$ and $\Gamma\in\L_{R}(\shC,A)$
we will say that $\Gamma$ is cohomologically left exact on $\Delta$
if it is cohomologically left exact on $\mu$. Notice that, using
Yoneda's lemma, we have an $A$-linear isomorphism
\[
\Hom_{\L_{R}(\shC,R)}(\oplus_{j}\Omega^{\E_{j}},\Gamma)\simeq\prod_{j}\Gamma_{\E_{j}}\text{ for all }\E_{j}\in\shC\comma\Gamma\in\L_{R}(\shC,A)
\]
In particular $\Hom_{\L_{R}(\shC,R)}(\Omega^{\E},-)$, which is the
evaluation in $\E$, is exact, which implies that $\Ext_{\L_{R}(\shC,R)}^{1}(\Omega^{\E},-)=0$. 

A map $\oplus_{j}\E_{j}\arr\oplus_{k}\E_{k}$ is called \emph{locally
finite }if, for all $j$, the restriction $\E_{j}\arr\oplus_{k}\E_{k}$
factors through a finite sub-sum. Composition of locally finite maps
is locally finite and, using Yoneda's lemma, we obtain a functorial
map 
\[
\Hom_{\L_{R}(\shC,R)}(\oplus_{j}\Omega^{\E_{j}},\oplus_{k}\Omega^{\E_{k}})\arr\Hom(\oplus_{j}\E_{j},\oplus_{k}\E_{k})
\]
which is an isomorphism onto the set of locally finite maps.

Let $\mu\colon\oplus_{j}\E_{j}\arr\E$ be a surjective map and set
$\Delta=\Imm(\Omega^{\mu})$. Notice that $\Delta_{\E'}$ is the set
of maps $\E'\arr\E$ which factors through $\mu$ via a locally finite
map $\E'\arr\oplus_{j}\E_{j}$. Given a map $u\colon\overline{\E}\arr\E$
in $\shC$ we set $u^{-1}(\Delta)=\Delta\times_{\Omega^{\E}}\Omega^{\overline{\E}}\subseteq\Omega^{\overline{\E}}$:
$u^{-1}(\Delta)_{\E'}$ is the set of maps $\E'\arr\overline{\E}$
such that $\E'\arr\overline{\E}\arr\E$ factors through $\mu$ via
a locally finite map $\E'\arr\oplus_{j}\E_{j}$. Notice that if $\shC\subseteq\QCoh^{\shC}\stX$
then $u^{-1}(\Delta)\in\Phi_{\shC}(\overline{\E})$. Indeed, if $\shH$
is the kernel of $\mu\oplus(-u)\colon(\oplus_{j}\E_{j})\oplus\overline{\E}\arr\E$
then all elements of $\shH$ on some object $\Spec B\arr\stX$ lies
in a finite subsum of $(\oplus_{j}\E_{j})\oplus\overline{\E}$. By
\ref{rem: local characterization of generation by D} $\shH$ is the
image of a locally finite map $\oplus_{q}\E_{q}\arr(\oplus_{j}\E_{j})\oplus\overline{\E}$
and, since $\shH\arr\overline{\E}$ is surjective, the induced map
$\mu'\colon\oplus_{\E'}\oplus_{\phi\in u^{-1}(\Delta)_{\E'}}\E'\arr\overline{\E}$
is surjective and clearly $u^{-1}(\Delta)=\Imm(\Omega^{\mu'})$.
\begin{lem}
\label{lem:almost left exact functors and hom} Let $\mu\colon\oplus_{j}\E_{i}\arr\E$
be a surjective map with $\E_{j},\E\in\shC$ and $\Gamma\in\L_{R}(\shC,A)$.
Then there is an exact sequence of $A$-modules
\[
0\arr\Hom_{\L_{R}(\shC,R)}(\Omega^{\E}/\Imm(\Omega^{\mu}),\Gamma)\arr\Gamma_{\E}\arr\prod_{j}\Gamma_{\E_{j}}
\]
If $\shT:\oplus_{k}\E_{k}\arr\oplus_{j}\E_{j}\arrdi{\mu}\E\arr0$
is a test sequence and $\Gamma$ is exact on $\shT$ then $\Gamma$
is cohomologically left exact on $\mu$. The converse holds if the
map 
\[
\Hom_{\L_{R}(\shC,R)}(\Ker(\Omega^{\mu}),\Gamma)\arr\prod_{k}\Gamma_{\E_{k}}
\]
obtained applying $\Hom_{\L_{R}(\shC,R)}(-,\Gamma)$ to the map $\oplus_{k}\Omega^{\E_{k}}\arr\Ker(\Omega^{\mu})$
is injective.\end{lem}
\begin{proof}
Set $\Delta=\Imm(\Omega^{\mu})$, $K=\Ker(\Omega^{\mu})$. Consider
the diagram   \[   \begin{tikzpicture}[xscale=2.7,yscale=-1.1]     
\node (A0_0) at (0.47, 0) {$\Hom(\Omega^\E/\Delta,\Gamma)$};     
\node (A0_1) at (1.3, 0) {$\Gamma_\E$};     
\node (A0_2) at (2, 0) {$\Hom(\Delta,\Gamma)$};     
\node (A1_2) at (2, 1) {$\prod_j \Gamma_{\E_j}$};     
\node (A2_2) at (2, 2) {$\Hom(K,\Gamma)$};     
\node (A0_3) at (3, 0) {$\Ext^1(\Omega^\E/\Delta,\Gamma)$};     
\node (A0_4) at (4, 0) {$\Ext^1(\Omega^\E,\Gamma)$};     
\node (A2_3) at (2.8, 2) {$\prod_k \Gamma_{\E_k}$};     
\path (A0_1) edge [->]node [auto] {$\scriptstyle{\alpha}$} (A1_2);     
\path (A1_2) edge [->]node [auto] {$\scriptstyle{\beta}$} (A2_3);     \path (A0_0) edge [right hook->]node [auto] {$\scriptstyle{}$} (A0_1);     \path (A0_1) edge [->]node [auto] {$\scriptstyle{}$} (A0_2);     \path (A0_2) edge [right hook->]node [auto] {$\scriptstyle{}$} (A1_2);
\path (A2_2) edge [->]node [auto] {$\scriptstyle{\lambda}$} (A2_3);     
\path (A0_3)edge [->]node [auto] {$\scriptstyle{}$} (A0_4);     \path (A0_2) edge [->]node [auto] {$\scriptstyle{}$} (A0_3);     \path (A1_2) edge [->]node [auto] {$\scriptstyle{}$} (A2_2);   \end{tikzpicture}   \]  The convention here is that $\beta$ and $\lambda$ are defined only
when a test sequence $\shT$ as in the statement exists and we will
not use them for the first statement. All the other maps are obtained
splitting $\Omega^{\E_{j}}\arr\Omega^{\E}\arr0$ into two exact sequences
and applying $\Hom(-,\Gamma)$, so that the first line and the central
column are exact. The map $\alpha$ obtained as composition is the
map defined in the first sequence in the statement. In particular
the first claim follows. So let's focus on the second one. The map
$\lambda$ is the second map in the statement while the map $\beta$
together with $\alpha$ are the maps defining the sequence (\ref{Gamma on a test sequence}).
Since $\Ext^{1}(\Omega^{\E},\Gamma)=0$ also the second claim follows.\end{proof}
\begin{lem}
\label{lem:key lemma for reducing test sequences} Let $\Gamma,K\in\L_{R}(\shC,R)$
and $u\colon\oplus_{q}\Omega^{\E_{q}}\arr K$ be a map, where $\E_{q}\in\shC$.
If for all $\Omega^{\E}\arr K$ with $\E\in\shC$ there exists a surjective
map $v\colon\oplus_{t}\E_{t}\arr\E$ with $\E_{t}\in\shC$ such that
the composition $\oplus_{t}\Omega^{\E_{t}}\arr\Omega^{\E}\arr K$
factors thorugh $u$ and $\Gamma$ is cohomologically left exact on
$v$ then the map 
\[
\Hom_{\L_{R}(\shC,R)}(K,\Gamma)\arr\Hom_{\L_{R}(\shC,R)}(\oplus_{q}\Omega^{\E_{q}},\Gamma)\simeq\prod_{q}\Gamma_{\E_{q}}
\]
is injective.\end{lem}
\begin{proof}
Let $\E\in\shC$ and $x\in K_{\E}$, which corresponds to a map $\Omega^{\E}\arr K$.
Consider the data given by hypothesis with respect to this last map.
We have commutative diagrams   \[   \begin{tikzpicture}[xscale=2.5,yscale=-1.2]     
\node (A0_0) at (0, 0) {$\oplus_t \Omega^{\E_t}$};     
\node (A0_1) at (1, 0) {$\Omega^{ \E}$};     
\node (A0_2) at (2, 0) {$\Hom(K,\Gamma)$};     
\node (A0_3) at (3, 0) {$\Gamma_{ \E}$};     
\node (A1_0) at (0, 1) {$\oplus_q \Omega^{\E_q}$};     
\node (A1_1) at (1, 1) {$K$};     
\node (A1_2) at (2, 1) {$\prod_q \Gamma_{\E_q}$};     
\node (A1_3) at (3, 1) {$\prod_t \Gamma_{\E_t}$};     
\path (A0_0) edge [->]node [auto] {$\scriptstyle{\Omega^v}$} (A0_1);     
\path (A0_1) edge [->]node [auto] {$\scriptstyle{x}$} (A1_1);     
\path (A1_0) edge [->]node [auto] {$\scriptstyle{u}$} (A1_1);     
\path (A0_3) edge [->]node [auto] {$\scriptstyle{\delta}$} (A1_3);     
\path (A0_2) edge [->]node [auto] {$\scriptstyle{\lambda}$} (A1_2);     
\path (A0_0) edge [->]node [auto] {$\scriptstyle{}$} (A1_0);     
\path (A0_2) edge [->]node [auto] {$\scriptstyle{\gamma}$} (A0_3);     
\path (A1_2) edge [->]node [auto] {$\scriptstyle{}$} (A1_3);   
\end{tikzpicture}   \]  where the second diagram is obtained by applying $\Hom(-,\Gamma)$
to the first one. The map $\lambda$ is the map in the statement,
while $\gamma$ is the evalutation in $x\in K_{\E}$. Thanks to \ref{lem:almost left exact functors and hom}
and since $\Gamma$ is cohomologically left exact on $v$ the map
$\delta$ is injective. So if $\phi\in\Hom(K,\Gamma)$ is such that
$\lambda(\phi)=0$ it follows that $\gamma(\phi)=\phi_{\E}(x)=0$,
as required.\end{proof}
\begin{thm}
\label{prop:test sequences and Extone} If $\shC\subseteq\QCoh^{\shC}\stX$
then $\Lex_{R}(\shC,A)$ coincides with the subcategory of $\L_{R}(\shC,A)$
of cohomologically left exact functors.\end{thm}
\begin{proof}
Let $\mu=\oplus_{j}\mu_{j}\colon\oplus_{j\in J}\E_{j}\arr\E$ be a
surjective map with $\E,\E_{j}\in\shC$ and set $\Delta=\Imm(\Omega^{\mu})$,
$K=\Ker(\Omega^{\mu})$. By \ref{rem: local characterization of generation by D}
there exists a test sequence $\oplus_{k}\E_{k}\arr\oplus_{j}\E_{j}\arrdi{\mu}\E\arr0$.
Using \ref{lem:almost left exact functors and hom} we have to prove
that if $\Gamma\in\L_{R}(\shC,R)$ is cohomologically left exact then
$\lambda\colon\Hom(K,\Gamma)\arr\prod_{k}\Gamma_{\E_{k}}$ is injective.
We are going to apply \ref{lem:key lemma for reducing test sequences}
with respect to the map $\oplus_{k}\Omega^{\E_{k}}\arr K$. If $\overline{\E}\in\shC$,
a map $\Omega^{\overline{\E}}\arr K$ is a locally finite map $\overline{\E}\arr\oplus_{j}\E_{j}$
which is zero composed by $\mu$, or, equivalently, mapping in the
image of $\oplus_{k}\E_{k}\arr\oplus_{j}\E_{j}$. Consider the kernel
$\shH$ of the difference of the maps $\oplus_{k}\E_{k}\arr\oplus_{j}\E_{j}$
and $\overline{\E}\arr\oplus_{j}\E_{j}$. Since this difference map
is locally finite, $\shC\subseteq\QCoh^{\shC}\stX$ and using \ref{rem: local characterization of generation by D}
there is a surjective map $\oplus_{t}\E_{t}\arr\shH$ with $\E_{t}\in\shC$
such that $\oplus_{t}\E_{t}\arr\oplus_{k}\E_{k}$ is locally finite
and $\oplus_{t}\E_{t}\arr\overline{\E}$ is surjective. This gives
the desired factorization for applying \ref{lem:key lemma for reducing test sequences}.
\end{proof}
We now show how to reduce the number of test sequences in order to
check when a $\Gamma\in\L_{R}(\shC,A)$ belongs to $\Lex_{R}(\shC,A)$.
The following is the key lemma:
\begin{lem}
\label{lem: reducing the set of test sequences}Let $\Phi'\subseteq\Phi_{\shC}$
such that, for all $\E\in\shC$ and $\Delta\in\Phi_{\shC}(\E)$ there
exists $\Delta'\in\Phi'\cap\Phi_{\shC}(\E)$ such that $\Delta'\subseteq\Delta$
(inside $\Omega^{\E}$). If $\Gamma\in\L_{R}(\shC,A)$, $\shC\subseteq\QCoh^{\shC}\stX$
and $\Gamma$ is cohomologically left exact on all the elements of
$\Phi'$ then $\Gamma$ is cohomologically left exact.\end{lem}
\begin{proof}
Consider $\Delta\in\Phi_{\shC}(\E)$, $\Delta'\subseteq\Delta$ with
$\Delta'\in\Phi'$ and the exact sequence $0\arr\Delta/\Delta'\arr\Omega^{\E}/\Delta'\arr\Omega^{\E}/\Delta\arr0$.
Applying $\Hom(-,\Gamma)$ and using that $\Ext^{1}(\Omega^{\E},\Gamma)=0$,
the only non trivial vanishing to check is $\Hom(\Delta/\Delta',\Gamma)=0$.
Write $\Delta'=\Imm(\Omega^{u})$, where $u\colon\oplus_{q}\E_{q}\arr\E$.
We can apply \ref{lem:key lemma for reducing test sequences} to $K=\Delta$
and the map $\Omega^{u}\colon\oplus_{q}\Omega^{\E_{q}}\arr\Delta$:
if $\Omega^{\E'}\arr\Delta\subseteq\Omega^{\E}$ is a map corresponding
to $\psi\colon\E'\arr\E$, then $\psi^{-1}(\Delta')\in\Phi_{\shC}$
and, by hypothesis, we can find $\Phi'\ni\Imm(\Omega^{v})\subseteq\psi^{-1}(\Delta')$;
the last inclusion tells us that $v$ is the factorization required
for \ref{lem:key lemma for reducing test sequences}. Thus the map
$\Hom(\Omega^{u},\Gamma)\colon\Hom(\Delta,\Gamma)\arr\Hom(\oplus_{q}\Omega^{\E_{q}},\Gamma)$
is injective and, since $\Delta'=\Imm(\Omega^{u})\subseteq\Delta$,
this is also true for the restriction $\Hom(\Delta,\Gamma)\arr\Hom(\Delta',\Gamma)$,
whose kernel is $\Hom(\Delta/\Delta',\Gamma)$, as required. \end{proof}
\begin{prop}
\label{ prop: Lex and left exact functors} Let $\shD\subseteq\QCoh\stX$
be a subcategory and $\Gamma\in\L_{R}(\shD,A)$. If $\Gamma\in\Lex_{R}(\shD,A)$
then $\Gamma$ is exact on finite test sequences in $\shD$ and transforms
arbitrary direct sums in $\shD$ into products. The converse holds
if one of the following conditions is satisfied:
\begin{itemize}
\item the category $\shD$ is stable by arbitrary direct sums;
\item all the sheaves in $\shD$ are finitely presented, $\shD\subseteq\QCoh^{\shD}\stX$
and $\stX$ is quasi-compact. In this case $\Gamma\in\Lex_{R}(\shD,A)$
if and only if it is cohomologically left exact on all surjective
maps $\E'\arr\E$ with $\E\in\shD$ and $\E'\in\shD^{\oplus}$.
\end{itemize}
In any of the above cases, if moreover $\shD$ is additive and all
surjections in $\shD$ have kernel in $\shD$ then $\Lex_{R}(\shD,A)$
is the subcategory of $\L_{R}(\shD,A)$ of left exact functors which
transforms arbitrary direct sums in products.\end{prop}
\begin{proof}
If $\Gamma\in\Lex_{R}(\shD,A)$ then it is clearly exact on finite
test sequences. Given a set $\{\E_{j}\in\shD\}_{j\in J}$ set $\E=\bigoplus_{j}\E_{j}$.
If $\E\in\shD$, then the sequence
\[
0\arr\bigoplus_{j\in J}\E_{j}\arrdi{\id_{\E}}\E\arr0
\]
is a test sequence and therefore we get that the natural map $\Gamma_{\E}\arr\prod_{j}\Gamma_{\E_{j}}$
is an isomorphism. If $\shD$ is stable by arbitrary direct sums is
easy to see that the converse holds. Moreover the last part of the
statement follows easily from the first part. 

So we focus on the second point and we assume that all the sheaves
in $\shD$ are finitely presented, $\shD\subseteq\QCoh^{\shD}\stX$
and that $\stX$ is quasi-compact. Since the class of finitely presented
quasi-coherent sheaves on $\stX$ modulo isomorphism is a set, we
can assume $\shD=\shC$ small. Let $\Phi'\subseteq\Phi_{\shC}$ be
the subset of functors of the form $\Imm(\Omega^{\mu})$ for some
surjective map $\mu\colon\E'\arr\E$ with $\E\in\shC$ and $\E'\in\shC^{\oplus}$.
The set $\Phi'$ satisfies the hypothesis of \ref{lem: reducing the set of test sequences}:
if $v\colon\oplus_{j\in J}\E_{j}\arr\E$ is a surjective map then
there exists a finite subset $J_{0}\subseteq J$ such that $v_{|\E'}\colon\E'=\oplus_{j\in J_{0}}\E_{j}\arr\E$
is surjective because $\E$ is of finite type and $\stX$ is quasi-compact.
In particular, taking into account \ref{prop:test sequences and Extone},
the last claim of the second point follows. It remains to show that
if $\Gamma\in\L_{R}(\shC,A)$ is exact on finite test sequences then
$\Gamma$ is cohomologically left exact on all the elements of $\Phi'$.
Let $\mu\colon\E'\arr\E$ be a surjective map with $\E\in\shC$ and
$\E'\in\shC^{\oplus}$. Since $\E$ is finitely presented and $\E'$
is of finite type it follows that $\Ker(\mu)$ is of finite type and,
since $\stX$ is quasi-compact and $\shC\subseteq\QCoh^{\shC}\stX$,
there exists $\E''\in\shC^{\oplus}$ and a surjective map $\E''\arr\Ker(\mu)$.
Thus $\E''\arr\E'\arrdi{\mu}\E\arr0$ is a finite test sequence and
by \ref{lem:almost left exact functors and hom} it follows that $\Gamma$
is cohomologically left exact on $\mu$ as required. %
\begin{comment}
\begin{proof}
Notice that a map $\E\arr\bigoplus_{s\in S}\E_{s}$ with $\E,\E_{s}\in\shC^{\oplus}$
always factors through a finite sum since $\stX$ is quasi-compact
and $\E$ is of finite type. In particular given a map $\bigoplus_{t\in T}\E_{t}\arr\bigoplus_{s\in S}\E_{s}$
with $\E_{s},\E_{t}\in\shC^{\oplus}$ induced by $\mu_{ts}\colon\E_{t}\arr\E_{s}$
we can always define the map
\[
\prod_{s\in S}\Gamma_{\E_{s}}\arr\prod_{t\in T}\Gamma_{\E_{t}}\comma(x_{s})_{s\in S}\longmapsto(\sum_{s\in S}\Gamma_{\mu_{ts}}(x_{s}))_{t\in T}
\]
Moreover this association is functorial on compositions of such maps. 

We start proving that given a commutative diagram \begin{equation}\label{eq: reducing to finite test sequences}
  \begin{tikzpicture}[xscale=2.3,yscale=-1.3]     \node (A0_0) at (0, 0) {$\overline \E$};     \node (A0_1) at (1, 0) {$\E_{t_0}$};     \node (A0_2) at (2, 0) {$\shG$};     \node (A0_3) at (3, 0) {$0$};     \node (A1_0) at (0, 1) {$\displaystyle \bigoplus_{s\in S}\E_s$};     \node (A1_1) at (1, 1) {$\displaystyle \bigoplus_{t\in T}\E_t$};     \node (A1_2) at (2, 1) {$\shG$};     \node (A1_3) at (3, 1) {$0$};     \path (A1_2) edge [->]node [auto] {$\scriptstyle{}$} (A1_3);     \path (A0_0) edge [->,dashed]node [auto] {$\scriptstyle{}$} (A0_1);     \path (A0_1) edge [->]node [auto] {$\scriptstyle{\beta}$} (A1_1);     \path (A1_0) edge [->]node [auto] {$\scriptstyle{\alpha}$} (A1_1);     \path (A1_1) edge [->]node [auto] {$\scriptstyle{}$} (A1_2);     \path (A0_2) edge [->]node [auto] {$\scriptstyle{\id_\shG}$} (A1_2);     \path (A0_0) edge [->,dashed]node [auto] {$\scriptstyle{}$} (A1_0);     \path (A0_1) edge [->]node [auto] {$\scriptstyle{\gamma}$} (A0_2);     \path (A0_2) edge [->]node [auto] {$\scriptstyle{}$} (A0_3);   \end{tikzpicture}   
\end{equation} with $t_{0}\in T$, $\E_{s},\E_{t}\in\shC^{\oplus}$ for all $s\in S$,
$t\in T$, $\shG\in\QCoh\stX$ finitely presented and with exact rows
there exists $\overline{\E}\in\shC^{\oplus}$ and dashed maps such
that the first row is exact. Indeed consider $\shK=\Ker(\alpha\oplus(-\beta))$.
Since $\shC\subseteq\QCoh^{\shC}\stX$, it is easy to see that $\shK$
is generated by $\shC$. So there exists $\{\E_{q}\in\shC\}_{q\in Q}$
and a surjective map $\bigoplus_{q}\E_{q}\arr\shK$. Note that the
image of the map $\shK\arr\E_{t_{0}}$ is exactly $\Ker\gamma$, which
is of finite type because $\E_{t_{0}}$ and $\shG$ are finitely presented.
In particular, since $\stX$ is quasi-compact, we can find a finite
sum $\overline{\E}\subseteq\bigoplus_{q}\E_{q}$ such that the induced
map $\overline{\E}\arr\Ker\gamma$ is surjective.

Let now $\Gamma\in\L_{R}(\shC,A)$ such that it is exact over finite
test sequences in $\shC$. Consider an exact sequence as in (\ref{eq:exact sequence for functor exactness}),
but with $\E_{j},\E_{k}\in\shC^{\oplus}$ instead of only in $\shC$.
We want to prove that the sequence (\ref{Gamma on a test sequence})
is exact. Since $\E$ is finitely presented and $\stX$ quasi-compact
there exists a finite sum $\E'\subseteq\bigoplus_{j}\E_{j}$ such
that the induced map $\E'\arr\E$ is surjective. Without loss of generality
we can assume that $\E'=\E_{j_{0}}$ with $j_{0}\in J$. Consider
the diagram (\ref{eq: reducing to finite test sequences}) with $S=K,$
$T=J$, $t_{0}=j_{0}$ and $\shG=\E$. In particular the sequence
\[
0\arr\Gamma_{\E}\arrdi{\Gamma_{\gamma}}\Gamma_{\E_{j_{0}}}\arr\Gamma_{\overline{\E}}
\]
is exact. This implies that the map $\xi\colon\Gamma_{\E}\arr\prod_{j}\Gamma_{\E_{j}}$
is injective. Now consider $x=(x_{j})_{j}\in\prod_{j}\Gamma_{\E_{j}}$
in the kernel of $\prod_{j}\Gamma_{\E_{j}}\arr\prod_{k}\Gamma_{\E_{k}}$.
It follows that $x_{j_{0}}\in\Ker(\Gamma_{\E_{j_{0}}}\arr\Gamma_{\overline{\E}})$
and therefore there exists $y\in\Gamma_{\E}$ such that $\Gamma_{\gamma}(y)=x_{j_{0}}$.
We can replace $x$ with $x-\xi(y)$ and therefore assume $x_{j_{0}}=0$.
Since $\E_{j_{0}}\arr\E$ is surjective, it follows that the map $\bigoplus_{k}\E_{k}\arr\bigoplus_{j\neq j_{0}}\E_{j}$
is also surjective. Moreover by construction $x\in\Ker(\prod_{j\neq j_{0}}\Gamma_{\E_{j}}\arr\prod_{k}\Gamma_{\E_{k}})$.
We want to show that $x=0$ and thus it is enough to prove that, if
$\bigoplus_{s}\E_{s}\arr\bigoplus_{t}\E_{t}$ is a surjective map,
with $\E_{s},\E_{t}\in\shC^{\oplus}$ then the map $\prod_{t}\Gamma_{\E_{t}}\arr\prod_{s}\Gamma_{\E_{s}}$
is injective. Here $S$ and $T$ are again any pair of sets. Moreover,
without loss of generality, we can assume that $\E_{s},\E_{t}\in\shC$.
Consider an element $z=(z_{t})_{t}\in\prod_{t}\Gamma_{\E_{t}}$ in
the kernel and $t_{0}\in T$. Again we consider the diagram (\ref{eq: reducing to finite test sequences})
with $\shG=0$, obtaining a sheaf $\overline{\E}\in\shC^{\oplus}$
and dashed maps. Since $\E_{t_{0}}\in\shC$ and the map $\overline{\E}\arr\E_{t_{0}}$
is surjective and can be extended to a finite test sequence in $\shC$,
the map $\Gamma_{\E_{t_{0}}}\arr\Gamma_{\overline{\E}}$ is injective.
But $z_{t_{0}}\in\Ker(\Gamma_{\E_{t_{0}}}\arr\Gamma_{\overline{\E}})$
by construction and we can conclude that $z_{t_{0}}=0$ and $z=0$.\end{proof}
\end{comment}

\end{proof}
There is another characterization of $\Lex_{R}(\shC,A)$ in terms
of sheaves on a site. Although we will not use it in this paper, I
think it is worth to point out. We refer to \cite[Tag 00YW]{SP014}
for general definitions and properties. We start by comparing $\Lex_{R}(\shC,A)$
and $\Lex_{R}(\shC^{\oplus},A)$.
\begin{prop}
\label{prop:Lex for Coplus and C} If $\shC\subseteq\QCoh^{\shC}\stX$
then the equivalence $\L_{R}(\shC^{\oplus},A)\simeq\L_{R}(\shC,A)$
maps $\Lex_{R}(\shC^{\oplus},A)$ to $\Lex_{R}(\shC,A)$.\end{prop}
\begin{proof}
We can assume $A=R$. Let $\Gamma\in\L_{R}(\shC^{\oplus},R)$ such
that $\Gamma\in\Lex_{R}(\shC,R)$ and consider $\Phi'\subseteq\Phi_{\shC^{\oplus}}$
the set of subfunctors $\Delta\subseteq\Omega^{\E}$ that can be written
as follows: $\E=\E_{1}\oplus\cdots\oplus\E_{r}$ and there are surjective
maps $\mu_{k}\colon\oplus_{q}\E_{q,k}\arr\E_{k}$ for $\E_{k},\E_{q,k}\in\shC$
such that $\Delta=\Imm(\Omega^{\mu_{1}})\oplus\cdots\oplus\Imm(\Omega^{\mu_{r}})$.
Since for such $\Delta$ we have 
\[
\Omega^{\E}/\Delta\simeq\oplus_{i}(\Omega^{\E_{i}}/\Delta_{i})
\]
 it follows that $\Gamma$ is cohomologically left exact on all the
elements of $\Phi'$. Taking into account \ref{prop:test sequences and Extone},
in order to conclude that $\Gamma\in\Lex_{R}(\shC,R)$ we can show
that $\Phi'\subseteq\Phi_{\shC^{\oplus}}$ satisfies the hypothesis
of \ref{lem: reducing the set of test sequences}. So let $\Delta=\Imm(\Omega^{\mu})\in\Phi_{\shC^{\oplus}}$
where $\mu\colon\oplus_{q}\E_{q}\arr\E$ where $\E,\E_{q}\in\shC^{\oplus}$.
If $\E=\E'_{1}\oplus\cdots\oplus\E'_{r}$ with $\E'_{i}\in\shC$ and
$\psi_{i}\colon\E'_{i}\arr\E$ are the inclusions then $\Delta_{i}=\psi_{i}^{-1}(\Delta)\in\Phi_{\shC^{\oplus}}(\E'_{i})=\Phi_{\shC}(\E'_{i})$
and it is easy to see that $\Phi'\ni\Delta_{1}\oplus\cdots\oplus\Delta_{r}\subseteq\Delta$
as required.\end{proof}
\begin{prop}
\label{prop:Lex and Grothendieck topology} If $\shC\subseteq\QCoh^{\shC}\stX$
and $\shJ$ is the smallest Grothendieck topology on $\shC^{\oplus}$
containing $\Phi_{\shC^{\oplus}}$ then $\Lex_{R}(\shC^{\oplus},A)$
is the category of sheaves of $A$-modules on $(\shC^{\oplus},\shJ)$
which are $R$-linear.\end{prop}
\begin{proof}
We can assume $A=R$ and $\shC=\shC^{\oplus}$. If $\Delta\subseteq\Omega^{\E}$
is a sieve and $f\colon\E'\arr\E$ we set $f^{-1}(\Delta)=\Delta\times_{\Omega^{\E}}\Omega^{\E'}\subseteq\Omega^{\E'}$.
Let $\widetilde{\shJ}$ be the set of sieves $\Delta\subseteq\Omega^{\E}$
of $\shC$ such that, for all $\Gamma\in\Lex_{R}(\shC,R)$ and maps
$f\colon\E'\arr\E$ the map
\[
\Hom_{\sets}(\Omega^{\E'},\Gamma)\arr\Hom_{\sets}(f^{-1}(\Delta),\Gamma)
\]
is bijective. Here $\Hom_{\sets}$ denotes the set of natural transformation
of functors with values in $\sets$. The set $\widetilde{\shJ}$ is
a Grothendieck topology on $\shC$ such that all functors in $\Lex_{R}(\shC,R)$
are sheaves. Notice that, by \ref{rem: natural transf are linear for additive categories},
if $A,B\in\L_{R}(\shC,R)$ then $\Hom_{\sets}(A,B)=\Hom_{\L_{R}(\shC,R)}(A,B)$.
Moreover if $\Delta\in\Phi_{C}$ and $f\colon\E'\arr\E$ is a map
in $\shC$ then $\Delta'=f^{-1}(\Delta)\in\Phi_{\shC}$ and if $\Gamma\in\L_{R}(\shC,R)$
then, applying $\Hom_{\L_{R}(\shC,R)}(-,\Gamma)$ on the exact sequence
$0\arr\Delta'\arr\Omega^{\E'}\arr\Omega^{\E'}/\Delta'\arr0$ and taking
into account that $\Ext^{1}(\Omega^{\E'},\Gamma)=0$ we obtain an
exact sequence 
\[
0\arr\Hom(\Omega^{\E'}/\Delta',\Gamma)\arr\Hom_{\sets}(\Omega^{\E'},\Gamma)\arr\Hom_{\sets}(\Delta',\Gamma)\arr\Ext^{1}(\Omega^{\E'}/\Delta',\Gamma)\arr0
\]
Thus, if $\shJ$ is the smallest topology on $\shC$ containing $\Phi_{\shC}$
then $\Phi_{\shC}\subseteq\shJ\subseteq\widetilde{\shJ}$ and everything
follows easily from \ref{prop:test sequences and Extone}.
\end{proof}
We now apply \ref{thm:general theorem for equivalences} and \ref{ prop: Lex and left exact functors}
in some (more) concrete situations.
\begin{thm}
\label{thm:equivalence for QCoh X} The category $\Lex_{R}(\QCoh\stX,A)$
is the category of contravariant, $R$-linear and left exact functors
$\Gamma\colon\QCoh\stX\arr\Mod A$ which transform arbitrary direct
sums in products. Moreover the functor
\[
\Omega^{*}\colon\QCoh_{A}\stX\arr\Lex_{R}(\QCoh\stX,A)
\]
is an equivalence of categories.\end{thm}
\begin{proof}
Follows from \ref{thm:general theorem for equivalences} and \ref{ prop: Lex and left exact functors}
with $\shD=\QCoh\stX$.\end{proof}
\begin{thm}
\label{thm:equivalence for noetherian algebraic stacks} Let $\stX$
be a noetherian algebraic stack. The category $\Lex_{R}(\Coh\stX,A)$
is the category of contravariant, $R$-linear and left exact functors
$\Coh\stX\arr\Mod A$. Moreover the functor
\[
\Omega^{*}\colon\QCoh_{A}\stX\arr\Lex_{R}(\Coh\stX,A)
\]
is an equivalence of categories.\end{thm}
\begin{proof}
Follows from \ref{thm:general theorem for equivalences} and \ref{ prop: Lex and left exact functors},
taking into account that in our assumptions $\Coh\stX$ is an abelian
category that generates $\QCoh\stX$.\end{proof}
\begin{thm}
\label{thm:equivalence for Loc X} Assume that $\stX$ is quasi-compact
and that $\Loc\stX$ generates $\QCoh\stX$. Then $\Lex_{R}(\Loc\stX,A)$
is the category of contravariant, $R$-linear and left exact functors
$\Loc\stX\arr\Mod A$. Moreover the functor
\[
\Omega^{*}\colon\QCoh_{A}\stX\arr\Lex_{R}(\Loc\stX,A)
\]
is an equivalence of categories.\end{thm}
\begin{proof}
Follows from \ref{thm:general theorem for equivalences} and \ref{ prop: Lex and left exact functors},
taking into account that all surjections in $\Loc\stX$ have kernels
in $\Loc\stX$.\end{proof}
\begin{thm}
\label{thm:equivalence for affine schemes} Let $B$ be an $R$-algebra
and $\shD\subseteq\Mod B$ be a subcategory that generates $\Mod B$,
that is there exists $\E_{1},\dots,\E_{r}\in\shD$ with a surjective
map $\bigoplus_{i}\E_{i}\arr B$. Then the functor
\[
\Omega^{*}\colon\Mod(A\otimes_{R}B)\arr\Lex_{R}(\shD,A)
\]
is an equivalence of categories. Moreover if $\shD\subseteq\Loc B$
then $\Lex_{R}(\shD,A)=\L_{R}(\shD,A)$.\end{thm}
\begin{proof}
If $\stX=\Spec B$, then $\QCoh_{A}\stX\simeq\Mod(A\otimes_{R}B)$
and the first part follows from \ref{thm:general theorem for equivalences}.
For the last claim, observe that any $\Gamma\colon\Loc B\arr\Mod A$
is exact because any short exact sequence in $\Loc B$ splits. By
\ref{thm:equivalence for Loc X} we can conclude that $\L_{R}(\Loc B,A)=\Lex_{R}(\Loc B,A)$.
If now $\shD\subseteq\Loc B$ and $\Gamma\in\L_{R}(\shD,A)$, we can
extend it to $\overline{\Gamma}\in\L_{R}(\Loc B,A)$ and therefore
$\Gamma=\overline{\Gamma}_{|\shD}\in\Lex_{R}(\shD,A)$.
\end{proof}
We want to extend Theorem \ref{thm:general theorem for equivalences}
to functors with monoidal structures.
\begin{defn}
If $\shD$ is a monoidal subcategory of $\QCoh\stX$ we define $\PMLex_{R}(\shD,A)$
(resp. $\MLex_{R}(\shD,A)$) as the subcategory of $\PML_{R}(\shD,A)$
(resp. $\ML_{R}(\shD,A)$) of functors $\Gamma$ such that $\Gamma\in\Lex_{R}(\shD,A)$.\end{defn}
\begin{thm}
\label{thm:fundamental for sheaf of algebras and monoidal functors}
Let $\shD$ be a monoidal subcategory of $\QCoh\stX$ that generates
it. Then the functors
\[
\Omega^{*}\colon\Rings_{A}\stX\arr\PMLex_{R}(\shD,A)\text{ and }\Omega^{*}\colon\QAlg_{A}\stX\arr\MLex_{R}(\shD,A)
\]
(see \ref{prop:Omegastar on algebras}) are equivalence of categories.
If $\shD$ is small a quasi inverse is given by $\alA_{*,\shD}\colon\PMLex_{R}(\shD,A)\arr\Rings_{A}\stX$
and $\alA_{*,\shD}\colon\MLex_{R}(\shD,A)\arr\QAlg_{A}\stX$ respectively
(see \ref{prop: from sheaves to algebras}). Moreover if $\overline{\shD}\subseteq\shD$
is a monoidal subcategory that generates $\QCoh\stX$ the restriction
functors $\PMLex_{R}(\shD,A)\arr\PMLex_{R}(\overline{\shD},A)$ and
$\MLex_{R}(\shD,A)\arr\MLex_{R}(\overline{\shD},A)$ are equivalences.\end{thm}
\begin{proof}
Assume that $\shD$ is small. Then $\Omega^{*}\colon\Rings_{A}\stX\arr\PMLex_{R}(\shD,A)$
and $\alA_{*,\shD}\colon\PMLex_{R}(\shD,A)\arr\Rings_{A}\stX$ are
quasi-inverses of each other because, by \ref{prop: from sheaves to algebras},
we have natural transformations $\id\arrdi{\gamma}\Omega^{*}\circ\alA_{*,\shD}$
and $\alA_{*,\shD}\circ\Omega^{*}\arrdi{\beta}\id$ which are isomorphisms
thanks to \ref{thm:general theorem for equivalences}. In particular
if $\overline{\shD}\subseteq\shD$ is a monoidal subcategory that
generates $\QCoh\stX$ then the restriction functor $\PMLex_{R}(\shD,A)\arr\PMLex_{R}(\overline{\shD},A)$
is an equivalence. Since $\gamma$ and $\beta$ preserve unities by
\ref{prop: from sheaves to algebras}, the same holds if we replace
$\Rings_{A}\stX$ by $\QAlg_{A}\stX$ and $\PMLex_{R}(\shD,A)$ by
$\MLex_{R}(\shD,A)$. Notice that there exists a small subcategory
$\shC'\subseteq\shD$ that generates $\QCoh\stX$ thanks to \ref{rem: finding small subcategory generating QCoh X}.
If $I\subseteq\shD$ is a set we set $\shC_{I}$ for the category
containing all tensor products with factors in $C\cup I$ and $\odi{\stX}$.
We have that $\shC_{I}$ is a collection of small monoidal subcategories
of $\shD$, that generate $\QCoh\stX$ and such that $\shC_{I}\subseteq\shC_{I'}$
if $I\subseteq I'$. If $\shC=\shC_{\emptyset}$ we can show that
the restrictions $\PMLex_{R}(\shD,A)\arr\PMLex_{R}(\shC,A)$ and $\MLex_{R}(\shD,A)\arr\MLex_{R}(\shC,A)$
are equivalences by proceeding as in the proof of \ref{thm:general theorem for equivalences}.
All the other claims in the statement follow easily from this fact.\end{proof}
\begin{thm}
\label{thm:rewriting of general equivalences for monoidal functors}
The theorems \ref{thm:equivalence for QCoh X}, \ref{thm:equivalence for noetherian algebraic stacks},
\ref{thm:equivalence for Loc X} continue to hold if we replace $\Lex_{R}$
by $\PMLex_{R}$ (resp. $\MLex_{R}$), $\QCoh_{A}\stX$ by $\QRings_{A}\stX$
(resp. $\QAlg_{A}\stX$) and the word ``functors'' by ``pseudo-monoidal
functors'' (resp. ``monoidal functors'').
\end{thm}

\section{\label{sec: group schemes representation} Group schemes and representations.}

Let $G$ be a flat and affine group scheme over $R$. In this section
we want to interpret the results obtained in the case $\stX=\Bi_{R}G$,
the stack of $G$-torsors for the \emph{fpqc} topology, which is a
quasi-compact fpqc stack with affine diagonal.

If $A$ is an $R$-algebra, by standard theory we have that $\QCoh_{A}\Bi_{R}G$
is the category $\Mod^{G}A$ of $G$-comodules over $A$. Recall that
the regular representation $R[G]$ of $G$ is by definition the $G$-comodule
$p_{*}\odi G$. By definition it comes equipped with a morphism of
$R$-algebras $\varepsilon\colon R[G]\arr R$ induced by the unit
section of $G$.
\begin{rem}
\label{rem: crucial isomorphism invariants RG and modules} If $M\in\Mod^{G}A$
then the composition
\[
(M\otimes_{R}R[G])^{G}\arr M\otimes_{R}R[G]\arrdi{\id_{M}\otimes\varepsilon}M
\]
 is an isomorphism. This follows from \cite[3.4]{Jantzen2003} applied
to $G=H$.
\end{rem}
We start with a criterion to find a set of generators for $\QCoh\Bi_{R}G$.
\begin{prop}
\label{thm:generating sheaves for BG} If the regular representation
$R[G]$ is a filtered direct limit of modules $B_{i}\in\Mod^{G}R$
which are finitely presented as $R$-modules then $\{\duale B_{i}\}_{i\in I}$
generates $\Mod^{G}R$.\end{prop}
\begin{proof}
Set $B=R[G]$ and $\varepsilon_{i}\colon B_{i}\arr R$ for the composition
$B_{i}\arr B\arrdi{\varepsilon}R$ and let $M\in\Mod^{G}R$. Since
filtered direct limits commute with tensor products and taking invariants,
by \ref{rem: crucial isomorphism invariants RG and modules} we have
that the limit of the maps $(\varepsilon_{i}\otimes\id_{M})_{|(B_{i}\otimes M)^{G}}\colon(B_{j}\otimes M)^{G}\arr M$
is an isomorphism. This means that for any $m\in M$ there exists
$i_{m}\in I$ and an element $\psi_{m}\in(B_{i_{m}}\otimes M)^{G}$
such that $(\varepsilon_{i}\otimes\id_{M})(\psi_{m})=m$. The map
$B_{i}\otimes M\arr\Hom(\duale{B_{i}},M)$ is $G$-equivariant and
therefore we obtain a $\delta_{m}\in\Hom^{G}(\duale{B_{i}},M)$ such
that $\delta_{m}(\varepsilon_{i})=m$. This implies that the map
\[
\bigoplus_{m\in M}\delta_{m}\colon\bigoplus_{m\in M}\duale B_{i_{m}}\arr M
\]
is surjective and therefore that $M$ is generated by $\{\duale{B_{i}}\}_{i\in I}$.\end{proof}
\begin{rem}
\label{rem: resolution property for BG cases} The class $\shG_{R}$
of flat, affine group schemes $G$ over $R$ such that $R[G]$ is
a direct limit of modules in $\Loc(\Bi_{R}G)$ is stable by arbitrary
products and projective limits. Moreover by construction contains
all groups which are flat, finite and finitely presented over $R$,
i.e. $R[G]\in\Loc(\Bi_{R}G)$, and thus all profinite groups. Since
any $G$-comodule is the union of the sub $G$-comodules which are
finitely generated $R$-modules (see \cite[2.13]{Jantzen2003}), we
see that $\shG_{R}$ contains all flat groups defined over a Dedekind
domain or a field, such as $\GL_{r}$, $\text{SL}_{r}$ and all diagonalizable
groups. Proposition \ref{thm:generating sheaves for BG} tells us
that if $G\in\shG_{R}$ then $\Bi_{R}G$ has the resolution property.
\end{rem}
Let $A$ be an $R$-algebra. We denote by $\Loc^{G}A$ the subcategory
of $\Mod^{G}A$ of $G$-comodules that are locally free of finite
rank (projective of finite type) as $A$-modules, so that $\Loc(\Bi_{R}G)\simeq\Loc^{G}R$.
We define $\QAdd^{G}A$ ($\QPMon^{G}A$, $\QMon^{G}A$) as the category
of \emph{covariant }$R$-linear (pseudo-monoidal, monoidal) functors
$\Loc^{G}R\arr\Mod A$. We set $\QRings^{G}A$ for the category of
$M\in\Mod^{G}A$ with a $G$-equivariant map $M\otimes_{A}M\arr M$
and $\QAlg^{G}A$ for the (not full) subcategory of $\QRings^{G}A$
of commutative $R$-algebras.
\begin{defn}
The group $G$ is called \emph{linearly reductive} if the functor
$(-)^{G}\colon\Mod^{G}R\arr\Mod R$ is exact.\end{defn}
\begin{rem}
\label{rem: LocG R semisimple if G linearly reductive} If $G$ is
linearly reductive then any short exact sequence in $\Loc^{G}R$ splits.
Indeed if $M\arr N$ is surjective then $\Hom_{R}^{G}(N,M)\arr\Hom_{R}^{G}(N,N)$
is surjective, yielding a $G$-equivariant section $N\arr M$.\end{rem}
\begin{thm}
\label{thm:general theorem for BG} If $\Bi_{R}G$ has the resolution
property then the functors
\[
\Mod^{G}A\arr\QAdd^{G}A\comma\QRings^{G}A\arr\QPMon^{G}A\comma\QAlg^{G}A\arr\QMon^{G}A
\]
which maps $M$ to the functor $(-\otimes_{R}M)^{G}\colon\Loc^{G}R\arr\Mod A$
are well defined, fully faithful and have essential image the subcategory
of functors which are left exact on short exact sequences in $\Loc^{G}R$.
In particular they are equivalences if $G$ is a linearly reductive
group.\end{thm}
\begin{proof}
Set $\shC=\Loc^{G}R$. The functor $\duale{(-)}\colon\Loc^{G}R\arr\Loc^{G}R$
is an equivalence and therefore we get equivalences $\QAdd^{G}A\simeq\L_{R}(\shC,A)$,
$\QPMon^{G}A\simeq\PML_{R}(\shC,A)$ and $\QMon^{G}A\simeq\ML_{R}(\shC,A)$.
Left exact functors are sent to left exact functors. Under those equivalences
$\Omega^{M}$ corresponds to $(-\otimes_{R}M)^{G}$ because $\Hom_{\Bi_{R}G}(\duale{\E},M)\simeq\Hl^{0}(\E\otimes_{\odi{\Bi_{R}G}}M)\simeq(\E\otimes_{R}M)^{G}$.
Thus the result follows from \ref{prop:excatness properties of Omega* and shF*},
\ref{thm:equivalence for Loc X}, \ref{thm:rewriting of general equivalences for monoidal functors}
and \ref{rem: LocG R semisimple if G linearly reductive}.
\end{proof}

\section{Tannaka reconstruction for stacks with the resolution property.}
\begin{defn}
A (contravariant) monoidal functor $\Omega\colon\catC\arr\catD$ between
symmetric monoidal categories is called \emph{strong }if the maps
\[
\Omega_{V}\otimes\Omega_{W}\arr\Omega_{V\otimes W}
\]
are isomorphisms for all $V,W\in\catC$ and the map $J\arr F_{I}$
is an isomorphism, where $I$ and $J$ are the unities of $\catC$
and $\catD$ respectively. 
\end{defn}
In this section we want to understand what sheaves of algebras correspond
to strong monoidal functors in the equivalence of \ref{thm:fundamental for sheaf of algebras and monoidal functors},
in the case where $\shD$ is a subcategory of locally free sheaves. 

We will consider only fpqc stacks with quasi-affine diagonal, for
instance algebraic stacks with quasi-affine diagonal (see \cite[Corollary 10.7]{Laumon1999})
and quasi-separated schemes. This is because resolution property is
somehow meaningless for other stacks, see for instance Remark $(1)$
in the introduction of \cite{Totaro2002}.
\begin{defn}
If $\stX$ is a fiber category over $R$, $\shC\subseteq\Loc\stX$
is a monoidal subcategory and $A$ is an $R$-algebra we define $\SMex_{R}(\shC,A)$
as the subcategory of $\MLex_{R}(\shC,A)$ of functors $\Gamma$ which
are strong monoidal and, for all geometric points $\Spec k\arr\Spec A$,
$\Gamma\otimes_{A}k\in\MLex_{R}(\shC,k)$. Given a fiber category
$\stY$ we denote by $\text{Fib}_{\stX,\shC}(\stY)$ the category
of \emph{covariant}, $R$-linear and strong monoidal functors $\Gamma\colon\shC\arr\Loc\stY$
which are exact on all exact sequences $\E''\arr\E'\arr\E\arr0$ with
$\E\in\shC$ and $\E',\E''\in\shC^{\oplus}$. We also define $\text{Fib}_{\stX,\shC}$
as the fiber category over $R$ whose fiber over an $R$-algebra $A$
is $\text{Fib}_{\stX,\shC}(\Spec A)$ and we call $\shP_{\shC}$ the
functor
\begin{equation}
\shP_{\shC}\colon\stX\arr\text{Fib}_{\stX,\shC}\comma(\Spec A\arrdi s\stX)\longmapsto(s^{*}\colon\shC\arr\Loc A)\label{eq:tannaka reconstruction functor}
\end{equation}

\end{defn}
We will prove the following:
\begin{thm}
\label{thm: Tannaka reconstrion for Loc X} Let $\stX$ be a quasi-compact
fpqc stack over $R$ with quasi-affine diagonal, $A$ be an $R$-algebra
and $\shC\subseteq\Loc\stX$ be a monoidal subcategory with duals
that generates $\QCoh\stX$. Then the functors   \[   \begin{tikzpicture}[xscale=4.2,yscale=-0.6]     \node (A0_0) at (0, 0) {$(\Spec \alA_{\Gamma,\shC }\arr \stX)$};     \node (A0_1) at (1, 0) {$(\Gamma \colon \shC\arr\Mod A)$};     \node (A1_0) at (0, 1) {$\stX(A)$};     \node (A1_1) at (1, 1) {$\SMex_R(\shC,A)$};     \node (A2_0) at (0, 2) {$(s\colon \Spec A\arr\stX)$};     \node (A2_1) at (1, 2) {$(\duale {(s^*_{|\shC})}\colon \shC\arr\Loc A)$};     \path (A1_0) edge [->]node [auto] {$\scriptstyle{}$} (A1_1);     \path (A2_0) edge [|->,gray]node [auto] {$\scriptstyle{}$} (A2_1);     \path (A0_1) edge [|->,gray]node [auto] {$\scriptstyle{}$} (A0_0);   \end{tikzpicture}   \] 
are well defined and quasi-inverses of each other. In particular the
functor $\shP_{\shC}\colon\stX\arr\text{Fib}_{\stX,\shC}$ is an equivalence
of stacks.
\end{thm}
An immediate corollary and generalization of Theorem \ref{thm: Tannaka reconstrion for Loc X}
is the following.
\begin{cor}
\label{cor:Lurie s point of view} Let $\stX$ be a quasi-compact
fpqc stack over $R$ with quasi-affine diagonal, $\shC\subseteq\Loc\stX$
be a monoidal subcategory with duals that generates $\QCoh\stX$ and
$\stY$ be a fibered category over $R$. Then the functor 
\[
\Hom(\stY,\stX)\arr\text{Fib}_{\stX,\shC}(\stY)\comma(\stY\arrdi f\stX)\longmapsto f_{|\shC}^{*}\colon\shC\arr\Loc(\stY)
\]
is an equivalence of categories.\end{cor}
\begin{proof}
The map in the statement is obtained applying $\Hom(\stY,-)$ to the
functor $\shP_{\shC}\colon\stX\arr\text{Fib}_{\stX,\shC}$, which
is an equivalence by \ref{thm: Tannaka reconstrion for Loc X}.
\end{proof}
One of the key points in the proof of statements above is a characterization
of the following stacks.
\begin{defn}
A \emph{pseudo-affine} stack is a quasi-compact fpqc stack with quasi-affine
diagonal such that all quasi-coherent sheaves on it are generated
by global sections. A map $f\colon\stY'\arr\stY$ of fibered categories
is called \emph{pseudo-affine }if for all maps $T\arr\stY$ from an
affine scheme the fiber product $T\times_{\stY}\stY'$ is pseudo-affine.\end{defn}
\begin{thm}
\label{prop:arbitrary intersection of quasi-compact open is pseudo-affine}
Intersection of quasi-compact open subschemes (thought of as sheaves)
of affine schemes are pseudo-affine. Conversely if $U$ is a pseudo-affine
stack then it is (equivalent to) a sheaf and it is the intersection
of the quasi-compact open subschemes of $\Spec\Hl^{0}(\odi U)$ containing
it.
\end{thm}
A quasi-affine scheme is pseudo-affine and, conversely, a pseudo-affine
stack which is algebraic is quasi-affine (see \cite[Proposition 3.1]{Gross2013}).
In general a pseudo-affine sheaf is not quasi-affine. An example is
the sheaf intersection of all the complement of closed points in $\Spec k[x,y]$,
where $k$ is a field.

The last statement of Theorem \ref{thm: Tannaka reconstrion for Loc X}
admits an almost converse. Let $\stX$ be a fibered category over
$R$ and $\shC\subseteq\Loc\stX$ be a full monoidal subcategory.
We have a functor $\shG_{*}\colon\shC\arr\Loc(\text{Fib}_{\stX,\shC})$
which maps a $\E\in\shC$ to the locally free sheaf
\[
\shG_{\E}\colon\text{Fib}_{\stX,\shC}^{\op}\arr\Ab\comma\shG_{\E}(\Gamma\in\text{Fib}_{\stX,\shC}(A))=\Gamma_{\E}\text{ for all }R\text{-algebras }A
\]
In particular $\shP_{\shC}^{*}\shG_{\E}\simeq\E$ for $\E\in\shC$.
Notice that, a priori, $\text{Fib}_{\stX,\shC}$ is not necessarily
fibered in groupoids and therefore the notion of a locally free sheaf
on it is not defined (although one can easily guess the definition).
Part of Theorems \ref{thm: Tannaka reconstrion for Loc X} and \ref{thm:when FibX,C has the resolution property}
(below) is to prove that, under suitable conditions on $\shC$, this
is indeed true.
\begin{thm}
\label{thm:when FibX,C has the resolution property} Let $\stX$ be
a quasi-compact fibered category over $R$ and $\shC\subseteq\Loc\stX$
be a full monoidal subcategory with duals and with $\Sym^{n}\E\in\shC$
for all $n\in\N$ if $\E$ has local rank not invertible in $R$.
If $f\colon\shC\arr\N$ is a function then $\text{Fib}_{\stX,\shC}^{f}$,
the sub-fibered category of $\text{Fib}_{\stX,\shC}$ of functors
$\Gamma$ such that $\rk\Gamma_{\E}=f(\E)$ for all $\E\in\shC$,
is a quotient stack of a pseudo-affine sheaf by the action of a (possibily
infinite) products of $\GL_{n}$; in particular it is a quasi-compact
fpqc stack in groupoids with affine diagonal. Moreover the subcategory
$\{\shG_{\E}\}_{\E\in\shC}\subseteq\Loc(\text{Fib}_{\stX,\shC}^{f})$
generates $\QCoh(\text{Fib}_{\stX,\shC}^{f})$ and, in particular,
$\text{Fib}_{\stX,\shC}^{f}$ has the resolution property.

Let $\shI$ be the set of functions $f\colon\shC\arr\N$ such that
there exists a geometric point $s\colon\Spec L\arr\stX$ with $f(\E)=\rk s^{*}\E$.
Given $f\colon\shC\arr\N$ then $\text{Fib}_{\stX,\shC}^{f}\neq\emptyset$
if and only if $f\in\shI$ and, if $\shI$ is finite, then 
\[
\text{Fib}_{\stX,\shC}=\bigsqcup_{f\in\shI}\text{Fib}_{\stX,\shC}^{f}
\]
In particular if $\shI$ is finite then $\text{Fib}_{\stX,\shC}$
is a quasi-compact fpqc stack in groupoids with affine diagonal, the
subcategory $\{\shG_{\E}\}_{\E\in\shC}\subseteq\Loc(\text{Fib}_{\stX,\shC})$
generates $\QCoh(\text{Fib}_{\stX,\shC})$ and $\text{Fib}_{\stX,\shC}$
has the resolution property.\end{thm}
\begin{rem}
The condition that $\shI$ is finite in \ref{thm:when FibX,C has the resolution property}
is not optimal, but at least it covers the case where $\stX$ admits
a surjective (on equivalence classes of geometric points) map $\stX'\arr\stX$
from an algebraic stack whose connected components are open (e.g.
$\stX$ is a connected or Noetherian algebraic stack). It is not clear
if $\text{Fib}_{\stX,\shC}$ always has the resolution property.
\end{rem}
We start with a first characterization of pseudo-affine stacks. Notice
that the condition that all quasi-coherent sheaves on $\stX$ are
generated by global section exactly means that the category $\{\odi{\stX}\}$
generates $\QCoh\stX$.
\begin{prop}
\label{lem:recognizing affine rings from global sections} Let $\stX\arrdi{\pi}\Spec R$
be a quasi-compact fpqc stack with quasi-affine diagonal. Then the
following conditions are equivalent:
\begin{enumerate}
\item the stack $\stX$ is pseudo-affine;
\item the map $\pi^{*}\pi_{*}\shF\arr\shF$ is surjective for all $\shF\in\QCoh\stX$;
\item the stack $\stX$ is equivalent to a sheaf and there exists a flat
monomorphism $\stX\arr\Spec B$, where $B$ is a ring.
\end{enumerate}
In this case the map $p\colon\stX\arr\Spec\Hl^{0}(\odi{\stX})$ is
a flat monomorphism, $p_{*}\colon\QCoh\stX\arr\Mod\Hl^{0}(\stX)$
is fully faithful and $p^{*}p_{*}\simeq\id$. Moreover if $\stX\times_{\Hl^{0}(\odi{\stX})}k\neq\emptyset$
for all geometric points $\Spec k\arr\Spec\Hl^{0}(\odi{\stX})$ then
$p$ is an isomorphism.\end{prop}
\begin{proof}
$2)\then1)$. Given $\shF\in\QCoh\stX$, take a surjective map $R^{(I)}\arr\pi_{*}\shF$.
In this case the composition
\[
\odi{\stX}^{(I)}\simeq\pi^{*}R^{(I)}\arr\pi^{*}\pi_{*}\shF\arr\shF
\]
 is surjective.

$1)\then2)$. A sheaf $\shF\in\QCoh\stX$ is generated by global sections
and the image of $\pi^{*}\pi_{*}\shF\arr\shF$ contains all of them.

$3)\then1)$. Denote by $p\colon\stX\arr\Spec B$ the flat monomorphism.
We are going to show that $\delta_{\shF}\colon p^{*}p_{*}\shF\arr\shF$
is an isomorphism for all $\shF\in\QCoh\stX$. Arguing as in $2)\then1)$
this will conclude the proof. By hypothesis there exists a representable
fpqc covering $h\colon\Spec C\arr\stX$ and we must prove that $h^{*}\delta_{\shF}$
is an isomorphism. Let $f=ph\colon\Spec C\arr\Spec B$ be the composition,
which is flat by hypothesis, and consider the commutative diagram
  \[   \begin{tikzpicture}[xscale=2.0,yscale=-1.0]     \node (A0_0) at (0, 0) {$\Spec C$};     \node (A0_1) at (1, 0) {$\stX\times_B C$};     \node (A0_2) at (2, 0) {$\Spec C$};     \node (A1_1) at (1, 1) {$\stX$};     \node (A1_2) at (2, 1) {$\Spec B$};     \path (A0_0) edge [->]node [auto] {$\scriptstyle{s}$} (A0_1);     \path (A0_1) edge [->]node [auto] {$\scriptstyle{\alpha}$} (A0_2);     \path (A1_1) edge [->]node [auto] {$\scriptstyle{p}$} (A1_2);     \path (A0_2) edge [->]node [auto] {$\scriptstyle{f}$} (A1_2);     \path (A0_0) edge [->]node [auto] {$\scriptstyle{h}$} (A1_1);     \path (A0_1) edge [->]node [auto] {$\scriptstyle{}$} (A1_1);   \end{tikzpicture}   \] 
Since $\alpha$ is a monomorphism with a section, $\alpha$ and $s$
are inverses of each other. Using the fact that $f$ is flat the map
$h^{*}\delta_{\shF}$ is given by
\[
p^{*}p_{*}\shF(\Spec C\arrdi h\stX)=p_{*}\shF(\Spec C\arrdi f\Spec B)=\shF(\stX\times_{B}C\arr\stX)\arrdi{s^{*}}\shF(\Spec C\arrdi h\stX)
\]
and therefore it is an isomorphism.

$1)\then3)$. Set $B=\Hl^{0}(\odi{\stX})$ and $p\colon\stX\arr\Spec B$
the induced map. Notice that $\L_{B}(\{\odi{\stX}\},B)\simeq\Mod B$
and under this isomorphism $\Omega^{*}\colon\QCoh\stX\arr\L_{B}(\{\odi{\stX}\},B)$
and $\shF_{*,\{\odi{\stX}\}}\colon\L_{B}(\{\odi{\stX}\},B)\arr\QCoh\stX$
correspond to $p_{*}\colon\QCoh\stX\arr\Mod B$ and $p^{*}\colon\Mod B\arr\QCoh\stX$
respectively. By hypothesis, \ref{prop:excatness properties of Omega* and shF*}
and \ref{thm:general theorem for equivalences} the map $p\colon\stX\arr\Spec B$
is flat, $p_{*}\colon\QCoh\stX\arr\Mod B$ is fully faithful and $p^{*}p_{*}\simeq\id$.

We want to show that $\stX\arr\Spec B$ is fully faithful or, equivalently,
that the diagonal $\stX\arr\stX\times_{B}\stX$ is an equivalence.
Let $V=\Spec C\arr\stX$ be a representable fpqc covering and denote
by $s\colon V\arr\stX\times_{B}V$ the graph of $h$. We have Cartesian
diagrams   \[   \begin{tikzpicture}[xscale=2.0,yscale=-1.2]     \node (A0_0) at (0, 0) {$V\times_\stX V$};     \node (A0_1) at (1, 0) {$V\times_B V$};     \node (A0_2) at (2, 0) {$V\times_\stX V$};     \node (A0_3) at (3, 0) {$V\times_B V$};     \node (A1_0) at (0, 1) {$\stX$};     \node (A1_1) at (1, 1) {$\stX\times_B \stX$};     \node (A1_2) at (2, 1) {$V$};     \node (A1_3) at (3, 1) {$\stX\times_B V$};     \path (A0_0) edge [->]node [auto] {$\scriptstyle{q}$} (A0_1);     \path (A0_1) edge [->]node [auto] {$\scriptstyle{}$} (A1_1);     \path (A1_0) edge [->]node [auto] {$\scriptstyle{\Delta}$} (A1_1);     \path (A0_3) edge [->]node [auto] {$\scriptstyle{}$} (A1_3);     \path (A0_2) edge [->]node [auto] {$\scriptstyle{}$} (A1_2);     \path (A1_2) edge [->]node [auto] {$\scriptstyle{s}$} (A1_3);     \path (A0_2) edge [->]node [auto] {$\scriptstyle{q}$} (A0_3);     \path (A0_0) edge [->]node [auto] {$\scriptstyle{}$} (A1_0);   \end{tikzpicture}   \] 
and the vertical arrows are representable fpqc coverings. By descent
it follows that $\Delta$ is an equivalence if we prove that $s$
is an equivalence.

Let $f\colon V\arr\stX\arrdi p\Spec B$ be the composition. Since
$p$ is flat also $f$ is flat and, in particular, $\Hl^{0}(\odi{\stX\times_{B}V})=C$.
Since $\stX\times_{B}V\arr\stX$ is affine it follows that $\{\odi{\stX\times_{B}V}\}$
generates $\QCoh(\stX\times_{B}V)$. Since $s$ is a section of $\stX\times_{B}V\arr V$,
we see that we can assume that $p\colon\stX\arr\Spec B$ has a section,
that we still denote by $s\colon\Spec B\arr\stX$. In this case we
have to prove that $p$ or $s$ is an equivalence. Notice that this
implies the last claim in the statement. Indeed if $\stX\times_{B}k\neq\emptyset$
for all geometric points then $f\colon V\arr\Spec B$ is an fpqc covering:
it follows that the map $\stX\arr\Spec B$ is an equivalence because
it has this property fpqc locally. 

We first prove that $p_{*}\colon\QCoh\stX\arr\Mod B$ is an equivalence.
It suffices to show that, if $M\in\Mod B$, then the map $\gamma_{M}\colon M\arr p_{*}p^{*}M$
is an isomorphism. Notice that $p^{*}\gamma_{M}$ is a section of
the map $(p*p_{*})p^{*}M\arr p^{*}M$ which is an isomorphism. So
$p^{*}\gamma_{M}$ and $\gamma_{M}=s^{*}p^{*}\gamma_{M}$ are isomorphisms.
Since $p_{*}s_{*}\colon\Mod B\arr\Mod B$ is the identity, we can
conclude that $s_{*}\simeq p^{*}$ and $s^{*}\simeq p_{*}$. Consider
the Cartesian diagram   \[   \begin{tikzpicture}[xscale=1.5,yscale=-1.2]     \node (A0_0) at (0, 0) {$U$};     \node (A0_1) at (1, 0) {$\Spec C$};     \node (A1_0) at (0, 1) {$\Spec B$};     \node (A1_1) at (1, 1) {$\stX$};     \path (A0_0) edge [->]node [auto] {$\scriptstyle{g}$} (A0_1);     \path (A1_0) edge [->]node [auto] {$\scriptstyle{s}$} (A1_1);     \path (A0_1) edge [->]node [auto] {$\scriptstyle{h}$} (A1_1);     \path (A0_0) edge [->]node [auto] {$\scriptstyle{t}$} (A1_0);   \end{tikzpicture}   \] 
where $h$ is a representable fpqc covering. Since $\stX$ has quasi-affine
diagonal it follows that $U$ is a quasi-affine scheme. Moreover
\[
\Hl^{0}(\odi U)\simeq g_{*}t^{*}\odi B\simeq h^{*}s_{*}\odi B\simeq h^{*}p^{*}\odi B\simeq C
\]
Thus $g\colon U\arr\Spec C$ and $s\colon\Spec B\arr\stX$ are open
immersions. We prove that $g$ is a surjective, which imply that $g$
and $s$ are isomorphisms. Let $Z$ be the complement of $U$ in $\Spec C$
with reduced structure. We have
\[
\Hl^{0}(\odi Z)\simeq p_{*}h_{*}\odi Z\simeq s^{*}h_{*}\odi Z=t_{*}g^{*}\odi Z=0
\]
Thus $Z$ is empty as required.\end{proof}
\begin{rem}
\label{rem: classifying stack of elliptic curve} The assumption on
the diagonal in \ref{lem:recognizing affine rings from global sections}
is necessary: the stack $\stX=\Bi_{k}E$, where $E$ is an elliptic
curve over a field $k$, is not a sheaf but $\QCoh\stX\simeq\QCoh k$.
\end{rem}

\begin{rem}
\label{rem: generating subcategory under pseudo-affine maps} If $f\colon\stY\arr\stX$
is a pseudo-affine map of pseudo-algebraic fiber categories then it
is quasi-compact with affine diagonal and the map $f^{*}f_{*}\shF\arr\shF$
is surjective for all $\shF\in\QCoh\stY$. In particular if $\shD\subseteq\QCoh\stX$
generates $\QCoh\stX$ then $f^{*}\shD$ generates $\QCoh\stY$. The
first claims follow by standard argument of descent while the second
by considering an atlas of $\stX$ and using the characterization
$3)$ of \ref{lem:recognizing affine rings from global sections}. 

If $\stY\arr\Spec A$ is a map of fibered categories then $\stY$
is pseudo-affine if and only if $\stY\arr\Spec A$ is pseudo-affine,
because if $\stY$ is pseudo-affine and $g\colon\stY'\arr\stY$ is
an affine map, then $\stY'$ is an fpqc sheaf with quasi-affine diagonal
and $\odi{\stY'}=g^{*}\odi{\stY}$ generates $\QCoh\stY'$. 

Finally, if $f\colon\stY'\arr\stY$ is a map of fpqc stacks and $Y\arr\stY$
is an fpqc atlas such that $Y'=Y\times_{\stY}\stY'\arr Y$ is pseudo-affine
then $f$ is pseudo-affine. Indeed by standard arguments of descent
we can assume $\stY=\Spec B$ and $Y=\Spec B'$ affine, which also
implies that $\stY'$ is quasi-compact with affine diagonal. Moreover
since $B\arr B'$ is flat we have $\Hl^{0}(\odi{Y'})\simeq\Hl^{0}(\odi{\stY'})\otimes_{B}B'$
and therefore we can assume that $\Hl^{0}(\odi{\stY'})=B$ and $\Hl^{0}(\odi{Y'})=B'$.
In this case $Y'\arr Y$ is flat and fully faithful and, since $\stY$
is an fpqc stack, it follows that also $\stY'\arr\stY$ is a flat
and fully faithful.
\end{rem}

\begin{rem}
\label{rem: cases for resolution property} Let $G$ be a flat and
affine group scheme over $R$ such that $\Bi_{R}G$ has the resolution
property. Taking into account \ref{rem: generating subcategory under pseudo-affine maps},
if $U$ is a pseudo-affine sheaf over $R$ with an action of $G$
then $[U/G]$ has the resolution property because the map $[U/G]\arr\Bi_{R}G$
is pseudo-affine. The same conclusion follows for a stack $\stX=[X/G]$,
where $X$ is a scheme, if there exists $\shL\in\Pic(\stX)$ whose
pullback $\shM$ to $X$ is very ample relatively to $R$. Indeed
$\stX$ can be written as $[U/G\times\Gm]$ where $U$ is the complement
of the zero section of $\shM\arr X$: the fact that $\shM$ is very
ample tells us that $U$ is quasi-affine. Moreover $\Bi(G\times\Gm)$
has the resolution property. Indeed let $N$ be the canonical invertible
sheaf on $\Bi_{R}\Gm$ and $\shF\in\QCoh(\Bi_{R}G)$. The action of
$\Gm$ yields a decomposition 
\[
\shF=\bigoplus_{n\in\Z}\shF_{n}
\]
and all $\shF_{n}=(\shF\otimes N^{\otimes-n})^{\Gm}$ are sub $G\times\Gm$-representations.
If $x\in\shF_{n}$ there exists $\E\in\Loc(\Bi_{R}G)$ and a $G$-equivariant
map $\E\arr F_{n}$ with $x$ in its image. Thus the map $\E\otimes N^{\otimes n}\arr F_{n}$,
where $\E$ has the trivial action of $\Gm$, is $G\times\Gm$-equivariant
and has $x$ in its image.
\end{rem}
The following property is known for algebraic stacks (see \cite[Corollary 5.11]{Gross2013}
and \cite[Proposition 1.3]{Totaro2002}).
\begin{cor}
\label{cor:quasi-affine and resolution implies affine} A quasi-compact
fpqc stack $\stX$ with quasi-affine diagonal and with the resolution
property has affine diagonal.\end{cor}
\begin{proof}
Since $\stX$ is pseudo-algebraic the category $\Loc\stX$ is essentially
small. Thus we can consider a set $\shR$ of representatives of locally
free sheaves over $\stX$. Given $\E\in\shR$ we define the sheaf
\[
\Fr(\E)\colon(\Sch/\stX)^{\text{op}}\arr\sets\comma\Fr(\E)(U\arrdi f\stX)=\bigsqcup_{n\in\N}\Iso_{U}(\odi U^{n},f^{*}\E)
\]
It is easy to see that $\Fr(\E)\arr\stX$ is an affine fpqc covering.
In particular $g\colon\Fr=\prod_{\E\in\shR}\Fr(\E)\arr\stX$ is also
an affine fpqc covering and $\Fr$ is quasi-compact with quasi-affine
diagonal. Thus it is enough to show that $\Fr$ has affine diagonal.
Since $g$ is affine by \ref{rem: generating subcategory under pseudo-affine maps}
$g^{*}\Loc\stX$ generates $\QCoh\Fr$. On the other hand if $\E\in\Loc\stX$
then by construction $\Fr$ is a (finite) disjoint union of open substacks
over which $g^{*}\E$ is free, which implies that $g^{*}\E$ is generated
by global sections. We can conclude that $\Fr$ is a pseudo-affine
sheaf and thus has affine diagonal.
\end{proof}

\begin{proof}
(of Theorem \ref{thm: Tannaka reconstrion for Loc X}). Since $\stX$
has affine diagonal by \ref{cor:quasi-affine and resolution implies affine},
all morphisms $\Spec A\arr\stX$ are affine. Therefore the functor
$\stX(A)\arr\QAlg_{A}\stX$ which maps $s\colon\Spec A\arr\stX$ to
$s_{*}\odi A$ is fully faithful. By \ref{thm:fundamental for sheaf of algebras and monoidal functors}
and the fact that 
\[
\Omega_{\E}^{s_{*}\odi A}=\Hom(\E,s_{*}\odi A)\simeq\duale{(s^{*}\E)}\text{ for }s\in\stX(A)\comma\E\in\shC
\]
we can conclude that the functor $\stX(A)\arr\SMex_{R}(\shC,A)$,
$s\longmapsto\duale{(s_{|\shC}^{*})}$ is well defined and fully faithful.
Thus everything follows if we prove that, given $\Gamma\in\SMex_{R}(\shC,A)$,
the composition $p\colon\Spec\alA_{\Gamma,\shC}\arr\stX_{A}\arr\Spec A$
is an isomorphism. Set $\alA=\alA_{\Gamma,\shC}$, $\stY=\Spec\alA$
and $f\colon\stY\arr\stX$ for the structure morphism. We want to
apply \ref{lem:recognizing affine rings from global sections} on
$p\colon\stY\arr\Spec A$. Notice that $\stY$ is quasi-compact and
has affine diagonal. Moreover by \ref{thm:fundamental for sheaf of algebras and monoidal functors}
\[
\Gamma_{\E}\simeq\Omega_{\E}^{\alA}=\Hom(\E,\alA)=\Hom(\E,f_{*}\odi{\stY})\simeq\Hom(f^{*}\E,\odi{\stY})=\Hl^{0}(\duale{(f^{*}\E)})
\]
In particular, since $\Gamma_{\odi{\stX}}\simeq A$, the map $A\arr\Hl^{0}(\alA)$
is an isomorphism. We show now that, if $\Spec k\arr\Spec A$ is a
geometric point, then $\stY\times_{A}k\neq\emptyset$. If $g\colon\stX_{k}\arr\stX_{A}$
is the projection then $\stY\times_{A}k\simeq\Spec(g^{*}\alA)$, while
by \ref{prop:base change for FstarC} we have $g^{*}\alA\simeq\alA_{\Gamma\otimes_{A}k,\shC}$.
By \ref{thm:fundamental for sheaf of algebras and monoidal functors}
we get $\Gamma\otimes_{A}k\simeq\Omega^{\alA_{\Gamma\otimes_{A}k,\shC}}$
and therefore $\alA_{\Gamma\otimes_{A}k,\shC}=0$, that is $\stY\times_{A}k=\emptyset$,
implies $\Gamma\otimes_{A}k=0$, while $\Gamma_{\odi{\stX}}\otimes_{A}k=A\otimes_{A}k\simeq k$.

It remains to show that $\{\odi{\stY}\}$ generates $\QCoh\stY$.
Since $\stY\arrdi f\stX$ is affine, $f^{*}\shC$ generates $\QCoh\stY$
by \ref{rem: generating subcategory under pseudo-affine maps}. Thus
we have to prove that all the sheaves $f^{*}\E$ are generated by
global sections. By hypothesis the map 
\[
\Hl^{0}(\duale{f^{*}\E'})\otimes\Hl^{0}(\duale{f^{*}\E})\arr\Hl^{0}(\duale{f^{*}(\E'\otimes\E)})
\]
is an isomorphism for all $\E,\E'\in\shC$. Since $\shC$ has duals,
choosing $\E'=\duale{\E}$ the above map became the evaluation
\[
\omega\colon\Hl^{0}(f^{*}\E)\otimes\Hom_{\stY}(f^{*}\E,\odi{\stY})\arr\End_{\stY}(f^{*}\E)
\]
Since $\omega$ is an isomorphism there exist $x_{1},\dots,x_{n}\in\Hl^{0}(f^{*}\E)$,
$\phi_{1},\dots,\phi_{n}\in\Hom(f^{*}\E,\odi{\stX})$ such that $\id_{f^{*}\E}=\omega(\sum_{i}x_{i}\otimes\phi_{i})$.
This implies that the map $\odi{\stY}^{n}\arr f^{*}\E$ given by the
global sections $x_{1},\dots,x_{n}$ is surjective, as required.

For the last statement we claim that $\duale -\colon\text{Fib}_{\stX,\shC}(A)\arr\SMex_{R}(\shC,A)$,
$\Gamma\longmapsto\duale{\Gamma}$ is an equivalence. Taking into
account \ref{ prop: Lex and left exact functors}, if $\Gamma\in\text{Fib}_{\stX,\shC}(A)$
then $\duale{\Gamma}\in\SMex_{R}(\shC,A)$ because $\Gamma$ is exact
on all finite test sequences and the dual of a right exact sequence
of locally free sheaves is again exact. Moreover the map $\stX(A)\arr\SMex_{R}(\shC,A)$
factors as 
\[
\stX(A)\arrdi{\shP_{\shC}}\text{Fib}_{\stX,\shC}(A)\arrdi{\duale -}\SMex_{R}(\shC,A)
\]
Thus $\duale -\colon\text{Fib}_{\stX,\shC}(A)\arr\SMex_{R}(\shC,A)$
is essentially surjective and $\duale -\colon\SMex_{R}(\shC,A)\arr\text{Fib}_{\stX,\shC}(A)$
is a quasi-inverse.
\end{proof}

\begin{proof}
(of Theorem \ref{prop:arbitrary intersection of quasi-compact open is pseudo-affine})
Let $X=\Spec B$ be an affine scheme, $\{U_{i}\}_{i\in I}$ be a set
of quasi-compact open subsets of $X$ and set $U=\cap_{i}U_{i}$.
If $i\in I$ the subscheme $U_{i}$ is the complement of the zero
locus of finitely many element of $B$ and thus there exists a free
$B$-module $E_{i}$ and a map $\phi_{i}\colon E_{i}\arr B$ such
that $U_{i}$ is the locus where $\phi_{i}$ is surjective. Let $V\colon\Aff/B\arr\sets$
be the functor 
\[
V(A)=\{(s_{i})_{i\in I}\st s_{i}\in E_{i}\otimes_{B}A\st\phi_{i}(s_{i})=1\}
\]
which is an affine scheme. The map $V\arr\Spec B$ factors through
$U$ and $V\arr U$ is surjective (as functors). Moreover if $\Spec A\arr U$
is a map then $V\times_{U}A=V\times_{X}A$ because $U\arr X$ is a
monomorphism. Since $V\times_{X}A$ is isomorphic to $(\prod_{i}\Ker\phi_{i})\times_{X}A$
we can conclude that $V\arr U$ is an affine fpqc epimorphism, so
that $U$ is quasi-compact, and that $U\arr X$ is flat. The result
then follows from \ref{lem:recognizing affine rings from global sections}.

For the converse, denote by $Z$ the intersection in the statement,
set $\shC=\{\odi U\}$, $B=\Hl^{0}(\odi U)$ and let $\alpha\colon\Spec A\arr\Spec B$
be a map which factors through $Z$ and $\shT_{*}\colon\odi U^{n}\arr\odi U^{m}\arr\odi U\arr0$
be an exact sequence on $U$. Since $C=\{\odi{\Spec B}\}$ as monoidal
categories, by \ref{thm: Tannaka reconstrion for Loc X} it is enough
to show that $\alpha_{|\shC}^{*}\colon\shC\arr\Loc A$ is exact on
$\shT_{*}$. The sequence $\shT_{*}$ defines a complex $\shW_{*}$
of free $A$-modules, namely $\shW_{*}=\Hl^{0}(\shT_{*})$, and the
locus $W$ in $\Spec A$ where $\shW_{*}$ is exact is quasi-compact,
open and contains $U$. Thus $Z\subseteq W$, the sequence $\shW_{*}$
become exact on $Z$ and therefore $\alpha^{*}$ maintains its exactness,
as required.\end{proof}
\begin{lem}
\label{lem:generators for GLn} Let $\shR$ be a set, $f\colon\shR\arr\N$
be a map and set 
\[
\GL_{f}=\prod_{i\in\shR}\GL_{f(i)}
\]
and $\shF_{i}$ for the locally free sheaf of rank $f(i)$ on $\Bi\GL_{f}$
coming from the universal one on $\Bi\GL_{f(i)}$. Then the subcategory
of $\Loc(\Bi\GL_{f})$ consisting of all tensor products of sheaves
\[
\{\shF_{i}\}_{i\in\shR}\comma\{\duale{(\Sym^{m}\shF_{i})}\}_{i\in\shR\text{ s.t. }f(i)\notin R^{*}\comma m\in\N}
\]
generates $\QCoh(\Bi\GL_{f})$.\end{lem}
\begin{proof}
We are going to apply \ref{thm:generating sheaves for BG}. Let $\shD$
be the subcategory of $\Loc\Bi\GL_{f}$ generated by direct sums and
tensor products of the sheaves $\Sym^{m}\shF_{i},(\det\shF_{i})^{-1}$
for $i\in\shR$, $m\in\N$. We claim that $R[\GL_{f}]$ is a direct
limit of representations in $\shD$. The representation $R[\GL_{f}]$
is a direct limit of tensor products of the regular representations
$R[\GL_{f(i)}]$ for $i\in\shR$. This allows us to reduce to the
case of $\GL_{n}$ (i.e. $\shR$ has one element). Call $\shF$ the
universal rank $n$ sheaf on $\Bi\GL_{n}$ and $\shG$ the free $R$-module
of rank $n$. As usual we can write $R[\GL_{n}]=R[X_{uv}]_{\det}$
for $1\leq u,v\leq n$, where $\det$ is the determinant polynomial.
The $R$-submodule generated by the $X_{u,v}$ is a $\GL_{n}$-subrepresentation
isomorphic to $\shF\otimes\shG$. In particular we obtain injective
maps
\[
\Sym^{nm}(\shF\otimes\shG)\otimes(\det\shF)^{-\otimes m}\arr R[\GL_{n}]
\]
whose image is the set of fractions $f/\det^{m}$ where $f$ is homogeneous
of degree $nm$. Those images form an incresing sequence of sub representations
saturating $R[\GL_{n}]$. Thus $R[\GL_{n}]$ is the direct limit of
the sheaves $\Sym^{nm}(\shF\otimes\shG)\otimes(\det\shF)^{-\otimes m}$
which belongs to $\shD$.

Coming back to the general framework, by \ref{thm:generating sheaves for BG}
and the existence of surjective maps $\shF_{i}^{\otimes f(i)}\arr\det\shF_{i}$
and, if $f(i)\in R^{*}$ so that $\GL_{f(i)}$ is linearly reductive,
$(\shF_{i}^{\vee})^{\otimes m}\arr(\Sym^{m}\shF_{i})^{\vee}$ we get
the desired result.
\end{proof}

\begin{proof}
(of Theorem \ref{thm:when FibX,C has the resolution property}). It
is easy to see that $\text{Fib}_{\stX,\shC}$ is a stack (not necessarily
in groupoids) for the fpqc topology on $\Aff/R$. To avoid problems
with disjoint unions we can assume that $\stX$ is a Zariski stack.
We will use notation from \ref{lem:generators for GLn}. Let $\shR$
be a set of representatives of $\shC/\simeq$ with $\odi{\stX}\in\shR$,
$J$ be a finite subset of $\shR$ and denote by $\shI_{J}$ the set
of $f\in\N^{J}$ extending to a function of $\shI$. Given $f\in\N^{J}$
we denote by $\stX_{f}$ the open locus of $\stX$ where $\rk\E=f(\E)$
for all $\E\in J$. Notice that $\stX=\bigsqcup_{f\in\shI_{J}}\stX_{f}$
and, since $\stX$ is quasi-compact, $\shI_{J}$ is finite. The sheaves
$(\E_{|\stX_{f}})_{\E\in J}$ induces a map $\stX_{f}\arr\Bi\GL_{f}=\shB_{f}$
and thus a map
\[
\omega_{J}\colon\stX\arr\bigsqcup_{f\in\shI_{J}}\shB_{f}=\shB_{J}
\]
For all $\E\in J$ there is a locally free sheaf $\shH_{\E,J}$ on
$\shB_{J}$ such that $(\shH_{\E,J})_{|\shB_{f}}$ is the canonical
locally free sheaf of rank $f(\E)$ (pullback from $\Bi\GL_{f(\E)}$),
so that $\omega_{J}^{*}\shH_{\E,J}\simeq\E$. Let $\shD_{J}$ be the
subcategory of $\Loc(\shB_{J})$ consisting of all tensor products
and duals of $\shH_{\E,J}$ and, if there exists $f\in\shI_{J}$ such
that $f(\E)\notin R^{*}$, also of $\Sym^{n}\shH_{\E,J}$ for $n\in\N$
and $\E\in J$. We have $\omega_{J}^{*}\shD_{J}\subseteq\shC$ by
construction, that $\shD_{J}$ generates $\QCoh(\shB_{J})$ by \ref{lem:generators for GLn}
and, in particular, that $\shP_{\shD_{J}}\colon\shB_{J}\arr\text{Fib}_{\shB_{J},\shD_{J}}$
is an equivalence of stacks by \ref{thm: Tannaka reconstrion for Loc X}.
When $J=\{\E\}$ we will replace $J$ by $\E$ in the subscripts.

We now show that $\text{Fib}_{\stX,\shC}$ is a stack in groupoids.
If $\Gamma,\Gamma'\in\text{Fib}_{\stX,\shC}(A)$, $\delta\colon\Gamma\arr\Gamma'$
is a morphism and $\E\in\shC$ then $\delta_{\E}\colon\Gamma_{\E}\arr\Gamma'_{\E}$
is an isomorphism because $\Gamma\circ\omega_{\E}^{*}\arrdi{\delta\circ\omega_{\E}^{*}}\Gamma'\circ\omega_{\E}^{*}$
is a morphism in the groupoid $\text{Fib}_{\stB_{\E},\shD_{\E}}(A)\simeq\shB_{\E}(A)$.

Given $f\colon\shC\arr\N$, we show that $\text{Fib}_{\stX,\shC}^{f}\neq\emptyset$
if and only if $f\in\shI$. For the if part, if $f\in\shI$ there
exists $s\colon\Spec L\arr\stX$ such that $\rk s^{*}\E=f(\E)$ for
all $\E\in\shC$. Thus $s_{|\shC}^{*}\in\text{Fib}_{\stX,\shC}^{f}$.
For the only if part we show that, if $L$ is an algebraically closed
field and $\Gamma\in\text{Fib}_{\stX,\shC}(L)$ then $ $$\rk\Gamma_{*}\colon\shC\arr\N$
belongs to $\shI$. This will also show that, if $\shI$ is finite,
then $\text{Fib}_{\stX,\shC}$ is the disjoint union of the $\text{Fib}_{\stX,\shC}^{f}$
for $f\in\shI$: indeed $\text{Fib}_{\stX,\shC}^{f}$ is the locus
where $\rk\shG_{\E}=f(\E)$ for all $\E\in\shC$.

So let $\Gamma\in\text{Fib}_{\stX,\shC}(L)$. Given $J\subseteq\shR$
finite consider $\Gamma\circ\omega_{J}^{*}\in\text{Fib}_{\shB_{J},\shD_{J}}(L)\simeq\shB_{J}(L)$,
so that there exists a map $\Spec L\arrdi s\shB_{J}$ such that $\Gamma\circ\omega_{J}^{*}\simeq s_{|\shD_{J}}^{*}$.
The map $s$ has image in some component $\shB_{f}$ with $f\in\shI_{J}$.
In particular if $\E\in J$ we have 
\[
\rk\Gamma_{\E}=\rk\Gamma_{\omega_{J}^{*}\shH_{\E,J}}=\rk(s^{*}(\shH_{\E,J})_{|\shB_{f}})=f(\E)
\]
This shows that for all finite subsets $J\subseteq\shR$ we have that
$f_{J}=(\rk\Gamma_{*})_{|J}$ belongs to $\shI_{J}$. By construction
$\stX_{f_{J}}$ is a non-empty (open and) closed substack of $\stX$.
Since $\stX_{f_{J}}\subseteq\stX_{f_{J'}}$ if $J'\subseteq J$ and
$\stX$ is quasi-compact it follows that $\bigcap_{J}\stX_{f_{J}}\neq\emptyset$
and thus that $\rk\Gamma_{*}\in\shI$.

It remains to show the claims about $\text{Fib}_{\stX,\shC}^{f}$
when $f\in\shI$. The sheaf $\shG_{\E}$ has rank $f(\E)$ on $\text{Fib}_{\stX,\shC}^{f}$
and thus we get a map $g\colon\text{\text{Fib}}_{\stX,\shC}^{f}\arr\Bi\GL_{f}$
such that $g^{*}\shF_{\E}\simeq\shG_{\E}$ for $\E\in\shR$ (here
we are using notation from \ref{lem:generators for GLn}). Let $\shC'$
be the subcategory of $\Loc\Bi\GL_{f}$ obtained by taking tensor
products and duals of the sheaves $\shF_{\E}$ and, if $f(\E)\notin R^{*}$,
also of $\Sym^{n}\shF_{\E}$ for $n\in\N$ and $\E\in\shR$. Moreover
set $\shD=\{\shG_{\E}\}_{\E\in\shC}\subseteq\Loc(\text{Fib}_{\stX,\shC})$.
We claim that $g^{*}\shC'\subseteq\shD_{|\text{Fib}_{\stX,\shC}^{f}}$.
It suffices to show that if $\E\in\shC$ then $\shG_{\duale{\E}}\simeq\duale{(\shG_{\E})}$
and, if $f(\E)\notin R^{*}$, $\shG_{\Sym^{n}\E}\simeq\Sym^{n}\shG_{\E}$
over the whole $\text{Fib}_{\stX,\shC}$. Consider $J=\{\E\}\subseteq\shR$.
The functor $\omega_{\E}\colon\stX\arr\shB_{\E}$ induces $\omega_{\E}^{*}\colon\shD_{\E}\arr\shC$
and a functor $\delta\colon\text{Fib}_{\stX,\shC}\arr\text{Fib}_{\shB_{\E},\shD_{\E}}\simeq\shB_{\E}$.
If $t$ is either $(-)^{\vee}$ or $\Sym^{n}$ and $\Gamma\in\text{Fib}_{\stX,\shC}$
we have isomorphisms
\[
\shG_{t(\E)}(\Gamma)=\Gamma_{t(\E)}\simeq(\Gamma\circ\omega_{\E}^{*})_{t(\shH_{\E,\{\E\}})}\simeq\delta(\Gamma)^{*}t(\shH_{\E,\{\E\}})\simeq t(\delta(\Gamma)^{*}\shH_{\E,\{\E\}})\simeq t(\Gamma_{\E})=t(\shG_{\E})(\Gamma)
\]
natural in $\Gamma$ and thus that $\shG_{t(\E)}\simeq t(\shG_{\E})$.

Using \ref{lem:generators for GLn} and \ref{rem: generating subcategory under pseudo-affine maps}
all the claims in the statement follow if we show that the fiber product
$U$ of $g\colon\text{\text{Fib}}_{\stX,\shC}^{f}\arr\Bi\GL_{f}$
along the canonical map $\Spec R\arr\Bi\GL_{f}$ is pseudo-affine.
Since $\Loc(\stX)$ is essentially small, it is easy to find a sub
monoidal subcategory $\tilde{\shC}$ of $\shC$ which is small and
such that $\widetilde{\shC}\arr\shC$ is an equivalence. Via restriction
we get an equivalence $\text{Fib}_{\stX,\shC}\arr\text{Fib}_{\stX,\tilde{\shC}}$
and we can assume that $\shC$ is small. If $A$ is an $R$-algebra
then $U(A)$ is the groupoid of $\Gamma\in\text{Fib}_{\stX,\shC}^{f}(A)$
together with basis of $\Gamma_{\E}$ for all $\E\in\shC$. In particular
it is easy to see that $U$ is (equivalent to) a sheaf. Let $V\colon\Aff/R\arr\sets$
be the functor which maps an $R$-algebra $A$ to the set of $R$-linear
and strong monoidal functors $\Gamma\colon\shC\arr\Loc A$ such that
$\rk\Gamma_{\E}=f(\E)$ together with a basis of $\Gamma_{\E}$ for
$\E\in\shC$. The sheaf $V$ is affine because it is a closed subscheme
of 
\[
\overline{V}=\prod_{\E\arr\E'\in\Arr(\shC)}\Homsh(R^{f(\E)},R^{f(\E')})\times\prod_{\E,\E'\in\shC}\Isosh(R^{f(\E)}\otimes R^{f(\E')},R^{f(\E\otimes\E')})\times\mathbb{G}_{m,R}\colon\Aff/R\arr\sets
\]
which is affine. Write $V=\Spec B$ and denote by $\Gamma\colon\shC\arr\Loc B$
the canonical $R$-linear and strong monoidal functor. Given a finite
test sequence $\shT$ in $\shC$ the sequence of maps $\Gamma_{\shT}$
is a complex of free $B$-modules and denote by $V_{\shT}$ the locus
in $V$ where this complex is exact. Clearly $U$ is the intersection
of the $V_{\shT}$ for all finite test sequences $\shT$. Since one
can easily check that $V_{\shT}$ is a quasi-compact open subscheme
of $V$ the result follows from \ref{prop:arbitrary intersection of quasi-compact open is pseudo-affine}.
\end{proof}
\begin{comment}
Denote by $\SMon_{R}^{G}$ the stack over $R$ whose fiber over an
$R$-algebra $A$ is the category of $R$-linear strong monoidal functors
$\Loc^{G}R\arr\Loc A$ which are exact on short exact sequences in
$\Loc^{G}R$. Taking into account \ref{rem: covariant to controvariant},
Theorem \ref{thm: Tannaka reconstrion for Loc X} with $\shC=\Loc^{G}R$
can be rephrased as follows. 
\begin{thm}
\label{thm:Tannaka reconstruction} If $\Bi_{R}G$ has the resolution
property the following functor is an equivalence of stacks   \[   \begin{tikzpicture}[xscale=3.2,yscale=-0.7]     \node (A0_0) at (0, 0) {$\Bi_R G $};     \node (A0_1) at (1, 0) {$\SMon^G_R$};     \node (A1_0) at (0, 1) {$( \Spec A\arrdi s \Bi_R G)$};     \node (A1_1) at (1, 1) {$s^*_{|\Loc^GR}$};     \path (A0_0) edge [->]node [auto] {$\scriptstyle{}$} (A0_1);     \path (A1_0) edge [|->,gray]node [auto] {$\scriptstyle{}$} (A1_1);   \end{tikzpicture}   \] 
\end{thm}
This result has already been proved in \cite[Theorem 1.3.2]{Schappi2011}
and it is a generalization of usual Tannaka's recontruction (\cite[Theorem 3.2]{Deligne1982}),
where $R$ is supposed to be a field.
\end{comment}
\begin{comment}
We now focus on the case in which $G$ is a finite, flat and finitely
presented group scheme over $R$. Since $R[G]\in\Loc^{G}R$, we see
that $\Bi_{R}G$ satisfies the resolution property.
\begin{defn}
\label{def: G cover} A $G$-\emph{cover }of an $R$-scheme $T$ is
a cover $f\colon X\arr T$ together with an action of $G$ on $X$
such that $f$ is invariant and $f_{*}\odi X$ and $R[G]\otimes\odi T$
are fppf locally isomorphic as $G$-comodules (not as rings). 

We denote by $\GCov$ the stack over $R$ of $G$-covers and by $\LMon_{R}^{G}$
the stack over $R$ whose fiber over an $R$-scheme $T$ is the groupoid
of $R$-linear monoidal functors $\Gamma\colon\Loc^{G}R\arr\Loc T$
such that $\rk\Gamma_{V}=\rk V$ for all $V\in\Loc^{G}R$, where $\rk$
is the rank function on the connected components, and $\Gamma$ is
exact on all short exact sequences in $\Loc^{G}R$.
\end{defn}
The stack $\GCov$ has been introduced in \cite{Tonini2011} and it
is algebraic and of finite type over $R$.
\begin{lem}
\label{lem: glrg} If $G$ is a linearly reductive and there exists
a set $I_{G}\subseteq\Loc^{G}R$ such that for all geometric points
(resp. for one geometric point if $\Spec R$ is connected) $\Spec k\arr\Spec R$
the map $(-\otimes_{R}k)\colon I_{G}\arr\Loc^{G}k$ is a bijection
onto (a set of representatives of) the irreducible representations
of $G\times_{R}k$ then 
\[
\eta_{M}\colon\bigoplus_{V\in I_{G}}V\otimes_{R}\Hom_{R}^{G}(V,M)\arr M\text{ for }M\in\Mod^{G}A
\]
is a $G$-equivariant and $A$-linear isomorphism.\end{lem}
\begin{proof}
If $V,W\in\Loc^{G}R$ then $\Hom^{G}(V,W)$ is locally free because
$G$ is linearly reductive (the map $\Hom^{G}(V,W)\arr\Hom(V,W)$
splits). Thus if $V,W\in I_{G}$ then $\Hom^{G}(V,W)=0$ if $V\neq W$
and $\Hom^{G}(V,V)=R\id_{V}$: since $G$ is linearly reductive taking
invariants commutes with arbitrary base change and it is enough to
check the rank on the geometric points (on the given geometric point
if $\Spec R$ is connected). In particular if $\Spec k\arrdi{\xi}\Spec R$
is any geometric point then $\xi^{*}\colon I_{G}\arr\Loc^{G}k$ is
injective with image a set of irreducible representations of $G\times k$.
This tells us that, if $\xi$ factors through $\Spec A$, then $\eta_{\xi^{*}M}$
is injective and, if $\xi^{*}(I_{G})$ is a set of representatives
of irreducible representations of $G\times k$, an isomorphism. If
$\Spec R$ is connected, so that $R[G]$ has constant rank, applying
this consideration to $M=R[G]$ we see that $\xi^{*}(I_{G})$ is a
set of representatives of irreducible representations of $G\times k$
for all geometric points $\xi$ of $\Spec R$. Thus $\eta_{M}$ is
an isomorphism on all geometric points of $\Spec A$. If $M$ is an
arbitrary direct sum of locally free $G$-comodules of finite rank
it follows that $\eta_{M}$ is an isomorphism. In general, using \ref{rem: crucial isomorphism invariants RG and modules},
we can find an exact sequence of $G$-comodules $V_{1}\arr V_{0}\arr M\arr0$
where the $V_{i}$ are sum of copies of $\duale{R[G]}\otimes A$.
Since $\eta_{V_{0}},\eta_{V_{1}}$ are isomorphisms, by functoriality
it follows that $\eta_{M}$ is an isomorphism.\end{proof}
\begin{thm}
\label{thm:gcovers and functors} The map of stacks   \[   \begin{tikzpicture}[xscale=3.2,yscale=-0.7]     \node (A0_0) at (0, 0) {$\GCov$};     \node (A0_1) at (1, 0) {$\LMon^G_R$};     \node (A1_0) at (0, 1) {$f\colon X\arr T$};     \node (A1_1) at (1, 1) {$(f_*\odi X\otimes -)^G$};     \path (A0_0) edge [->]node [auto] {$\scriptstyle{}$} (A0_1);     \path (A1_0) edge [|->,gray]node [auto] {$\scriptstyle{}$} (A1_1);   \end{tikzpicture}   \] 
is well defined, an open immersion and it extends the equivalence
in \ref{thm:Tannaka reconstruction}. It is an equivalence if $G$
satisfies the hypothesis of \ref{lem: glrg}.\end{thm}
\begin{proof}
Let $A$ be an $R$-algebra. If $\alA\in\Loc^{G}A$ is fppf locally
isomorphic to $R[G]\otimes A$ then the functor $(\alA\otimes-)^{G}\colon\Loc^{G}R\arr\Mod A$
is fppf locally $R$-linearly isomorphic to the forgetful functor
$(-\otimes_{R}A)\colon\Loc^{G}R\arr\Loc A$ (without monoidal structure)
thanks to \ref{rem: crucial isomorphism invariants RG and modules}.
This shows that the functor in the statement is well defined. It is
fully faithful thanks to \ref{rem: resolution property for BG cases}
and \ref{thm:general theorem for BG}. It extends the functor in \ref{thm:Tannaka reconstruction}
because if $f\colon X\arr\Spec A$ is a $G$-torsor corresponding
to $s\colon\Spec A\arr\Bi_{R}G$ then $s_{*}\odi A\simeq f_{*}\odi X$
as sheaves of algebras on $\Bi_{R}G$ and
\[
(s_{*}\odi A\otimes_{R}\shF)^{G}\simeq\Hom_{\Bi_{R}G}(\duale{\shF},s_{*}\odi A)\simeq\Hom_{A}(s^{*}\duale{\shF},A)\simeq s^{*}\shF\text{ for }\shF\in\Loc(\Bi_{R}G)=\Loc^{G}R
\]
We now prove that it is an open immersion. Let $\Gamma\in\LMon_{R}^{G}(A)$
and set $\alA=\alA_{\Gamma,\Loc(\Bi G)}\in\QAlg^{G}(A)$, which, by
definition of $\LMon_{R}^{G}$, is a locally free sheaf of algebras.
By \ref{thm:general theorem for BG} we have $\Gamma\simeq\Omega^{\alA}$
and we have to prove that the locus in $\Spec A$ where $\alA$ is
fppf locally the regular representation is open. Concretely, if $\xi\colon\Spec k\arr\Spec A$
is a geometric point and $\alA\otimes k\in\GCov(k)$ we will prove
that there exists a flat and finitely presented map $\Spec A'\arr\Spec A$
through which $\xi$ factors and such that $\alA\otimes A'\simeq A'[G]$.
Denote by $p\in\Spec A$ the image of $\xi$. The stack $\GCov$ is
algebraic and it is easy to see that it has affine diagonal. In particular
it is an fpqc stack. Moreover $\LMon_{R}^{G}$ is another fpqc stack
since the stack of locally free sheaves over $R$ is so. Since $\GCov$
is a full subcategory of $\LMon_{R}^{G}$, by descent along the fpqc
covering $\Spec k\arr\Spec\overline{k(p)}$ we obtain that $\alA\otimes\overline{k(p)}\in\GCov(\overline{k(p)})$
and, in particular, the existence of a $G$-equivariant isomorphism
$\overline{\omega}\colon\duale{(\overline{k(p)}[G])}\arr\duale{(\alA\otimes\overline{k(p)})}$.
Thus we can assume $k=\overline{k(p)}$. Since 
\[
\Hom^{G}(\duale{A[G]},M)\simeq\Hom^{G}(A,A[G]\otimes M)\simeq(A[G]\otimes M)^{G}\simeq M
\]
for all $G$-comodule $M$ over $A$ by \ref{rem: crucial isomorphism invariants RG and modules},
the map $\overline{\omega}$ is completely determined by a $\overline{\phi}\in\duale{\alA}\otimes\overline{k(p)}$.
There exists a finite field extension $L/k(p)$ such that $\overline{\phi}$
comes from some element in $\duale{\alA}\otimes L$ and it is a general
fact that we can find an fppf neighborhood $\Spec A'$ of $p$ in
$\Spec A$ with a point $p'\in\Spec A'$ over $p$ such that $k(p')=L$.
Up to shrink $\Spec A'$ around $p'$ we can assume to have $\phi\in\duale{\alA}$
inducing $\overline{\phi}$. The element $\phi$ defines a $G$-equivariant
map $\omega\colon\duale{A[G]}\arr\duale{\alA}$ of locally free sheaves
on $A$ inducing $\overline{\omega}$. Since $\overline{\omega}$
is an isomorphism it follows that $\omega$ is an isomorphism in a
Zariski open neighborhood of $p$ as required.

Assume now that $G$ satisfies the hypothesis of \ref{lem: glrg}
and let $\Gamma\in\LMon_{R}^{G}(A)$. We have that $\Omega=\Gamma_{\duale{(-)}}\in\MLex_{R}(\Loc^{G}R,A)$.
In particular, if $\alA=\alA_{\Omega,\Loc^{G}R}$ then $\Gamma_{V}\simeq(\alA\otimes_{R}V)^{G}$.
We have to prove that $\alA$ is fppf locally isomorphic to $R[G]\otimes A$.
This follows by the rank condition and \ref{lem: glrg}.\end{proof}
\begin{rem}
\label{rem: counterexample for gcovers and monoidal} In general the
functor in \ref{thm:gcovers and functors} is not an equivalence even
if $G$ is linearly reductive. Consider $R=\Q$, $G=\Z/3\Z$, $\alA=\overline{\Q}[x,y]/(x,y)^{2}$
with action of $G\times\overline{\Q}\simeq\mu_{3}$ given by $\deg x=\deg y=1$
and $\Gamma=(\alA\otimes_{\Q}-)^{G}\colon\Loc^{G}\Q\arr\Loc\overline{\Q}$.
It is easy to see that $\Gamma\in\LMon_{\Q}^{G}(\overline{\Q})$ and
it does not correspond to a $G$-cover, that is $\alA\notin\GCov(\overline{\Q})$.
\end{rem}

\begin{rem}
If $G$ satisfies the hypothesis of \ref{lem: glrg} there is a more
explicit way to pass from functors to sheaves, which may be useful
in concrete examples. Define
\[
\shF_{*}\colon\QAdd^{G}A\arr\Mod^{G}A\comma\shF_{\Gamma}=\bigoplus_{V\in I_{G}}\duale V\otimes\Gamma_{V}
\]
Using \ref{lem: glrg} it is easy to see that $\shF_{*}$ is a quasi-inverse
of $\Mod^{G}A\arr\QAdd^{G}A$, $M\longmapsto(-\otimes M)^{G}=\Hom^{G}(\duale{(-)},M)$,
the other natural isomorphism being
\[
\beta_{U}\colon\Hom^{G}(\duale U,\shF_{\Gamma})\simeq\bigoplus_{V\in I_{G}}\Hom^{G}(V,U)\otimes\Gamma_{V}\arr\Gamma_{U}\text{ for }\Gamma\in\QAdd^{G}A\comma U\in\Loc^{G}R
\]
The map $\beta_{U}$ is determined by $\alpha_{U}\colon\duale U\otimes\Gamma_{U}\arr\shF_{\Gamma}$.
It is easy to see that: $1)$ if $U\in I_{G}$ then $\alpha_{U}$
is the inclusion; $2)$ if $U=U_{1}\oplus U_{2}$ then $\alpha_{U}$
is zero on $\duale{U_{i}}\otimes\Gamma_{U_{j}}$for $i\neq j\in\{1,2\}$
and coincides with $\alpha_{U_{i}}$ on $\duale{U_{i}}\otimes\Gamma_{U_{i}}$
for all $i=1,2$; $3)$ if $U=\shH\otimes U'$ for $\shH\in\Loc R$
and $U'\in\Loc^{G}R$ then $\alpha_{U}$ is
\[
\duale U\otimes\Gamma_{U}\simeq\duale{\shH}\otimes\shH\otimes U'\otimes\Gamma_{U'}\arrdi{\text{ev}_{\shH}\otimes\alpha_{U'}}\arr\shF_{\Gamma}
\]
where $\text{ev}_{\shH}\colon\duale{\shH}\otimes\shH\arr R$ is the
evaluation; $4)$ if $\gamma\colon U\arr V$ is a $G$-equivariant
isomorphism then $\alpha_{U}=[(\duale{\gamma})^{-1}\otimes\Gamma_{\gamma}]\circ\alpha_{V}$.
By going through the definitions it is easy to show that if $\Gamma$
has a pseudo-monoidal structure the induced ring structure on $\shF_{\Gamma}$
is given by
\[
\duale V\otimes\Gamma_{V}\otimes\duale W\otimes\Gamma_{W}\arr\duale{(V\otimes W)}\otimes\Gamma_{V\otimes W}\arrdi{\alpha_{V\otimes W}}\shF_{\Gamma}\text{ for }V,W\in I_{G}
\]

\end{rem}
Let $R$ be a base commutative ring and $G$ be a flat, finite and
finitely presented group scheme over $R$. In \cite{Tonini2011} I
have introduced the notion of ramified Galois cover with group $G$,
briefly a $G$-cover, and the stack $\GCov$ of such objects (see
\ref{def: G cover} for details), which is algebraic and of finite
type over $R$ and contains $\Bi G$, the stack of $G$-torsors, as
an open substack. If $G$ is diagonalizable, its nice representation
theory makes possible to study $G$-covers in terms of simplified
data (collection of invertible sheaves and morphisms between them)
and to investigate the geometry of the moduli $\GCov$. 

The general case is much harder, even when $G$ is a constant group
over an algebraically closed field of characteristic zero: a direct
approach as in the diagonalizable case fails because of the complexity
of the representation theory of $G$. Thus in order to handle general
$G$-covers one need a different perspective and Tannaka's duality
comes into play. The $G$-torsors are very special $G$-covers and
the solution of Tannaka's reconstruction problem asserts that they
can be described in terms of particular strong monoidal functors with
domain $\Loc^{G}R$, the category of $G$-comodules over $R$ which
are projective and finitely generated as $R$-modules. Since a $G$-cover
is a ``weak'' version of a $G$-torsor it is natural to look at
a ``weak'' version of a strong monoidal functor, that is, as the
words suggest, a monoidal functor (see \ref{def: pseudo monoidal commutativi associative}
for a precise definition). If we denote by $\LMon_{R}^{G}$ the stack
over $R$ whose fiber over an $R$-scheme $S$ is the groupoid of
$R$-linear, exact (on short exact sequences) monoidal functors $\Gamma\colon\Loc^{G}R\arr\Loc S$
such that $\rk\Gamma_{V}=\rk V$ (pointwise) for all $V\in\Loc^{G}R$
we have:
\begin{thm*}
[\ref{thm:gcovers and functors}] The map of stacks 
\[
\GCov\arr\LMon_{R}^{G}\comma(X\arrdi fS)\longmapsto(f_{*}\odi X\otimes-)^{G}
\]
is an open immersion and it restricts to an equivalence from $\Bi G$
to the substack of $\LMon_{R}^{G}$ of strong monoidal functors. It
is an equivalence if $G$ is linearly reductive and there exists a
set $I_{G}\subseteq\Loc^{G}R$ such that for all geometric points
(resp. for one geometric point if $\Spec R$ is connected) $\Spec k\arr\Spec R$
the map $(-\otimes_{R}k)\colon I_{G}\arr\Loc^{G}k$ is a bijection
onto a set of representatives of the irreducible representations of
$G\times_{R}k$.
\end{thm*}
The second part of the statement is not true in general, even when
$G$ is linearly reductive (see \ref{rem: counterexample for gcovers and monoidal}). 

This paper in not about applications of the above theorem to $G$-covers,
although those exist and they are described in my Ph.D. thesis \cite[Section 1.3]{Tonini2013},
quite the contrary: we will deduce the above theorem as an application
of a much more general theory which is the heart of this paper.
\end{comment}


\bibliographystyle{amsalpha}
\addcontentsline{toc}{section}{\refname}\bibliography{biblio}

\end{document}
